\chapter{Circuits: KVL, KCL, impedance}
\label{vol04:ch09}

\epigraph{The current in any conductor is the algebraic sum of the
currents which the several electromotive forces would produce, each
acting alone in the otherwise unaltered circuit.}{Gustav Kirchhoff,
1845, paraphrased in the form in which his two laws have entered
engineering practice; see Alexander and Sadiku, \emph{Fundamentals of
Electric Circuits}, 6th ed.\ (McGraw-Hill, 2017), Section 2.4
\cite[Section~2.4]{text:alexander-sadiku-circuits}.}

\chapmeta{Half-life: durable. Archetypes: balance, transport, control. Exercise target: 42. Project track: physical build plus measurement (hybrid). Hazard class: low (bench-scale signals, mains avoided). Reader path default: standard, with sections explicitly tagged. Named cases: none directly; grounding and EMI failure modes cross-referenced from Chapter 8.}

\noindent
The lumped-parameter circuit is an engineering approximation to
Maxwell's equations valid when the device is small compared to a
wavelength of any signal it carries. Within that approximation
Kirchhoff's voltage law and Kirchhoff's current law become the
operational identities of the discipline: every circuit obeys an
energy balance around every loop and a charge balance at every node.
This chapter develops the working machinery (KVL, KCL, Th\'evenin and
Norton equivalents, superposition, transient response of first- and
second-order networks, AC phasor analysis, power, and three-phase
delivery) at the depth the practising engineer carries to a bench or
a SPICE deck.

\section{Lumped-parameter circuits}\pathtag{core}

A circuit is a network of two-terminal devices connected by ideal
conductors. The devices store, release, or dissipate energy; the
conductors carry current without resistance. The terminals of each
device sit at \engterm{nodes}; the conductors connecting them are
\engterm{branches}. The number of nodes $n$ and branches $b$ obeys
the topological identity $b - n + 1 = \ell$, where $\ell$ is the
number of independent loops. In a circuit with $b$ branches we have
$b$ unknown branch currents and $b$ unknown branch voltages; the
constitutive equations of the $b$ devices provide $b$ relations, and
the $n-1$ KCL equations plus $\ell$ KVL equations provide the
remaining $b$ topological relations. The system closes.

The lumped-parameter approximation requires that the physical extent
$L$ of every device and conductor is much smaller than the wavelength
$\lambda = c/f$ at the highest frequency $f$ of interest. At $1\,\si{
\mega\hertz}$, $\lambda = 300\,\si{\meter}$, and a $30\,\si{\centi
\meter}$ printed circuit board is lumped to within $0.1\,\%$. At
$1\,\si{\giga\hertz}$, $\lambda = 30\,\si{\centi\meter}$, and the same
board is no longer lumped; it requires transmission-line analysis
(Chapter 10). The lumped-circuit analysis of this chapter is the
working tool from DC to the low VHF band; above that, transmission-
line and distributed analysis takes over.

The five canonical two-terminal devices are:
\begin{itemize}[leftmargin=*]
  \item \engterm{Resistor}: $v = iR$, with $R$ in $\si{\ohm}$.
  \item \engterm{Capacitor}: $i = C\,\dd v/\dd t$, with $C$ in $\si{
    \farad}$.
  \item \engterm{Inductor}: $v = L\,\dd i/\dd t$, with $L$ in $\si{
    \henry}$.
  \item \engterm{Independent voltage source}: $v = V_{s}$ (specified),
    $i$ unconstrained.
  \item \engterm{Independent current source}: $i = I_{s}$ (specified),
    $v$ unconstrained.
\end{itemize}
Dependent (controlled) sources, in which the source value depends on
a current or voltage elsewhere in the circuit, are the building
blocks of small-signal models for transistors and operational
amplifiers; they appear in Section 9.4.

The sign convention for two-terminal devices is the
\engterm{passive sign convention}: the current $i$ enters the
positive terminal of $v$. Under this convention, $P = vi$ is the
instantaneous power absorbed by the device; positive $P$ means the
device dissipates or stores energy, negative $P$ means it delivers
energy back to the rest of the circuit. A source operating in its
forward mode has negative absorbed power; a resistor always has
non-negative absorbed power.

\section{Kirchhoff's voltage and current laws}\pathtag{core}

\begin{principle}[Kirchhoff's current law (KCL)]
The algebraic sum of currents leaving any node is zero,
\[
\sum_{k} i_{k} = 0.
\]
KCL is the lumped-circuit statement of the continuity equation
$\partial \rho/\partial t + \divg \vect{J} = 0$ in the limit that no
charge accumulates at any node. The lumped approximation embeds the
no-accumulation assumption.
\end{principle}

\begin{principle}[Kirchhoff's voltage law (KVL)]
The algebraic sum of voltages around any closed loop is zero,
\[
\sum_{k} v_{k} = 0.
\]
KVL is the lumped-circuit statement of Faraday's law $\oint \vect{E}
\cdot \dd\vect{\ell} = -\dd \Phi_{B}/\dd t$ in the limit that no
time-varying magnetic flux links the loop. The lumped approximation
embeds the no-stray-flux assumption.
\end{principle}

KCL and KVL together provide enough equations to solve any
lumped-parameter circuit once the device constitutive equations are
written. Two systematic solution methods follow.

The \engterm{nodal analysis} method uses KCL. Choose a reference node
(usually called ground), label the other $n-1$ nodes with their
unknown voltages $v_{1}, v_{2}, \ldots, v_{n-1}$, and write a KCL
equation at each non-reference node expressing every branch current
in terms of the node voltages and the device equations. The result is
$n-1$ equations in $n-1$ unknowns. For a circuit of only resistors
and current sources, the equations are linear and the solution is the
inverse of an $(n-1)\times(n-1)$ matrix; SPICE and every other circuit
simulator does this internally.

The \engterm{mesh analysis} method uses KVL. Identify the $\ell$
independent loops; assign a clockwise mesh current to each; write a
KVL equation around each loop expressing every branch voltage in
terms of the mesh currents and the device equations. The result is
$\ell$ equations in $\ell$ unknowns.

The two methods are duals. For a circuit dominated by current sources
and parallel connections, nodal is faster; for a circuit dominated by
voltage sources and series connections, mesh is faster. The working
engineer learns both and chooses by inspection.

% Figure: a two-loop circuit annotated for both nodal analysis (one unknown
% node voltage) and mesh analysis (two mesh currents). Plain TikZ.
\begin{figure}[htbp]
\centering
\begin{tikzpicture}[
  line width=0.8pt,
  res/.style={draw, fill=white, minimum width=12mm, minimum height=4.5mm, font=\scriptsize},
  font=\small,
]
% outer rectangle corners
\coordinate (tl) at (0,3);
\coordinate (tm) at (3,3);
\coordinate (tr) at (6,3);
\coordinate (bl) at (0,0);
\coordinate (bm) at (3,0);
\coordinate (br) at (6,0);
% left source (drawn as a circle with +/-)
\draw[engblue] (bl) -- (tl);
\node[draw, circle, minimum size=8mm, fill=white] (Vs) at (0,1.5) {};
\node[font=\scriptsize] at (0,1.5) {$V_s$};
\node[font=\scriptsize, anchor=east] at (-0.45,2.0) {$+$};
\node[font=\scriptsize, anchor=east] at (-0.45,1.0) {$-$};
% top resistor R1 (left to middle)
\node[res] (R1) at (1.5,3) {$R_1$};
\draw[engblue] (tl) -- (R1) -- (tm);
% top resistor R3 (middle to right)
\node[res] (R3) at (4.5,3) {$R_3$};
\draw[engblue] (tm) -- (R3) -- (tr);
% middle vertical resistor R2 (the node-voltage branch)
\node[res, rotate=90] (R2) at (3,1.5) {};
\node[font=\scriptsize, anchor=west] at (3.25,1.5) {$R_2$};
\draw[engblue] (tm) -- (R2) -- (bm);
% right vertical resistor R4
\node[res, rotate=90] (R4) at (6,1.5) {};
\node[font=\scriptsize, anchor=west] at (6.25,1.5) {$R_4$};
\draw[engblue] (tr) -- (R4) -- (br);
% bottom rail (ground)
\draw[engblue] (bl) -- (bm) -- (br);
\node[font=\scriptsize] at (3,-0.45) {reference (ground)};
% node-voltage label
\node[circle, fill=engblack, inner sep=1pt] at (tm) {};
\node[engred, font=\scriptsize, anchor=south] at (tm) {$v_1$};
% mesh-current arrows (two clockwise loops)
\draw[engorange, -{Stealth}] (1.5,1.5) ++(0:0.5) arc (0:300:0.5);
\node[engorange, font=\scriptsize] at (1.5,1.5) {$i_1$};
\draw[engorange, -{Stealth}] (4.5,1.5) ++(0:0.5) arc (0:300:0.5);
\node[engorange, font=\scriptsize] at (4.5,1.5) {$i_2$};
\end{tikzpicture}
\caption[A two-loop circuit for nodal and mesh analysis]{A two-loop ladder
network. Nodal analysis treats the bottom rail as reference and writes one
KCL equation at node $v_1$, giving a single unknown. Mesh analysis assigns
clockwise currents $i_1$ and $i_2$ to the two windows and writes one KVL
equation per loop, giving two unknowns. Both methods return the same branch
currents; nodal is shorter here because the circuit has a single
non-reference node, while mesh is shorter when voltage sources dominate.}
\label{fig:vol04ch09:nodal-mesh}
\end{figure}


Figure~\ref{fig:vol04ch09:nodal-mesh} shows a two-loop ladder annotated
for both methods. Series-parallel reduction handles the simpler case
when the network has no sources buried inside it:
Figure~\ref{fig:vol04ch09:series-parallel} reduces to a single
equivalent resistance by collapsing the parallel pair, then summing the
series chain.

% Figure: a resistor network combining series and parallel sections, drawn
% with plain TikZ paths and rectangular resistor symbols (no circuitikz).
\begin{figure}[htbp]
\centering
\begin{tikzpicture}[
  line width=0.8pt,
  res/.style={draw, fill=white, minimum width=14mm, minimum height=5mm},
  font=\small,
]
% source terminals on the left
\node[circle, fill=engblack, inner sep=1pt] (a) at (0,2) {};
\node[circle, fill=engblack, inner sep=1pt] (b) at (0,0) {};
\node[anchor=east, font=\scriptsize] at (a.west) {$+$};
\node[anchor=east, font=\scriptsize] at (b.west) {$-$};
\node[anchor=east] at (-0.4,1) {$V_s$};
% top wire to node n1
\node[circle, fill=engblack, inner sep=1pt] (n1) at (3,2) {};
\node[res] (R1) at (1.5,2) {$R_1$};
\draw[engblue] (a) -- (R1) -- (n1);
\node[font=\scriptsize, anchor=south] at (n1.north) {node 1};
% parallel branch R2 and R3 from n1 down to bottom rail
\node[res, rotate=90] (R2) at (3,1) {};
\node[font=\scriptsize, anchor=west] at (3.25,1) {$R_2$};
\node[res, rotate=90] (R3) at (4.5,1) {};
\node[font=\scriptsize, anchor=west] at (4.75,1) {$R_3$};
\node[circle, fill=engblack, inner sep=1pt] (n2) at (4.5,2) {};
\draw[engblue] (n1) -- (R2);
\draw[engblue] (R2) -- (3,0);
\draw[engblue] (n1) -- (n2) -- (R3);
\draw[engblue] (R3) -- (4.5,0);
% bottom return rail
\draw[engblue] (b) -- (6,0);
% series R4 on the way back along top
\node[res] (R4) at (5.5,2) {$R_4$};
\draw[engblue] (n2) -- (R4) -- (6,2) -- (6,0);
% labels for currents
\node[engred, font=\scriptsize, anchor=south] at (0.9,2.05) {$I$};
\draw[engred, -{Stealth}] (0.6,2.25) -- (1.0,2.25);
\end{tikzpicture}
\caption[Series-parallel resistor network]{A resistor network mixing series
and parallel connections. The source current $I$ flows through $R_1$ in
series, splits between $R_2$ and $R_3$ in parallel at node 1, then recombines
and passes through $R_4$ in series on the return. The equivalent resistance
seen by the source is $R_1 + (R_2 \parallel R_3) + R_4$, where
$R_2 \parallel R_3 = R_2 R_3/(R_2+R_3)$. Reduction proceeds by collapsing the
parallel pair to one resistor, then summing the series chain.}
\label{fig:vol04ch09:series-parallel}
\end{figure}


\section{Resistive networks: Th\'evenin, Norton, superposition}\pathtag{core}

For purely resistive linear circuits, three theorems compress the
analysis. Their proofs follow directly from linearity and from the
two Kirchhoff laws; the working consequences are powerful.

\begin{theorem}[Superposition]
In a linear circuit driven by multiple independent sources, the
response (any branch voltage or current) is the algebraic sum of the
responses to each source acting alone, with the other sources
deactivated. Deactivating a voltage source means short-circuiting it;
deactivating a current source means open-circuiting it.
\end{theorem}

Superposition reduces a multi-source circuit to a sum of single-source
sub-circuits, each easier to analyse. The price is that the method
fails on nonlinear elements (diodes, transistors above their small-
signal regime); for nonlinear elements, either linearise about an
operating point first (small-signal model) or use a numerical solver.

\begin{theorem}[Th\'evenin]
Any linear two-terminal network containing independent sources,
dependent sources, and resistors can be replaced, at its accessible
terminals, by an equivalent circuit consisting of a single voltage
source $V_{\text{Th}}$ in series with a single resistance $R_{\text{Th}}$.
\begin{itemize}[leftmargin=*]
  \item $V_{\text{Th}}$ is the open-circuit voltage at the terminals.
  \item $R_{\text{Th}}$ is the resistance seen looking back into the
    terminals with all independent sources deactivated. For circuits
    with dependent sources, $R_{\text{Th}}$ is found by applying a
    test voltage and measuring the resulting current.
\end{itemize}
\end{theorem}

\begin{theorem}[Norton]
Any linear two-terminal network can be replaced, at its accessible
terminals, by an equivalent current source $I_{\text{N}}$ in parallel
with a resistance $R_{\text{N}} = R_{\text{Th}}$. The Norton current
is the short-circuit current at the terminals: $I_{\text{N}} =
V_{\text{Th}}/R_{\text{Th}}$.
\end{theorem}

% Figure: Thevenin equivalent (left) and Norton equivalent (right) shown side
% by side, illustrating source transformation. Plain TikZ.
\begin{figure}[htbp]
\centering
\begin{tikzpicture}[
  line width=0.8pt,
  res/.style={draw, fill=white, minimum width=12mm, minimum height=4.5mm, font=\scriptsize},
  font=\small,
]
% --- Thevenin (left) ---
\begin{scope}[xshift=0cm]
  \draw[engblue] (0,0) -- (0,1.2);
  \node[draw, circle, minimum size=8mm, fill=white] (Vth) at (0,1.8) {};
  \node[font=\scriptsize] at (0,1.8) {$V_{\text{Th}}$};
  \node[font=\scriptsize, anchor=east] at (-0.45,2.3) {$+$};
  \node[font=\scriptsize, anchor=east] at (-0.45,1.3) {$-$};
  \draw[engblue] (0,2.4) -- (0,3.2);
  \node[res] (Rth) at (1.4,3.2) {$R_{\text{Th}}$};
  \draw[engblue] (0,3.2) -- (Rth) -- (2.8,3.2) -- (2.8,2.6);
  \draw[engblue] (0,0) -- (2.8,0) -- (2.8,0.6);
  % open terminals a, b
  \node[circle, fill=engblack, inner sep=1pt] (a1) at (2.8,2.6) {};
  \node[circle, fill=engblack, inner sep=1pt] (b1) at (2.8,0.6) {};
  \node[font=\scriptsize, anchor=west] at (2.95,2.6) {$a$};
  \node[font=\scriptsize, anchor=west] at (2.95,0.6) {$b$};
  \node[font=\small] at (1.4,-0.7) {Th\'evenin};
\end{scope}
% --- Norton (right) ---
\begin{scope}[xshift=6cm]
  \node[draw, circle, minimum size=8mm, fill=white] (In) at (0,1.6) {};
  \node[font=\scriptsize] at (0,1.6) {$I_{\text{N}}$};
  \draw[engorange, -{Stealth}] (0,1.1) -- (0,2.1);
  \draw[engblue] (0,2.0) -- (0,3.2) -- (2.8,3.2) -- (2.8,2.6);
  \draw[engblue] (0,1.2) -- (0,0) -- (2.8,0) -- (2.8,0.6);
  \node[res, rotate=90] (Rn) at (1.4,1.6) {};
  \node[font=\scriptsize, anchor=south] at (1.4,2.15) {$R_{\text{N}}$};
  \draw[engblue] (1.4,3.2) -- (Rn);
  \draw[engblue] (Rn) -- (1.4,0);
  \node[circle, fill=engblack, inner sep=1pt] (a2) at (2.8,2.6) {};
  \node[circle, fill=engblack, inner sep=1pt] (b2) at (2.8,0.6) {};
  \node[font=\scriptsize, anchor=west] at (2.95,2.6) {$a$};
  \node[font=\scriptsize, anchor=west] at (2.95,0.6) {$b$};
  \node[font=\small] at (1.4,-0.7) {Norton};
\end{scope}
\end{tikzpicture}
\caption[Th\'evenin and Norton equivalents]{The two equivalent
representations of any linear two-terminal network. The Th\'evenin form is a
voltage source $V_{\text{Th}}$ in series with $R_{\text{Th}}$; the Norton
form is a current source $I_{\text{N}}$ in parallel with
$R_{\text{N}} = R_{\text{Th}}$. Source transformation converts one to the
other through $I_{\text{N}} = V_{\text{Th}}/R_{\text{Th}}$. At terminals
$a$-$b$ the two networks are indistinguishable to any external load.}
\label{fig:vol04ch09:thevenin-norton}
\end{figure}


Th\'evenin and Norton equivalents are interconvertible by source
transformation. The pair is the working language of interface
specifications: when a designer says ``the source impedance is
$50\,\si{\ohm}$,'' the speaker means $R_{\text{Th}} = 50\,\si{\ohm}$
at the relevant frequency. When the speaker says ``the load is a
$1\,\si{\kilo\ohm}$ input,'' the listener treats it as the load
resistance into which the Th\'evenin source delivers power.

The \engterm{maximum power transfer theorem} follows from Th\'evenin:
a source with open-circuit voltage $V_{\text{Th}}$ and source
resistance $R_{\text{Th}}$ delivers maximum power to a load $R_{L}$
when $R_{L} = R_{\text{Th}}$, with the maximum power equal to
$V_{\text{Th}}^{2}/(4R_{\text{Th}})$. The matched-source efficiency is
$50\,\%$, which is acceptable for signal transmission (RF, audio) but
unacceptable for power delivery (where one designs for $R_{L} \gg
R_{\text{Th}}$ to keep the efficiency close to $100\,\%$).

\begin{archetype}[Balance and transport]
A linear circuit is two patterns nested. The balance pattern enforces
KCL at every node and KVL around every loop, giving the topological
structure. The transport pattern is the constitutive equation of each
two-terminal device: resistors transport instantly ($v = iR$),
capacitors transport with a time derivative ($i = C\,\dd v/\dd t$),
inductors transport with the other time derivative ($v = L\,\dd i/
\dd t$). Every circuit problem decomposes into the two patterns; an
engineer who keeps them separate in the bookkeeping makes fewer sign
errors.
\end{archetype}

\section{Capacitors and inductors}\pathtag{core}

A capacitor stores energy in its electric field. The constitutive
equation $i = C\,\dd v/\dd t$ implies that the voltage across a
capacitor cannot change instantaneously without infinite current; in
practice, capacitor voltage is the slow state variable of a circuit
containing a capacitor. The energy stored at voltage $v$ is
$U_{C} = \tfrac{1}{2}Cv^{2}$, the integral of $vi$ from $v = 0$ to the
present voltage. Real capacitors have parasitic resistance
(\engterm{equivalent series resistance}, ESR) that dissipates power
during charging and discharging, and parasitic inductance (ESL) that
limits high-frequency performance.

An inductor stores energy in its magnetic field. The constitutive
equation $v = L\,\dd i/\dd t$ implies that the current through an
inductor cannot change instantaneously without infinite voltage; the
current is the slow state variable. The energy stored at current $i$
is $U_{L} = \tfrac{1}{2}Li^{2}$. Real inductors have parasitic
resistance (winding resistance) and parasitic capacitance (between
turns of the winding) that limits their useful frequency range to
below the self-resonant frequency.

Two capacitors in parallel sum: $C_{\text{tot}} = C_{1}+C_{2}$. Two
capacitors in series invert and add: $1/C_{\text{tot}} = 1/C_{1} +
1/C_{2}$. Two inductors in series sum (assuming no mutual coupling):
$L_{\text{tot}} = L_{1}+L_{2}$; two in parallel invert and add. The
rules are the duals of the resistor rules: resistors in series sum,
in parallel invert and add.

\begin{principle}[Pick the state variable]
A circuit's order is the number of independent capacitors and
inductors. Capacitor voltage and inductor current are the state
variables. A circuit with one capacitor or one inductor is first-order
and has an exponential response with a single time constant. A circuit
with one capacitor and one inductor (or two of either, configured to
be independent) is second-order with one of three responses:
overdamped, critically damped, underdamped. A circuit with $n$
states has $n$th-order response described by an $n$th-order linear
ODE. State-space methods of Volume IX systematise this; the procedure
here is the entry point.
\end{principle}

The op-amp introduces \engterm{dependent sources} into the working
vocabulary. An ideal op-amp has infinite input impedance, zero output
impedance, and infinite open-loop gain; its working tools are the two
``virtual short'' constraints in negative-feedback operation:
$v_{+} = v_{-}$ (the two input terminals are at the same voltage) and
$i_{+} = i_{-} = 0$ (no current flows into the input terminals).
Under these constraints, the working op-amp circuits (inverting
amplifier, non-inverting amplifier, summing amplifier, integrator,
differentiator, instrumentation amplifier, active filter) follow from
nodal analysis applied at the inverting input. The detailed small-
signal model is in Sze \cite{text:sze-semiconductor} or Alexander-
Sadiku \cite{text:alexander-sadiku-circuits}; for this chapter the
virtual-short approximation is the working tool.

% Figure: two op-amp configurations, inverting (left) and non-inverting
% (right), drawn with a plain TikZ triangle for the amplifier symbol.
\begin{figure}[htbp]
\centering
\begin{tikzpicture}[
  line width=0.8pt,
  res/.style={draw, fill=white, minimum width=11mm, minimum height=4mm, font=\scriptsize},
  font=\small,
]
% --- Inverting amplifier (left) ---
\begin{scope}[xshift=0cm]
  % op-amp triangle
  \draw[engblue] (1.6,0.6) -- (1.6,2.4) -- (3.4,1.5) -- cycle;
  \node[font=\scriptsize, anchor=west] at (1.7,2.0) {$-$};
  \node[font=\scriptsize, anchor=west] at (1.7,1.0) {$+$};
  % input resistor Rin
  \node[res] (Rin) at (0.4,2.0) {$R_{\text{in}}$};
  \draw[engblue] (-0.8,2.0) -- (Rin) -- (1.6,2.0);
  \node[font=\scriptsize, anchor=east] at (-0.85,2.0) {$v_{\text{in}}$};
  % feedback resistor Rf
  \node[res] (Rf) at (2.5,3.2) {$R_f$};
  \draw[engblue] (1.6,2.0) -- (1.6,3.2) -- (Rf) -- (3.9,3.2) -- (3.9,1.5);
  % output
  \draw[engblue] (3.4,1.5) -- (4.4,1.5);
  \node[font=\scriptsize, anchor=west] at (4.45,1.5) {$v_{\text{out}}$};
  % non-inverting input grounded
  \draw[engblue] (1.6,1.0) -- (1.0,1.0) -- (1.0,0.4);
  \node[font=\scriptsize] at (1.0,0.15) {gnd};
  \node[font=\small] at (1.8,-0.4) {inverting, gain $-R_f/R_{\text{in}}$};
\end{scope}
% --- Non-inverting amplifier (right) ---
\begin{scope}[xshift=7cm]
  \draw[engblue] (1.6,0.6) -- (1.6,2.4) -- (3.4,1.5) -- cycle;
  \node[font=\scriptsize, anchor=west] at (1.7,2.0) {$-$};
  \node[font=\scriptsize, anchor=west] at (1.7,1.0) {$+$};
  % input to non-inverting terminal
  \draw[engblue] (0.4,1.0) -- (1.6,1.0);
  \node[font=\scriptsize, anchor=east] at (0.35,1.0) {$v_{\text{in}}$};
  % feedback divider Rf and Rg
  \node[res] (Rf2) at (2.5,3.2) {$R_f$};
  \draw[engblue] (1.6,2.0) -- (1.6,3.2) -- (Rf2) -- (3.9,3.2) -- (3.9,1.5);
  \node[res, rotate=90] (Rg) at (1.6,2.9) {};
  \node[font=\scriptsize, anchor=east] at (1.35,2.9) {$R_g$};
  \draw[engblue] (1.6,2.0) -- (Rg);
  \draw[engblue] (Rg) -- (1.6,3.5) -- (1.0,3.5) -- (1.0,3.9);
  \node[font=\scriptsize] at (1.0,4.15) {gnd};
  % output
  \draw[engblue] (3.4,1.5) -- (4.4,1.5);
  \node[font=\scriptsize, anchor=west] at (4.45,1.5) {$v_{\text{out}}$};
  \node[font=\small] at (1.8,-0.4) {non-inverting, gain $1+R_f/R_g$};
\end{scope}
\end{tikzpicture}
\caption[Inverting and non-inverting op-amp configurations]{The two basic
negative-feedback op-amp amplifiers. Both follow from the virtual-short
constraints of the ideal op-amp: the two input terminals sit at the same
voltage and draw no current. The inverting configuration has gain
$-R_f/R_{\text{in}}$ and input impedance $R_{\text{in}}$; the non-inverting
configuration has gain $1 + R_f/R_g$ and very high input impedance. The
choice between them turns on the sign of the gain and on the source impedance
the amplifier must present to its driver.}
\label{fig:vol04ch09:opamp-configs}
\end{figure}


Figure~\ref{fig:vol04ch09:opamp-configs} draws the two basic
negative-feedback configurations. The inverting amplifier has gain
$-R_f/R_{\text{in}}$ and presents input impedance $R_{\text{in}}$; the
non-inverting amplifier has gain $1+R_f/R_g$ and presents very high
input impedance. Both follow from the virtual short by writing KCL at
the inverting input.

\subsection{Feedback gain, loop gain, and the gain-bandwidth product}
\pathtag{standard}

The virtual short is the limit of a quantitative statement worth
carrying in full, because it tells the engineer how much error the
finite gain of a real op-amp leaves behind and where the closed-loop
bandwidth lands. A real op-amp produces an output $v_o = A(v_+ - v_-)$
with a large but finite open-loop gain $A$. The non-inverting amplifier
feeds a fraction $\beta = R_g/(R_f + R_g)$ of the output back to the
inverting input, so $v_- = \beta v_o$ and $v_+ = v_{\text{in}}$. Solving
$v_o = A(v_{\text{in}} - \beta v_o)$ for the closed-loop gain gives
\[
G = \frac{v_o}{v_{\text{in}}} = \frac{A}{1 + A\beta}
= \frac{1}{\beta}\cdot\frac{A\beta}{1 + A\beta}.
\]
The product $A\beta$ is the \engterm{loop gain}, the gain a signal sees
travelling once around the feedback loop, and $1 + A\beta$ is the
\engterm{return difference}. When the loop gain is large, $A\beta \gg
1$, the closed-loop gain collapses to $G \to 1/\beta = 1 + R_f/R_g$, the
ideal result, set by the resistor ratio alone and independent of the
op-amp. The fractional error of the ideal formula is $1/(1 + A\beta)$:
an op-amp with $A = 10^{5}$ in a circuit with $\beta = 1/100$ has a loop
gain of $10^{3}$, so the gain is low by about $0.1\,\%$, and the
$0.1\,\%$ shrinks further as $A$ rises. The same return difference
divides every imperfection referred to the output: distortion,
output-impedance variation, and the effect of load changes all fall by
the factor $1 + A\beta$. Feedback buys accuracy by trading away the
surplus open-loop gain.

The open-loop gain is not constant with frequency. A general-purpose
op-amp is deliberately compensated to roll off as a single pole,
\[
A(\omega) = \frac{A_0}{1 + j\omega/\omega_p},
\]
with a DC gain $A_0$ and a dominant pole at $\omega_p$ placed low (often
a few hertz to a few tens of hertz) so that the phase stays well short
of $180^\circ$ until the gain has fallen below unity, the condition for
stable feedback. Substituting this $A(\omega)$ into the closed-loop
expression and keeping $\beta$ frequency-independent,
\[
G(\omega) = \frac{A_0/(1 + j\omega/\omega_p)}
{1 + \beta A_0/(1 + j\omega/\omega_p)}
= \frac{A_0}{(1 + \beta A_0) + j\omega/\omega_p}.
\]
Dividing through by $(1 + \beta A_0)$ puts this in single-pole form with
DC gain $G_0 = A_0/(1 + \beta A_0) \approx 1/\beta$ and a closed-loop
corner at
\[
\omega_c = \omega_p(1 + \beta A_0) \approx \beta A_0\,\omega_p.
\]
The product of the closed-loop DC gain and the closed-loop bandwidth is
therefore $G_0\,\omega_c \approx (1/\beta)(\beta A_0\omega_p) =
A_0\omega_p$, a constant independent of $\beta$. That constant is the
\engterm{gain-bandwidth product} (GBW), the frequency at which the
open-loop gain itself falls to unity. The working rule drops out: a
feedback amplifier set to a closed-loop gain $G$ has a bandwidth
$\text{GBW}/G$, so raising the gain narrows the bandwidth in exact
proportion. Figure~\ref{fig:vol04ch09:gbw} draws the open-loop roll-off
and two closed-loop gain lines breaking where they meet it.

% Figure: op-amp open-loop gain rolling off at -20 dB/dec from a low-frequency
% plateau, intersecting two closed-loop gain lines at the gain-bandwidth
% product. Illustrates the constant gain-bandwidth product of a
% single-pole-compensated op-amp. pgfplots semilog.
\begin{figure}[htbp]
\centering
\begin{tikzpicture}
\begin{semilogxaxis}[
  width=11.5cm, height=5.4cm,
  xlabel={frequency $f$ (Hz)},
  ylabel={gain (dB)},
  xmin=1, xmax=10000000, ymin=-10, ymax=130,
  axis lines=left,
  tick label style={font=\tiny},
  label style={font=\scriptsize},
  ytick={0,20,40,60,80,100,120},
  clip=false,
]
% open-loop: A0 = 120 dB (1e6), dominant pole at fp = 10 Hz, unity-gain at
% GBW = A0 * fp = 1e7 Hz. |A| dB = 20 log10( A0 / sqrt(1+(f/fp)^2) )
% = 120 - 10 log10(1 + (f/10)^2)
\addplot[engblue, very thick, domain=1:10000000, samples=400]
  {120 - 10*log10(1 + (x/10)^2)};
% closed-loop gain 40 dB (=100): flat to fc1 = GBW/100 = 1e5 Hz
\addplot[engred, thick] coordinates {(1,40) (100000,40)};
\addplot[engred, thick, dashed, domain=100000:10000000, samples=2]
  {40 - 20*log10(x/100000)};
\node[engred, font=\tiny, anchor=south west] at (axis cs:1.5,41)
  {closed loop, $G=100$};
% closed-loop gain 20 dB (=10): flat to fc2 = GBW/10 = 1e6 Hz
\addplot[engorange, thick] coordinates {(1,20) (1000000,20)};
\addplot[engorange, thick, dashed, domain=1000000:10000000, samples=2]
  {20 - 20*log10(x/1000000)};
\node[engorange, font=\tiny, anchor=south west] at (axis cs:1.5,21)
  {closed loop, $G=10$};
% GBW marker at unity gain crossing
\addplot[enggray, dotted] coordinates {(10000000,-10) (10000000,130)};
\node[enggray, font=\tiny, anchor=south east, align=right] at (axis cs:10000000,0)
  {GBW\\$=10\,$MHz};
\node[engblue, font=\tiny, anchor=west] at (axis cs:1.5,118) {open loop};
\end{semilogxaxis}
\end{tikzpicture}
\caption[Gain-bandwidth product of a compensated op-amp]{The open-loop gain
of a single-pole-compensated op-amp (blue) starts at its DC value, breaks at
the dominant pole, and falls at $-20\,\si{\decibel}$ per decade to the
unity-gain frequency, the gain-bandwidth product. A feedback network that
sets a closed-loop gain holds that gain flat until the closed-loop line meets
the open-loop roll-off, where the amplifier runs out of loop gain and the
response breaks. A gain of one hundred breaks at one hundredth of the
gain-bandwidth product; a gain of ten breaks a decade higher. The product of
closed-loop gain and closed-loop bandwidth is the constant the data sheet
calls GBW.}
\label{fig:vol04ch09:gbw}
\end{figure}


The inverting amplifier obeys the same arithmetic with one wrinkle. Its
ideal gain is $-R_f/R_{\text{in}}$, but the bandwidth is governed by the
\engterm{noise gain}, the gain the op-amp's own input error sees, which
is $1 + R_f/R_{\text{in}}$ rather than $R_f/R_{\text{in}}$. An inverting
stage of gain $-10$ has a noise gain of $11$, so its bandwidth is
$\text{GBW}/11$, not $\text{GBW}/10$, and its output noise and offset
are amplified by $11$. The distinction between signal gain and noise
gain is the kind of detail that separates a circuit that meets its
bandwidth spec on the bench from one that falls a few percent short, and
it is the reason the design exercises in this chapter size the op-amp by
its noise gain.

\section{Transient response: RC, RL, RLC}\pathtag{core}
\label{sec:vol04ch09:transient}

A first-order RC circuit with a resistor $R$ in series with a
capacitor $C$, driven by a voltage source $V_{s}$, obeys
\[
RC\dv{v_{C}}{t} + v_{C} = V_{s}.
\]
The solution with initial condition $v_{C}(0) = v_{0}$ is
\[
v_{C}(t) = V_{s} + (v_{0}-V_{s})e^{-t/\tau}, \quad \tau = RC.
\]
The time constant $\tau$ has units of seconds. After $\tau$ the
voltage has moved $63\,\%$ of the way from $v_{0}$ to $V_{s}$; after
$3\tau$, $95\,\%$; after $5\tau$, $99.3\,\%$. The practising-engineer
rule of thumb is that the circuit has settled to the new steady state
after $5\tau$.

% Figure: RC step response, charging curve approaching V_s with the 63/95/99
% percent marks at one, three, and five time constants. pgfplots.
\begin{figure}[htbp]
\centering
\begin{tikzpicture}
\begin{axis}[
  width=12cm, height=7.5cm,
  xlabel={time $t/\tau$}, ylabel={capacitor voltage $v_C/V_s$},
  xmin=0, xmax=6, ymin=0, ymax=1.1,
  axis lines=left,
  tick label style={font=\scriptsize},
  label style={font=\small},
  legend pos=south east,
  legend style={font=\scriptsize},
  clip=false,
]
% charging curve 1 - exp(-t)
\addplot[engblue, very thick, domain=0:6, samples=120] {1 - exp(-x)};
\addlegendentry{$v_C(t) = V_s(1 - e^{-t/\tau})$}
% asymptote
\addplot[enggray, dashed] coordinates {(0,1) (6,1)};
% one tau: 0.632
\draw[engred, dashed] (axis cs:1,0) -- (axis cs:1,0.632);
\draw[engred, dashed] (axis cs:0,0.632) -- (axis cs:1,0.632);
\filldraw[engred] (axis cs:1,0.632) circle [radius=1.6pt];
\node[engred, font=\scriptsize, anchor=north west] at (axis cs:1.05,0.60) {$63\%$ at $\tau$};
% three tau: 0.950
\filldraw[engorange] (axis cs:3,0.950) circle [radius=1.6pt];
\node[engorange, font=\scriptsize, anchor=north west] at (axis cs:3.05,0.92) {$95\%$ at $3\tau$};
% five tau: 0.993
\filldraw[englime!60!black] (axis cs:5,0.993) circle [radius=1.6pt];
\node[font=\scriptsize, anchor=south] at (axis cs:5,1.0) {$99.3\%$ at $5\tau$};
\end{axis}
\end{tikzpicture}
\caption[RC charging transient]{Step response of a first-order RC circuit.
The capacitor voltage rises exponentially toward the source voltage with
time constant $\tau = RC$. The working settling marks are $63\%$ after one
time constant, $95\%$ after three, and $99.3\%$ after five. The
engineer's rule that the circuit has reached steady state after $5\tau$
follows directly. The discharge curve is the mirror image,
$v_C = V_s e^{-t/\tau}$.}
\label{fig:vol04ch09:rc-transient}
\end{figure}


Figure~\ref{fig:vol04ch09:rc-transient} plots the charging curve with
the three settling marks. The discharge curve is the mirror image,
falling from the initial voltage toward zero with the same time
constant.

An RL circuit with $R$ in series with $L$ obeys
\[
\frac{L}{R}\dv{i_{L}}{t} + i_{L} = \frac{V_{s}}{R},
\]
with time constant $\tau = L/R$. The duality of capacitor voltage
with inductor current is general: every result for the RC circuit
maps to the RL circuit under the substitution $v \leftrightarrow i$,
$C \leftrightarrow L$, $R \leftrightarrow 1/R$.

A second-order RLC series circuit with source $V_{s}$ driving $R$,
$L$, $C$ in series obeys
\[
LC\frac{\dd^{2} v_{C}}{\dd t^{2}} + RC\dv{v_{C}}{t} + v_{C} = V_{s},
\]
or in standard form
\[
\frac{\dd^{2} v_{C}}{\dd t^{2}} + 2\zeta\omega_{n}\dv{v_{C}}{t}
+ \omega_{n}^{2} v_{C} = \omega_{n}^{2}V_{s},
\]
with \engterm{natural frequency} $\omega_{n} = 1/\sqrt{LC}$ and
\engterm{damping ratio} $\zeta = (R/2)\sqrt{C/L}$. Three regimes
follow from the roots of the characteristic polynomial
$s^{2}+2\zeta\omega_{n}s+\omega_{n}^{2}=0$:
\begin{itemize}[leftmargin=*]
  \item \engterm{Overdamped} $\zeta>1$: two real negative roots; the
    response is the sum of two decaying exponentials.
  \item \engterm{Critically damped} $\zeta=1$: a repeated real root;
    the response is the fastest decay without overshoot.
  \item \engterm{Underdamped} $\zeta<1$: complex-conjugate roots; the
    response oscillates at frequency $\omega_{d} = \omega_{n}\sqrt{
    1-\zeta^{2}}$ and decays at envelope rate $\zeta\omega_{n}$.
\end{itemize}
The underdamped response with low damping (high $Q = 1/(2\zeta)$ in
filter terminology) is the working ringing seen in oscilloscope traces
of digital signals that overshoot their target voltage. Critical
damping is the design target for control systems that need fast
response without oscillation.

The unit-step response of an underdamped second-order system with
$\zeta=0.3$, $\omega_{n}=1$ overshoots by about $37\,\%$, settles to
within $5\,\%$ of its final value in about $9$ time units, and
oscillates at $\omega_{d}\approx 0.95$. The same dynamics appear in
mechanical systems (Volume III, Chapter 12) and in control (Volume
IX); the transfer-function language unifies them.

% Figure: unit-step response of a second-order system in three damping
% regimes (underdamped, critically damped, overdamped). pgfplots.
% Underdamped zeta=0.3, wn=1: y = 1 - exp(-0.3 t)/sqrt(1-0.09) sin(wd t + phi)
% with wd = sqrt(1-0.09)=0.954, phi = acos(zeta)=72.54 deg.
\begin{figure}[htbp]
\centering
\begin{tikzpicture}
\begin{axis}[
  width=12cm, height=7.5cm,
  xlabel={time $\omega_n t$}, ylabel={step response $v_C/V_s$},
  xmin=0, xmax=14, ymin=0, ymax=1.45,
  axis lines=left,
  tick label style={font=\scriptsize},
  label style={font=\small},
  legend pos=north east,
  legend style={font=\scriptsize},
  clip=false,
]
% final value
\addplot[enggray, dashed] coordinates {(0,1) (14,1)};
% underdamped zeta = 0.3
\addplot[engred, very thick, domain=0:14, samples=200]
  {1 - exp(-0.3*x)/0.9539*sin(deg(0.9539*x) + 72.54)};
\addlegendentry{$\zeta=0.3$ underdamped}
% critically damped zeta = 1: y = 1 - (1+t) e^{-t}
\addplot[engblue, very thick, domain=0:14, samples=200]
  {1 - (1+x)*exp(-x)};
\addlegendentry{$\zeta=1$ critically damped}
% overdamped zeta = 2: roots -0.2679, -3.7321; y = 1 - (s2 e^{s1 t} - s1 e^{s2 t})/(s2 - s1)
% s1=-0.2679, s2=-3.7321
\addplot[englime!60!black, very thick, domain=0:14, samples=200]
  {1 - (-3.7321*exp(-0.2679*x) + 0.2679*exp(-3.7321*x))/(-3.7321 + 0.2679)};
\addlegendentry{$\zeta=2$ overdamped}
\end{axis}
\end{tikzpicture}
\caption[Second-order step response in three damping regimes]{Unit-step
response of a series RLC circuit (capacitor voltage) at three damping ratios.
The underdamped case ($\zeta=0.3$) overshoots the final value and rings at
$\omega_d = \omega_n\sqrt{1-\zeta^2}$ before settling. The critically damped
case ($\zeta=1$) is the fastest approach without overshoot. The overdamped
case ($\zeta=2$) approaches more slowly as a sum of two decaying
exponentials. Control design that needs speed without ringing targets
critical damping; high-$Q$ tuned circuits live deep in the underdamped
regime.}
\label{fig:vol04ch09:rlc-damping}
\end{figure}


Figure~\ref{fig:vol04ch09:rlc-damping} plots the step response at three
damping ratios. The underdamped curve overshoots and rings; the
critically damped curve is the fastest monotone approach; the
overdamped curve approaches more slowly as the sum of two exponentials.

For the input ramp, sinusoid, and impulse the response is found by
the same characteristic-polynomial method or by Laplace transform;
the Laplace approach is developed in Volume II, Chapter 18 and is the
natural language of linear time-invariant analysis.

\begin{estimation}
\textbf{Question.} A camera flash uses a $1000\,\si{\micro\farad}$
capacitor charged to $300\,\si{\volt}$ and discharged through a xenon
flash tube of effective resistance $\sim 1\,\si{\ohm}$ during the
discharge. Estimate the peak discharge current and the duration of
the flash; estimate the energy delivered to the flash and the peak
electrical power.

\textbf{Estimate, before reading on.} State a factor-of-three-
confidence answer for each quantity before reading the working.

\textbf{Working.} The energy stored is $U = \tfrac{1}{2}CV^{2} =
0.5\times 10^{-3}\times 9\times 10^{4} = 45\,\si{\joule}$. The
discharge through a $1\,\si{\ohm}$ resistor has time constant
$\tau = RC = 1\times 10^{-3} = 1\,\si{\milli\second}$; the flash
duration is several time constants, of order $1$ to $3\,\si{\milli
\second}$, consistent with published xenon-flash specifications. The
peak current is $V_{0}/R = 300\,\si{\ampere}$ initially, falling
exponentially. The peak electrical power is $V_{0}^{2}/R = 9\times
10^{4}\,\si{\watt} = 90\,\si{\kilo\watt}$.

\textbf{Postmortem.} The xenon arc's V-I characteristic is far from
linear: it ignites at high voltage, drops to a low arc voltage during
conduction, and extinguishes when the current falls below a critical
value. The $1\,\si{\ohm}$ effective resistance is a rough average
over the discharge; the actual peak current can exceed $1\,\si{\kilo
\ampere}$ at the moment of ignition before settling to the arc
regime. A reader who guessed within a factor of three on energy and
flash duration has internalised the time scale; the peak current is
harder to estimate without the tube's V-I curve.
\end{estimation}

\section{AC steady state and impedance}\pathtag{core}
\label{sec:vol04ch09:ac}

In sinusoidal steady state, every voltage and current in a linear
circuit is a sinusoid of the same frequency $\omega$. The phasor
notation represents a sinusoid $v(t) = V_{m}\cos(\omega t+\phi)$ by
the complex amplitude
\[
\tilde V = V_{m}e^{j\phi}, \quad v(t) = \Re(\tilde V e^{j\omega t}).
\]
Differentiation in the time domain becomes multiplication by $j\omega$
in the phasor domain; integration becomes division. The constitutive
equations of the three energy-storage elements become algebraic:
\begin{itemize}[leftmargin=*]
  \item Resistor: $\tilde V = R\tilde I$, impedance $Z_{R} = R$.
  \item Capacitor: $\tilde V = (1/j\omega C)\tilde I$, impedance
    $Z_{C} = 1/(j\omega C)$.
  \item Inductor: $\tilde V = j\omega L\,\tilde I$, impedance
    $Z_{L} = j\omega L$.
\end{itemize}
\engterm{Impedance} $Z = R+jX$ generalises resistance to AC; its real
part $R$ is the resistance, its imaginary part $X$ is the
\engterm{reactance}. \engterm{Admittance} $Y = 1/Z = G+jB$ is its
reciprocal; its real part $G$ is the conductance, its imaginary part
$B$ is the susceptance. KVL and KCL apply to phasor voltages and
currents unchanged; Th\'evenin, Norton, and superposition apply in
phasor form. The DC machinery transports to AC at the price of
working with complex numbers instead of real ones.

% Figure: phasor diagram for an RL series load showing V, I, and the phase
% angle phi by which voltage leads current. Plain TikZ vectors.
\begin{figure}[htbp]
\centering
\begin{tikzpicture}[
  line width=0.9pt,
  font=\small,
  >={Stealth[length=2.5mm]},
]
% axes
\draw[enggray, ->] (-0.5,0) -- (5,0) node[anchor=west, font=\scriptsize] {Re};
\draw[enggray, ->] (0,-0.5) -- (0,3.5) node[anchor=south, font=\scriptsize] {Im};
\coordinate (O) at (0,0);
% current phasor along the real axis (reference), magnitude 4
\draw[engblue, ->, very thick] (O) -- (4,0);
\node[engblue, anchor=north] at (4,-0.1) {$\tilde I$};
% voltage phasor leading by phi = 37 deg, magnitude ~ 4 / cos(37) ~ 5
\draw[engred, ->, very thick] (O) -- (37:5);
\node[engred, anchor=south west] at (37:5) {$\tilde V$};
% resistive component V_R along I
\draw[engorange, ->] (O) -- (4,0);
% reactive component V_L vertical (from tip of V_R to tip of V)
\draw[engorange, ->] (4,0) -- (4,3.0);
\node[engorange, font=\scriptsize, anchor=west] at (4.05,1.5) {$\tilde V_L = j\omega L\,\tilde I$};
\node[engorange, font=\scriptsize, anchor=north] at (2,-0.3) {$\tilde V_R = R\,\tilde I$};
% angle arc
\draw[engblack] (1.2,0) arc (0:37:1.2);
\node[font=\scriptsize] at (1.6,0.45) {$\phi$};
\end{tikzpicture}
\caption[Phasor diagram of a series RL load]{Phasor diagram for a series RL
load with the current taken as the reference along the real axis. The
resistor voltage $\tilde V_R = R\tilde I$ is in phase with the current; the
inductor voltage $\tilde V_L = j\omega L\,\tilde I$ leads by $90^\circ$. The
total voltage $\tilde V = \tilde V_R + \tilde V_L$ leads the current by the
impedance angle $\phi = \arctan(\omega L/R)$. The power factor is
$\cos\phi$, and the load is inductive ($\phi > 0$), so it absorbs reactive
power.}
\label{fig:vol04ch09:phasor}
\end{figure}


The phasor diagram in Figure~\ref{fig:vol04ch09:phasor} draws the
voltage and current of a series RL load as rotating vectors frozen at
$t=0$. The resistor voltage sits in phase with the current; the
inductor voltage leads it by $90^\circ$; their sum leads the current by
the impedance angle $\phi$.

Impedances in series add: $Z_{\text{tot}} = Z_{1}+Z_{2}$. Impedances in
parallel invert and add: $1/Z_{\text{tot}} = 1/Z_{1}+1/Z_{2}$. The RLC
series impedance is $Z(\omega) = R + j\omega L + 1/(j\omega C) = R +
j(\omega L - 1/\omega C)$, which has zero reactance at the
\engterm{resonant frequency} $\omega_{0} = 1/\sqrt{LC}$. At resonance
the circuit is purely resistive; below resonance it is capacitive;
above resonance it is inductive.

The \engterm{transfer function} $H(\omega) = V_{\text{out}}/V_{\text{in}}$
of a filter is the phasor ratio of output to input. A low-pass first-
order filter (single RC stage) has $H(\omega) = 1/(1+j\omega RC)$; a
high-pass single-stage filter has $H(\omega) = j\omega RC/(1+j\omega
RC)$. A bandpass filter is the cascade of high-pass and low-pass with
non-overlapping cutoffs, or an RLC tuned circuit. The
\engterm{quality factor} $Q$ of a tuned circuit is $\omega_{0}L/R$
for a series RLC, $R/(\omega_{0}L)$ for a parallel RLC; $Q$ governs
the bandwidth ($\Delta\omega = \omega_{0}/Q$) and the peak amplitude
(peak amplification of $Q$ at resonance for an underdamped system).

A Bode plot is a plot of $20\log_{10}|H(\omega)|$ in decibels and
$\angle H(\omega)$ in degrees against $\log\omega$. The Bode plot is
the working representation of a filter's frequency response;
asymptotic-construction rules give an approximate Bode plot from the
filter's pole-zero locations, and the working engineer carries the
$-20\,\si{\decibel}$-per-decade-per-pole and $+20\,\si{\decibel}$-per-
decade-per-zero conventions to bench measurements.

% Figure: Bode magnitude plot of a first-order RC low-pass filter, with the
% -20 dB/decade asymptote and the -3 dB cutoff marked. pgfplots, semilog x.
% H = 1/sqrt(1 + (f/fc)^2), fc = 1 (normalised). dB = 20 log10|H|.
\begin{figure}[htbp]
\centering
\begin{tikzpicture}
\begin{semilogxaxis}[
  width=12cm, height=7cm,
  xlabel={frequency $f/f_c$}, ylabel={magnitude $20\log_{10}|H|$ (dB)},
  xmin=0.01, xmax=100, ymin=-45, ymax=5,
  axis lines=left,
  tick label style={font=\scriptsize},
  label style={font=\small},
  legend pos=south west,
  legend style={font=\scriptsize},
  clip=false,
]
% exact magnitude in dB
\addplot[engblue, very thick, domain=0.01:100, samples=200]
  {10*log10(1/(1 + x^2))};
\addlegendentry{exact $|H|$}
% low-frequency asymptote (0 dB)
\addplot[enggray, dashed] coordinates {(0.01,0) (1,0)};
% high-frequency asymptote -20 dB/decade
\addplot[engred, dashed, domain=1:100, samples=2]
  {-20*log10(x)};
\addlegendentry{$-20$ dB/decade asymptote}
% -3 dB point at f = fc
\filldraw[engorange] (axis cs:1,-3.0103) circle [radius=1.8pt];
\node[engorange, font=\scriptsize, anchor=west] at (axis cs:1.1,-3) {$-3$ dB at $f_c$};
\end{semilogxaxis}
\end{tikzpicture}
\caption[Bode magnitude plot of a first-order low-pass filter]{Bode
magnitude response of a first-order RC low-pass filter with cutoff $f_c =
1/(2\pi RC)$. The pass band is flat at $0$ dB; the response falls at
$-20$ dB per decade above the cutoff. The two straight-line asymptotes meet
at the corner frequency, where the exact response is $-3$ dB (a factor of
$1/\sqrt{2}$ in amplitude). The asymptotic construction lets the engineer
sketch the response of any pole-zero filter by inspection before refining it
on the bench.}
\label{fig:vol04ch09:bode-lowpass}
\end{figure}


Figure~\ref{fig:vol04ch09:bode-lowpass} shows the Bode magnitude of a
first-order low-pass filter, with the two asymptotes meeting at the
corner frequency where the exact response is $-3\,\si{\decibel}$.
Figure~\ref{fig:vol04ch09:bandpass-resonance} shows the magnitude of a
series RLC bandpass filter at two quality factors; the higher-$Q$ curve
is the narrower one, and its $-3\,\si{\decibel}$ bandwidth is
$\Delta f = f_0/Q$.

% Figure: magnitude response of a series-RLC bandpass filter at two quality
% factors, showing the resonant peak and the bandwidth. pgfplots semilog x.
% |H| = (f/(Q f0)) / sqrt( (1 - (f/f0)^2)^2 + (f/(Q f0))^2 ), normalised f0=1.
\begin{figure}[htbp]
\centering
\begin{tikzpicture}
\begin{semilogxaxis}[
  width=12cm, height=7cm,
  xlabel={frequency $f/f_0$}, ylabel={$|H(f)|$},
  xmin=0.1, xmax=10, ymin=0, ymax=1.1,
  axis lines=left,
  tick label style={font=\scriptsize},
  label style={font=\small},
  legend pos=north east,
  legend style={font=\scriptsize},
  clip=false,
]
% Q = 5
\addplot[engblue, very thick, domain=0.1:10, samples=300]
  {(x/5) / sqrt( (1 - x^2)^2 + (x/5)^2 )};
\addlegendentry{$Q = 5$}
% Q = 1.5
\addplot[engred, very thick, domain=0.1:10, samples=300]
  {(x/1.5) / sqrt( (1 - x^2)^2 + (x/1.5)^2 )};
\addlegendentry{$Q = 1.5$}
% -3 dB level (1/sqrt2 = 0.707)
\addplot[enggray, dashed] coordinates {(0.1,0.7071) (10,0.7071)};
\node[enggray, font=\scriptsize, anchor=south west] at (axis cs:0.1,0.71) {$1/\sqrt{2}$};
\end{semilogxaxis}
\end{tikzpicture}
\caption[Bandpass resonance at two quality factors]{Magnitude response of a
series RLC bandpass filter at two quality factors. The peak occurs at the
resonant frequency $f_0 = 1/(2\pi\sqrt{LC})$, where the reactances of the
inductor and capacitor cancel and the impedance is purely resistive. The
higher-$Q$ curve is narrower; the $-3$ dB bandwidth is $\Delta f = f_0/Q$.
The bandpass project at the end of this chapter targets $Q = 5$, the blue
curve.}
\label{fig:vol04ch09:bandpass-resonance}
\end{figure}


A filter with more than one pole stacks the slopes. A two-pole low-pass
response is flat in the pass band, falls at $-20\,\si{\decibel}$ per
decade past the first pole, and steepens to $-40\,\si{\decibel}$ per
decade past the second; its phase swings $-45^\circ$ per decade through
each pole toward a final $-180^\circ$.
Figure~\ref{fig:vol04ch09:bode-twopole} draws the magnitude and phase of
a two-pole filter with its straight-line asymptotes, the construction
the working engineer sketches before refining the $-3\,\si{\decibel}$
rounding at each corner on the bench.

% Figure: asymptotic Bode magnitude and phase of a two-pole low-pass transfer
% function, with the straight-line asymptotes overlaid on the exact curves.
% H(s) = 1 / ((1 + s/w1)(1 + s/w2)), w1 = 1, w2 = 10 (normalised). pgfplots.
\begin{figure}[htbp]
\centering
% === magnitude ===
\begin{tikzpicture}
\begin{semilogxaxis}[
  width=11.5cm, height=4.6cm,
  xlabel={frequency $\omega/\omega_1$},
  ylabel={$|H|$ (dB)},
  xmin=0.01, xmax=1000, ymin=-90, ymax=10,
  axis lines=left,
  tick label style={font=\tiny},
  label style={font=\scriptsize},
  ytick={0,-20,-40,-60,-80},
  clip=false,
]
% exact magnitude: 20 log10 [ 1/sqrt((1+x^2)(1+(x/10)^2)) ]
\addplot[engblue, very thick, domain=0.01:1000, samples=300]
  {-10*log10((1 + x^2)*(1 + (x/10)^2))};
% asymptotes: 0 dB to w1, -20 dB/dec from w1 to w2, -40 dB/dec beyond w2
\addplot[engred, dashed] coordinates {(0.01,0) (1,0)};
\addplot[engred, dashed] coordinates {(1,0) (10,-20)};
\addplot[engred, dashed, domain=10:1000, samples=2] {-20 - 40*log10(x/10)};
\node[engorange, font=\tiny, anchor=south west] at (axis cs:1,-3) {$\omega_1$};
\node[engorange, font=\tiny, anchor=south west] at (axis cs:10,-23) {$\omega_2$};
\end{semilogxaxis}
\end{tikzpicture}

\vspace{0.3cm}
% === phase ===
\begin{tikzpicture}
\begin{semilogxaxis}[
  width=11.5cm, height=4.0cm,
  xlabel={frequency $\omega/\omega_1$},
  ylabel={$\angle H$ (deg)},
  xmin=0.01, xmax=1000, ymin=-190, ymax=10,
  axis lines=left,
  tick label style={font=\tiny},
  label style={font=\scriptsize},
  ytick={0,-45,-90,-135,-180},
  clip=false,
]
% exact phase: -(atan(x) + atan(x/10)), in degrees
\addplot[engblue, very thick, domain=0.01:1000, samples=300]
  {-(atan(x) + atan(x/10))};
% straight-line phase asymptotes: -45 deg/dec around each pole
\addplot[engred, dashed] coordinates {(0.01,0) (0.1,0)};
\addplot[engred, dashed] coordinates {(0.1,0) (10,-90)};
\addplot[engred, dashed] coordinates {(1,-45) (100,-135)};
\addplot[engred, dashed, domain=100:1000, samples=2] {-180};
\end{semilogxaxis}
\end{tikzpicture}
\caption[Asymptotic Bode plot of a two-pole filter]{Magnitude (top)
and phase (bottom) of a two-pole low-pass transfer function with poles
at $\omega_1$ and $\omega_2 = 10\,\omega_1$, exact curves in blue and
straight-line asymptotes dashed. The magnitude is flat to $\omega_1$,
falls at $-20\,\si{\decibel}$ per decade between the poles, then at
$-40\,\si{\decibel}$ per decade beyond $\omega_2$; the phase swings
$-45^\circ$ per decade through each pole toward a final $-180^\circ$.
Reading poles off the asymptote breakpoints, then refining the
$-3\,\si{\decibel}$ rounding at each corner, lets the engineer sketch
the response of any pole-zero filter before measuring it.}
\label{fig:vol04ch09:bode-twopole}
\end{figure}


\begin{mastery}
Above the lumped band the two-port view replaces the loop-and-node view.
A linear two-port is described by a $2\times 2$ matrix relating its port
voltages and currents: the impedance ($Z$), admittance ($Y$), hybrid
($h$), or transmission ($ABCD$) parameters, each convenient for a
different connection (series, parallel, cascade). At radio frequencies,
where open and short terminations are hard to realise without parasitic
reactance, engineers switch to scattering ($S$) parameters, which
describe the two-port by the ratios of reflected to incident travelling
waves at a reference impedance (usually $50\,\si{\ohm}$). The $S$-matrix
is what a vector network analyser measures directly, and it connects
this chapter's impedance language to the transmission-line analysis of
Chapter 10: $S_{11}$ is the input reflection coefficient, and a
well-matched port has $|S_{11}|$ small.
\end{mastery}

\section{The Laplace transform of circuit transients}\pathtag{standard}
\label{sec:vol04ch09:laplace}

The characteristic-polynomial method of
Section~\ref{sec:vol04ch09:transient} solves a transient by guessing the
form of the homogeneous solution and matching initial conditions. The
Laplace transform does the same work with the initial conditions built in
from the start, and it turns the differential equation into an algebraic
one whose solution carries both the transient and the steady state in a
single ratio of polynomials. The transform of a time function $f(t)$ is
\[
F(s) = \mathcal{L}\{f(t)\} = \int_{0}^{\infty} f(t)\,e^{-st}\,\dd t,
\quad s = \sigma + j\omega,
\]
and its two working properties are the transform of a derivative,
$\mathcal{L}\{\dd f/\dd t\} = sF(s) - f(0^{-})$, which folds the initial
value into the algebra, and the transform of an integral, which divides by
$s$. Each circuit element acquires an $s$-domain impedance: a resistor
keeps $Z_R = R$; an inductor with initial current $i_L(0)$ becomes an
impedance $sL$ in series with a voltage source $Li_L(0)$ that injects the
stored flux; a capacitor with initial voltage $v_C(0)$ becomes an
impedance $1/(sC)$ in series with a voltage source $v_C(0)/s$ that injects
the stored charge. With every element replaced by its $s$-domain model,
the circuit obeys KVL and KCL in $s$ exactly as the phasor circuit obeyed
them in $j\omega$, and the unknown node voltages and branch currents are
the solution of a linear algebraic system.

A series RLC circuit driven from rest by a step $V_s/s$ is the cleanest
example. The voltage across the capacitor is the source times the divider
formed by the capacitor impedance over the total series impedance,
\[
V_C(s) = \frac{V_s}{s}\cdot
\frac{1/(sC)}{R + sL + 1/(sC)}
= \frac{V_s}{s}\cdot
\frac{\omega_n^{2}}{s^{2} + 2\zeta\omega_n s + \omega_n^{2}},
\]
with $\omega_n = 1/\sqrt{LC}$ and $\zeta = (R/2)\sqrt{C/L}$ as before. The
denominator factors into the pole at the origin, contributed by the step,
and the quadratic whose roots are the natural frequencies of the circuit.
Partial-fraction expansion separates the response into a term per pole,
\[
V_C(s) = \frac{V_s}{s}
- \frac{V_s\,(s + 2\zeta\omega_n)}{s^{2} + 2\zeta\omega_n s + \omega_n^{2}},
\]
and the inverse transform reads each term back into the time domain: the
$V_s/s$ term is the constant final value $V_s$, and the quadratic term,
completed to the square $(s + \zeta\omega_n)^{2} + \omega_d^{2}$ with
$\omega_d = \omega_n\sqrt{1 - \zeta^{2}}$, inverts to the decaying
sinusoid
\[
v_C(t) = V_s\left[1 - \frac{e^{-\zeta\omega_n t}}{\sqrt{1 - \zeta^{2}}}
\sin\!\left(\omega_d t + \arccos\zeta\right)\right],
\]
the same underdamped ringing the characteristic method produced, now
arriving without a separate initial-condition step. The location of the
poles, $-\zeta\omega_n \pm j\omega_d$, fixes both the envelope decay rate
(the real part) and the ringing frequency (the imaginary part), so the
$s$-plane map of the poles is the response read at a glance.

The poles and zeros of a transfer function carry the same information as
the time response and the Bode plot, and reading among the three
representations is the working fluency the rest of the chapter assumes. A
zero of $H(s)$ is a frequency, real or complex, at which the network
transmits nothing; a pole is a natural frequency at which the unforced
network rings or decays. A simple pole on the negative real axis at
$s = -1/\tau$ is a first-order decay of time constant $\tau$; a
complex-conjugate pole pair is second-order ringing; a zero at the origin
is the differentiating action of a high-pass stage, which kills the DC
term. The residue of each pole, the coefficient the partial-fraction
expansion attaches to it, sets how strongly that natural mode appears in
the response to a given input, so a pole close to a zero contributes a
small residue and a nearly cancelled mode, the basis of pole-zero
cancellation in compensator design.

\begin{principle}[Poles are natural frequencies, zeros are
transmission nulls]
The denominator roots of a circuit's transfer function are the
frequencies at which the unforced circuit can sustain a response, the
natural modes; the numerator roots are the frequencies at which the
driven circuit transmits nothing. Stability requires every pole in the
open left half-plane; a pole on the imaginary axis is a sustained
oscillation and a pole in the right half-plane is runaway growth. The
forced steady-state response at frequency $\omega$ is $H(j\omega)$, the
transfer function evaluated on the imaginary axis, which is why the Bode
plot, the pole-zero map, and the impulse response are three views of one
object.
\end{principle}

\begin{mastery}
The transfer function $H(s)$ is the Laplace transform of the circuit's
\engterm{impulse response} $h(t)$, the output when the input is a unit
impulse, and that pairing turns the response to any input into a single
integral. A linear time-invariant circuit obeys superposition in time as
well as across sources: an arbitrary input $x(t)$ is a continuum of
scaled, shifted impulses, and the output is the same continuum of scaled,
shifted impulse responses summed, the \engterm{convolution}
\[
y(t) = \int_0^{t} h(t - \xi)\,x(\xi)\,\dd\xi.
\]
Convolution in time is multiplication in the $s$-domain, $Y(s) = H(s)X(s)$,
which is the entire reason the transform is worth the trouble: the messy
overlap integral becomes a product of polynomials, solved by the
partial-fraction recipe and inverted term by term. The step response is
the convolution of the impulse response with a unit step, equivalently the
running integral of the impulse response, so the three standard test
responses, impulse, step, and sinusoid, are one function viewed under
three inputs. The same convolution is the discrete-time filter of digital
signal processing, where the integral becomes a sum over taps and $H(s)$
becomes the $z$-transform $H(z)$; the lumped circuit and the digital
filter are the continuous and sampled faces of one linear-systems
algebra, the bridge this chapter builds toward the signal-processing and
control volumes.
\end{mastery}

\section{Fourier series of periodic waveforms and harmonic power}
\pathtag{standard}
\label{sec:vol04ch09:fourier}

Phasor analysis handles one frequency at a time; a periodic waveform that
is not a pure sinusoid carries many frequencies at once, and the Fourier
series is the tool that resolves it into them. Any periodic voltage
$v(t)$ of period $T$ and fundamental angular frequency
$\omega_0 = 2\pi/T$ expands as a sum of sinusoids at integer multiples of
the fundamental,
\[
v(t) = a_0 + \sum_{n=1}^{\infty}
\left[a_n\cos(n\omega_0 t) + b_n\sin(n\omega_0 t)\right],
\]
with coefficients recovered by projecting the waveform onto each harmonic,
\[
a_n = \frac{2}{T}\int_0^T v(t)\cos(n\omega_0 t)\,\dd t,
\qquad
b_n = \frac{2}{T}\int_0^T v(t)\sin(n\omega_0 t)\,\dd t,
\]
and $a_0$ the average value. The orthogonality of the sinusoids over a
period, $\int_0^T \cos(m\omega_0 t)\cos(n\omega_0 t)\,\dd t = 0$ for
$m \neq n$, is what makes the projection clean: each integral picks out
exactly one coefficient and discards the rest.

A symmetric square wave of amplitude $V$ swinging about zero has only odd
harmonics, with amplitudes that fall as $1/n$,
\[
v(t) = \frac{4V}{\pi}\sum_{k=0}^{\infty}
\frac{\sin\!\big((2k+1)\omega_0 t\big)}{2k+1}
= \frac{4V}{\pi}\left[\sin\omega_0 t + \tfrac{1}{3}\sin 3\omega_0 t
+ \tfrac{1}{5}\sin 5\omega_0 t + \cdots\right].
\]
Figure~\ref{fig:vol04ch09:fourier-square} plots the partial sums against
the target square wave and the harmonic-amplitude spectrum. The partial
sum keeps a fixed overshoot of about nine percent at the discontinuity,
the Gibbs phenomenon, which narrows as the harmonic count rises but never
falls below that height, so a finite-bandwidth channel rounds a square
edge and rings at it rather than reproducing the corner. The triangular
wave, by contrast, has odd harmonics falling as $1/n^{2}$, so it
reconstructs far faster and rings less, the harmonic signature of its
continuous slope.

% Figure: the partial-sum reconstruction of a square wave from its odd
% harmonics, showing the Gibbs overshoot at the discontinuity, and the
% harmonic-amplitude spectrum. pgfplots. The square wave has period 2 and
% unit amplitude; harmonics fall as 1/n for odd n only.
\begin{figure}[htbp]
\centering
% === left: partial sums in the time domain ===
\begin{tikzpicture}
\begin{axis}[
  width=7.2cm, height=5.4cm,
  xlabel={$t/T$}, ylabel={$v(t)$},
  xmin=0, xmax=1, ymin=-1.5, ymax=1.5,
  axis lines=left,
  tick label style={font=\tiny},
  label style={font=\scriptsize},
  legend style={font=\tiny, at={(0.98,0.02)}, anchor=south east, draw=none, fill=none},
  clip=false,
]
% one term: (4/pi) sin(2 pi t)
\addplot[enggray, thick, domain=0:1, samples=200]
  {(4/pi)*sin(deg(2*pi*x))};
\addlegendentry{$N=1$}
% three terms (n = 1,3,5): add (4/pi)[ sin/1 + sin3/3 + sin5/5 ]
\addplot[engorange, thick, domain=0:1, samples=300]
  {(4/pi)*( sin(deg(2*pi*x)) + sin(deg(6*pi*x))/3 + sin(deg(10*pi*x))/5 )};
\addlegendentry{$N=5$}
% many terms (n up to 19): full Gibbs ear visible
\addplot[engblue, very thick, domain=0:1, samples=600]
  {(4/pi)*( sin(deg(2*pi*x)) + sin(deg(6*pi*x))/3 + sin(deg(10*pi*x))/5
   + sin(deg(14*pi*x))/7 + sin(deg(18*pi*x))/9 + sin(deg(22*pi*x))/11
   + sin(deg(26*pi*x))/13 + sin(deg(30*pi*x))/15 + sin(deg(34*pi*x))/17
   + sin(deg(38*pi*x))/19 )};
\addlegendentry{$N=19$}
\draw[engred, dashed] (axis cs:0,1) -- (axis cs:0.5,1);
\draw[engred, dashed] (axis cs:0.5,-1) -- (axis cs:1,-1);
\end{axis}
\end{tikzpicture}
\hspace{0.4cm}
% === right: harmonic amplitude spectrum ===
\begin{tikzpicture}
\begin{axis}[
  width=6.2cm, height=5.4cm,
  xlabel={harmonic $n$}, ylabel={amplitude $|c_n|$},
  xmin=0, xmax=12, ymin=0, ymax=1.4,
  axis lines=left,
  tick label style={font=\tiny},
  label style={font=\scriptsize},
  ybar, bar width=3pt,
  xtick={1,3,5,7,9,11},
  clip=false,
]
% (4/pi)/n for odd n; even harmonics are zero
\addplot[engblue, fill=engblue] coordinates {
  (1,1.2732) (3,0.4244) (5,0.2546) (7,0.1819) (9,0.1415) (11,0.1157)};
\end{axis}
\end{tikzpicture}
\caption[Fourier reconstruction and spectrum of a square wave]{A unit
square wave reconstructed from its Fourier series (left) and its
harmonic-amplitude spectrum (right). The square wave contains only odd
harmonics, whose amplitudes fall as $4/(\pi n)$, drawn as the bar
spectrum. Adding more terms sharpens the edges, but the partial sum
keeps a fixed overshoot of about nine percent at the discontinuity, the
Gibbs phenomenon, which narrows but never shrinks as the harmonic count
rises. The $1/n$ amplitude roll-off is why the power in the higher
harmonics falls as $1/n^2$ and the total harmonic content of a
fast-edged waveform concentrates in its first few harmonics.}
\label{fig:vol04ch09:fourier-square}
\end{figure}


The power a periodic waveform carries is the sum of the powers of its
harmonics, a result that follows from the same orthogonality and is named
for Parseval. For a periodic voltage across a resistor $R$ the average
power is
\[
P = \frac{1}{R}\,\overline{v^{2}}
= \frac{1}{R}\left[a_0^{2} + \tfrac{1}{2}\sum_{n=1}^{\infty}
(a_n^{2} + b_n^{2})\right]
= \frac{V_0^{2}}{R} + \sum_{n=1}^{\infty}\frac{V_{n,\text{rms}}^{2}}{R},
\]
each harmonic contributing independently because the cross terms
integrate to zero. The decomposition gives the engineer two daily
quantities. The first is the \engterm{RMS value} of a non-sinusoidal
waveform, $V_{\text{rms}} = \sqrt{V_0^{2} + \sum_n V_{n,\text{rms}}^{2}}$,
the root of the sum of the squared harmonic RMS values, which a
true-RMS meter computes and an averaging meter calibrated for sinusoids
does not. The second is \engterm{total harmonic distortion}, the ratio of
the RMS of the harmonics above the fundamental to the RMS of the
fundamental,
\[
\text{THD} = \frac{\sqrt{\sum_{n\geq 2} V_{n,\text{rms}}^{2}}}
{V_{1,\text{rms}}},
\]
the figure of merit for an amplifier's linearity, a power supply's
current waveform, and a utility's voltage quality. A load that draws a
distorted current loads the supply with harmonic power that does no useful
work and heats the neutral conductor of a three-phase system, which is why
power-quality standards cap the harmonic current a device may inject.

\section{Two-port network parameters}\pathtag{mastery}
\label{sec:vol04ch09:two-port}

A circuit block accessed through two terminal pairs, an input port and an
output port, is a \engterm{two-port}, and the working description of such
a block is a $2\times 2$ matrix relating its four port variables: the
input voltage $V_1$ and current $I_1$, and the output voltage $V_2$ and
current $I_2$. Which two of the four are taken as the independent
variables fixes which parameter set results, and each set is the
convenient one for a particular way of connecting blocks together.
Figure~\ref{fig:vol04ch09:two-port} fixes the sign convention and shows
the cascade that motivates the transmission parameters.

% Figure: a generic linear two-port with its port voltages and currents
% labelled to the standard convention, and the cascade of two two-ports
% that motivates the ABCD parameters. Plain TikZ, no circuitikz.
\begin{figure}[htbp]
\centering
\begin{tikzpicture}[
  line width=0.8pt,
  font=\small,
]
  % === single two-port box ===
  \draw[engblue] (0,0) rectangle (2.4,2.0);
  \node[font=\small] at (1.2,1.0) {two-port};
  % port 1 terminals (left)
  \draw[engblue] (-1.2,1.7) -- (0,1.7);
  \draw[engblue] (-1.2,0.3) -- (0,0.3);
  \node[circle, fill=engblack, inner sep=0.7pt] at (-1.2,1.7) {};
  \node[circle, fill=engblack, inner sep=0.7pt] at (-1.2,0.3) {};
  % port 1 voltage and current
  \draw[engred, -{Latex[length=2mm]}] (-1.05,1.55) -- (-0.65,1.55);
  \node[engred, font=\tiny, anchor=south] at (-0.85,1.55) {$I_1$};
  \draw[enggray, -{Latex[length=1.6mm]}] (-1.2,0.45) -- (-1.2,1.55);
  \node[enggray, font=\tiny, anchor=east] at (-1.25,1.0) {$V_1$};
  \node[font=\tiny, anchor=east] at (-1.25,1.7) {$+$};
  \node[font=\tiny, anchor=east] at (-1.25,0.3) {$-$};
  % port 2 terminals (right)
  \draw[engblue] (2.4,1.7) -- (3.6,1.7);
  \draw[engblue] (2.4,0.3) -- (3.6,0.3);
  \node[circle, fill=engblack, inner sep=0.7pt] at (3.6,1.7) {};
  \node[circle, fill=engblack, inner sep=0.7pt] at (3.6,0.3) {};
  % port 2 voltage and current (into the port)
  \draw[engred, -{Latex[length=2mm]}] (3.05,1.55) -- (2.65,1.55);
  \node[engred, font=\tiny, anchor=south] at (2.85,1.55) {$I_2$};
  \draw[enggray, -{Latex[length=1.6mm]}] (3.6,0.45) -- (3.6,1.55);
  \node[enggray, font=\tiny, anchor=west] at (3.65,1.0) {$V_2$};
  \node[font=\tiny, anchor=west] at (3.65,1.7) {$+$};
  \node[font=\tiny, anchor=west] at (3.65,0.3) {$-$};

  % === cascade of two two-ports ===
  \begin{scope}[xshift=6.0cm]
  \draw[engblue] (0,0) rectangle (1.6,2.0);
  \node[font=\scriptsize] at (0.8,1.0) {$A_1 B_1$};
  \node[font=\scriptsize] at (0.8,0.55) {$C_1 D_1$};
  \draw[engblue] (1.6,0) rectangle (3.2,2.0);
  \node[font=\scriptsize] at (2.4,1.0) {$A_2 B_2$};
  \node[font=\scriptsize] at (2.4,0.55) {$C_2 D_2$};
  % external leads
  \draw[engblue] (-0.9,1.7) -- (0,1.7);
  \draw[engblue] (-0.9,0.3) -- (0,0.3);
  \draw[engblue] (3.2,1.7) -- (4.1,1.7);
  \draw[engblue] (3.2,0.3) -- (4.1,0.3);
  % internal connection nodes
  \draw[engblue] (1.6,1.7) -- (1.6,1.7);
  \node[circle, fill=engblack, inner sep=0.7pt] at (1.6,1.7) {};
  \node[circle, fill=engblack, inner sep=0.7pt] at (1.6,0.3) {};
  \node[font=\tiny, anchor=south] at (-0.45,1.75) {in};
  \node[font=\tiny, anchor=south] at (3.65,1.75) {out};
  \node[font=\small] at (1.6,-0.7) {cascade: $T = T_1 T_2$};
  \end{scope}
\end{tikzpicture}
\caption[A linear two-port and a cascade]{Left: a linear two-port with
its port voltages and currents in the standard convention, both port
currents drawn flowing into the upper terminal. A choice of which two of
the four variables are taken as independent fixes which parameter set
($Z$, $Y$, $h$, or $ABCD$) describes the box. Right: two two-ports in
cascade, the output of the first driving the input of the second. The
$ABCD$ (transmission) description is the one that makes a cascade a
matrix product, $T = T_1 T_2$, which is why a chain of stages, filter
sections, or transmission-line segments is analysed in $ABCD$ form.}
\label{fig:vol04ch09:two-port}
\end{figure}


The \engterm{impedance} or $Z$ parameters take the two currents as
independent and give the two voltages,
\[
\begin{bmatrix} V_1 \\ V_2 \end{bmatrix}
=
\begin{bmatrix} z_{11} & z_{12} \\ z_{21} & z_{22} \end{bmatrix}
\begin{bmatrix} I_1 \\ I_2 \end{bmatrix},
\]
where each parameter is measured with one port open-circuited:
$z_{11} = V_1/I_1$ with the output open ($I_2 = 0$) is the input
impedance with the output unloaded, and $z_{12} = V_1/I_2$ with the input
open is the reverse transfer impedance. Open-circuit measurements make the
$Z$ parameters natural for a series-series interconnection, where the
impedance matrices of the connected blocks add. The \engterm{admittance}
or $Y$ parameters are the dual, taking the voltages as independent and
measuring with ports short-circuited; the $Y$ matrices of two blocks in
parallel-parallel connection add.

The \engterm{hybrid} or $h$ parameters mix the two, taking $I_1$ and $V_2$
as independent,
\[
\begin{bmatrix} V_1 \\ I_2 \end{bmatrix}
=
\begin{bmatrix} h_{11} & h_{12} \\ h_{21} & h_{22} \end{bmatrix}
\begin{bmatrix} I_1 \\ V_2 \end{bmatrix},
\]
which suits a transistor: $h_{11}$ is the input impedance with the output
shorted, $h_{21}$ is the forward current gain with the output shorted (the
$\beta$ of a bipolar transistor in the small-signal model), $h_{12}$ is
the reverse voltage feedback, and $h_{22}$ is the output admittance with
the input open. The hybrid description is why a transistor data sheet
quotes $h_{fe}$ for current gain and $h_{ie}$ for input impedance.

The \engterm{transmission} or $ABCD$ parameters express the input
quantities in terms of the output,
\[
\begin{bmatrix} V_1 \\ I_1 \end{bmatrix}
=
\begin{bmatrix} A & B \\ C & D \end{bmatrix}
\begin{bmatrix} V_2 \\ -I_2 \end{bmatrix},
\]
with the output current sign flipped so that $-I_2$ is the current
flowing on into the next stage. The transmission form has the property
the others lack: when two two-ports are cascaded, the output of the first
feeding the input of the second, the overall $ABCD$ matrix is the matrix
product of the two, $T = T_1 T_2$. A chain of stages, a ladder filter, or
a string of transmission-line segments is therefore analysed by
multiplying the transmission matrices in order, the same bookkeeping that
makes the transfer-matrix method the workhorse of filter synthesis and of
the transmission-line cascades of Chapter 10. A simple series impedance
$Z$ has $T = \left[\begin{smallmatrix} 1 & Z \\ 0 & 1
\end{smallmatrix}\right]$ and a simple shunt admittance $Y$ has $T =
\left[\begin{smallmatrix} 1 & 0 \\ Y & 1 \end{smallmatrix}\right]$, so an
L-section, a T-section, or a pi-section is built up by multiplying these
elementary blocks.

The four parameter sets describe the same box and convert into one
another by algebra; an engineer picks the set that makes the
interconnection at hand a simple addition or product, computes in it, and
converts back if the answer is wanted in another form. A reciprocal,
passive two-port satisfies $z_{12} = z_{21}$ in the $Z$ set and
$AD - BC = 1$ in the transmission set, a constraint that halves the number
of independent parameters and is the first check on a measured parameter
table.

\section{Noise and signal-to-noise ratio of an RC stage}\pathtag{standard}
\label{sec:vol04ch09:noise}

Every resistor at a temperature above absolute zero is a source of random
voltage, the thermal agitation of its charge carriers, and that voltage
sets the noise floor below which no signal in the circuit can be
distinguished. The \engterm{Johnson-Nyquist} noise of a resistor $R$ at
temperature $T$ has a flat (white) voltage spectral density
\[
S_v(f) = 4k_B T R \quad [\si{\volt\squared\per\hertz}],
\]
with $k_B$ the Boltzmann constant, so the mean-square noise voltage
measured over a bandwidth $\Delta f$ is $\overline{v_n^{2}} = 4k_B T
R\,\Delta f$ and the RMS noise voltage is $\sqrt{4k_B T R\,\Delta f}$. The
density does not depend on frequency, which is why the noise a resistor
delivers to a circuit is set by the bandwidth the rest of the circuit
allows it, not by the resistor alone.

The first-order RC low-pass stage makes this concrete and yields a result
of unexpected simplicity. The resistor injects white noise of density
$4k_B T R$, and the capacitor's low-pass response shapes it: the transfer
function from the resistor noise to the capacitor voltage is $H(j\omega) =
1/(1 + j\omega RC)$, of magnitude-squared $|H|^{2} = 1/(1 + (\omega
RC)^{2})$. The total noise voltage at the capacitor is the density times
the squared response integrated over all frequency,
\[
\overline{v_C^{2}} = \int_0^{\infty} 4k_B T R\,
\frac{1}{1 + (2\pi f RC)^{2}}\,\dd f
= 4k_B T R\cdot\frac{1}{4RC}
= \frac{k_B T}{C}.
\]
The integral of the Lorentzian low-pass shape is $1/(4RC)$, the
\engterm{noise bandwidth} of a first-order filter (which is $\pi/2$ times
its $-3\,\si{\decibel}$ bandwidth $1/(2\pi RC)$), and the resistance
cancels: the noise voltage on the capacitor is $\sqrt{k_B T/C}$,
independent of $R$. This is the \engterm{kTC noise} that floors every
sampled-data system, because the act of sampling a voltage onto a
capacitor freezes exactly this much noise onto it regardless of the switch
resistance. The resistance sets how fast the noise settles, not how much
of it there is; a larger sampling capacitor is the only way to lower the
floor, at the cost of a slower acquisition.

\begin{principle}[The capacitor sets the sampled-noise floor]
The thermal noise sampled onto a capacitor through any resistance is
$\sqrt{k_B T/C}$, set by the capacitance and the temperature alone. A
designer who needs a quieter sample-and-hold or a quieter pixel does not
lower the switch resistance, which changes only the settling time; the
designer raises the sampling capacitance, trading acquisition speed and
silicon area for a lower noise floor. At room temperature a
$1\,\si{\pico\farad}$ capacitor floors at about $64\,\si{\micro\volt}$
RMS, and a $1\,\si{\nano\farad}$ capacitor at about $2\,\si{\micro\volt}$
RMS, the square-root scaling that governs every charge-domain sensor.
\end{principle}

The \engterm{signal-to-noise ratio} of the stage is the ratio of the
signal power to this noise power, $\text{SNR} = V_{\text{sig,rms}}^{2}/
\overline{v_C^{2}}$, usually quoted in decibels as $10\log_{10}$ of the
power ratio or $20\log_{10}$ of the voltage ratio. The SNR sets the number
of bits a downstream converter can use: a converter resolving below the
noise floor wastes its lower bits on digitised noise, so the working rule
pairs the sampling capacitor, the signal swing, and the converter
resolution into one budget, the same budget the sensor-chain case study
draws for its strain front end. When several noise sources act together
and are uncorrelated, their mean-square voltages add, so the total RMS
noise is the root-sum-square of the contributions, and the design lever is
always the largest single contributor.

\section{Worked examples}\pathtag{standard}

The worked examples below carry the machinery from formula to number.
Each is the kind of calculation that recurs daily at a bench or in a
SPICE deck, and together they touch every section of the chapter at
least once.

\paragraph{Node voltage by inspection.}
A $12\,\si{\volt}$ source drives a $100\,\si{\ohm}$ resistor into a
node, and a $50\,\si{\ohm}$ resistor returns that node to ground. The
single non-reference node voltage $v_{1}$ follows from one KCL
equation,
\[
\frac{v_{1}-12}{100} + \frac{v_{1}}{50} = 0
\quad\Longrightarrow\quad
v_{1}\left(\frac{1}{100}+\frac{1}{50}\right) = \frac{12}{100},
\]
so $v_{1} = 12\cdot(1/100)/(3/100) = 4\,\si{\volt}$, the voltage
divider result. The companion script
\texttt{code/nodal\_solver.py} stamps the
same circuit into a modified-nodal-analysis matrix and returns
$v_{1}=4.0\,\si{\volt}$, the procedure every simulator uses on
networks too large to solve by hand.

\paragraph{Th\'evenin reduction of a bridge leg.}
Take a $10\,\si{\volt}$ source in series with $1\,\si{\kilo\ohm}$
feeding the parallel pair $2\,\si{\kilo\ohm}\parallel 4\,\si{\kilo
\ohm}$, with terminal $a$ at the far end of the $4\,\si{\kilo\ohm}$
resistor. The open-circuit voltage is the voltage across the parallel
pair, $V_{\text{Th}} = 10\cdot R_{p}/(1\,\text{k}+R_{p})$ with $R_{p} =
(2)(4)/(2+4) = 1.33\,\si{\kilo\ohm}$, giving $V_{\text{Th}} =
10\cdot 1.33/2.33 = 5.71\,\si{\volt}$. The Th\'evenin resistance is the
source resistance deactivated and seen from $a$-$b$: $R_{\text{Th}} =
1\,\text{k}\parallel(2\,\text{k}\parallel 4\,\text{k})$ wait, the
$4\,\si{\kilo\ohm}$ resistor is in series to terminal $a$, so
$R_{\text{Th}} = 4\,\text{k} + (2\,\text{k}\parallel 1\,\text{k}) =
4\,\text{k}+0.67\,\text{k} = 4.67\,\si{\kilo\ohm}$. The reduction
replaces a four-element network with one source and one resistor for
every downstream calculation.

\paragraph{Bandpass numbers from a sweep.}
The project filter ($L=10\,\si{\milli\henry}$, $C=24\,\si{\nano\farad}$,
$R=126\,\si{\ohm}$) has
$f_{0}=1/(2\pi\sqrt{LC})=10.3\,\si{\kilo\hertz}$ and
$Q=\omega_{0}L/R\approx 5$, so the $-3\,\si{\decibel}$ bandwidth is
$f_{0}/Q\approx 2\,\si{\kilo\hertz}$. The script
\texttt{code/ac\_sweep.py} computes the
complex transfer function over a logarithmic sweep and confirms these
figures; a synthetic measured sweep with bench-style scatter is in
\texttt{data/bandpass\_sweep.csv} for the
reader who wants to practise fitting the resonance from noisy data. The
transient script
\texttt{code/rlc\_transient.py} integrates
the same second-order circuit in the time domain and agrees with the
closed-form step response to under a millivolt at
$\omega_{n}\Delta t = 0.01$, while
\texttt{data/rc\_step\_measured.csv} tabulates
a first-order charging transient against theory.

\paragraph{Second-order RLC step by characteristic roots.}
A series RLC circuit with $L = 1\,\si{\milli\henry}$, $C = 1\,\si{\micro
\farad}$, and $R = 12\,\si{\ohm}$ is driven by a $5\,\si{\volt}$ step
from rest. The natural frequency is $\omega_{n} = 1/\sqrt{LC} =
1/\sqrt{10^{-3}\cdot 10^{-6}} = 3.16\times 10^{4}\,\si{\per\second}$, and
the damping ratio is $\zeta = (R/2)\sqrt{C/L} = 6\cdot\sqrt{10^{-6}/
10^{-3}} = 6\cdot 0.0316 = 0.19$, so the response is underdamped. The
characteristic roots are $s = -\zeta\omega_{n}\pm j\omega_{n}\sqrt{1-
\zeta^{2}} = -6.0\times 10^{3}\pm j\,3.10\times 10^{4}\,\si{\per
\second}$. The damped frequency is $\omega_{d} = 3.10\times 10^{4}\,\si{
\per\second}$ ($f_{d} = 4.94\,\si{\kilo\hertz}$), the envelope time
constant is $1/(\zeta\omega_{n}) = 167\,\si{\micro\second}$, and the
peak overshoot for $\zeta = 0.19$ is $e^{-\pi\zeta/\sqrt{1-\zeta^{2}}} =
0.54$, so the capacitor voltage rings up to $5(1+0.54) = 7.7\,\si{\volt}$
before settling to $5\,\si{\volt}$. The capacitor sees half again its
supply voltage at the first peak, a transient that a part rated only to
the nominal rail would not survive. Figure~\ref{fig:vol04ch09:rlc-roots}
places the roots of this circuit and its two heavier-damped siblings on
the $s$-plane and shows the step response each produces.

% Figure: the three damping regimes of a second-order RLC step response
% located on the s-plane by the characteristic roots, with the matching
% time-domain step responses. Left: root loci; right: step curves. pgfplots.
\begin{figure}[htbp]
\centering
% === left: s-plane root positions ===
\begin{tikzpicture}
\begin{axis}[
  width=6.2cm, height=6.2cm,
  xlabel={$\Re(s)$ ($\times 10^{4}\,\si{\per\second}$)},
  ylabel={$\Im(s)$ ($\times 10^{4}\,\si{\per\second}$)},
  xmin=-12, xmax=2, ymin=-7, ymax=7,
  axis lines=middle,
  tick label style={font=\tiny},
  label style={font=\tiny},
  clip=false,
]
% underdamped pair (complex)
\addplot[only marks, mark=x, engblue, mark size=3pt]
  coordinates {(-1,5.916) (-1,-5.916)};
\node[engblue, font=\tiny, anchor=south west] at (axis cs:-1,5.916) {underdamped};
% critically damped (repeated real)
\addplot[only marks, mark=square*, engorange, mark size=2pt]
  coordinates {(-6,0)};
\node[engorange, font=\tiny, anchor=south] at (axis cs:-6,0.5) {critical};
% overdamped (two distinct real)
\addplot[only marks, mark=*, engred, mark size=2pt]
  coordinates {(-1.0,0) (-11.0,0)};
\node[engred, font=\tiny, anchor=north] at (axis cs:-8,-0.6) {overdamped};
\end{axis}
\end{tikzpicture}
\hspace{0.6cm}
% === right: step responses ===
\begin{tikzpicture}
\begin{axis}[
  width=6.2cm, height=6.2cm,
  xlabel={$\omega_n t$}, ylabel={$v_C/V_s$},
  xmin=0, xmax=12, ymin=0, ymax=1.6,
  axis lines=left,
  tick label style={font=\tiny},
  label style={font=\tiny},
  legend style={font=\tiny, at={(0.98,0.03)}, anchor=south east},
  clip=false,
]
% underdamped zeta=0.17: omega_d = sqrt(1-zeta^2)=0.9854, phase=acos(zeta)=80.21 deg
\addplot[engblue, very thick, domain=0:12, samples=200]
  {1 - exp(-0.17*x)/0.9854*sin(deg(0.9854*x) + 80.21)};
\addlegendentry{$\zeta=0.17$}
% critically damped zeta=1
\addplot[engorange, thick, domain=0:12, samples=200]
  {1 - (1 + x)*exp(-x)};
\addlegendentry{$\zeta=1$}
% overdamped zeta=2: roots -0.268, -3.732 (units of omega_n)
\addplot[engred, thick, domain=0:12, samples=200]
  {1 - (3.732*exp(-0.268*x) - 0.268*exp(-3.732*x))/3.464};
\addlegendentry{$\zeta=2$}
\draw[enggray, dashed] (axis cs:0,1) -- (axis cs:12,1);
\end{axis}
\end{tikzpicture}
\caption[Second-order RLC roots and step response]{The characteristic
roots of a series RLC circuit (left, on the $s$-plane) and the step
response each root configuration produces (right). A complex-conjugate
pair off the real axis (underdamped) gives a ringing response that
overshoots; a repeated real root (critically damped) gives the fastest
monotone approach; two distinct real roots (overdamped) give a slower
sum of two decaying exponentials. The real part of the dominant root
sets the envelope decay rate; the imaginary part sets the ringing
frequency $\omega_d$.}
\label{fig:vol04ch09:rlc-roots}
\end{figure}


\paragraph{Power-factor correction of an inductive load.}
A $400\,\si{\volt}$, $50\,\si{\hertz}$ single-phase load draws
$30\,\si{\ampere}$ at a power factor of $0.70$ lagging. The apparent
power is $S = VI = 400\cdot 30 = 12.0\,\si{\kilo\volt\ampere}$; the real
power is $P = S\cos\phi = 0.70\cdot 12.0 = 8.4\,\si{\kilo\watt}$; the
reactive power is $Q_{1} = S\sin\phi = 12.0\cdot\sin(45.6^\circ) =
8.57\,\text{kvar}$. To raise the power factor to $0.95$ ($\phi_{2} =
18.2^\circ$) the load reactive power must fall to $Q_{2} = P\tan\phi_{2}
= 8.4\cdot 0.329 = 2.76\,\text{kvar}$, so a parallel capacitor must
supply $Q_{C} = Q_{1}-Q_{2} = 5.81\,\text{kvar}$. A shunt capacitor at
line voltage absorbs $Q_{C} = \omega C V^{2}$, so $C = Q_{C}/(\omega
V^{2}) = 5810/(2\pi\cdot 50\cdot 400^{2}) = 116\,\si{\micro\farad}$. The
real power and the load current it commands are unchanged; the line
current falls from $30\,\si{\ampere}$ to $P/(V\cos\phi_{2}) = 8400/
(400\cdot 0.95) = 22.1\,\si{\ampere}$, a $26\,\%$ reduction in the
current the supply wiring carries, which is the saving the correction
buys.

\paragraph{Instrumentation amplifier gain.}
The three-op-amp instrumentation amplifier sets its differential gain
with a single resistor and rejects the common-mode voltage that a plain
difference amplifier would pass. The two input buffers share a gain
resistor $R_{g}$ between their inverting inputs; each buffer feeds a
feedback resistor $R_{1}$. The virtual short forces the full
differential input $v_{\text{in}}^{+}-v_{\text{in}}^{-}$ across $R_{g}$,
so the current through $R_{g}$ is $(v_{\text{in}}^{+}-v_{\text{in}}^{-})/
R_{g}$, and the same current flows through both $R_{1}$ resistors. The
first-stage differential output is therefore
\[
v_{o1}-v_{o2} = \left(1+\frac{2R_{1}}{R_{g}}\right)
\left(v_{\text{in}}^{+}-v_{\text{in}}^{-}\right),
\]
and the output difference amplifier of gain $R_{3}/R_{2}$ multiplies
this, giving a total gain $G = (1+2R_{1}/R_{g})(R_{3}/R_{2})$. With
$R_{1} = 25\,\si{\kilo\ohm}$, $R_{2} = R_{3} = 10\,\si{\kilo\ohm}$, and a
single $R_{g} = 1\,\si{\kilo\ohm}$, the gain is $G = (1+50)\cdot 1 = 51$;
changing $R_{g}$ to $500\,\si{\ohm}$ raises it to $101$. One resistor
sets the gain while the matched pairs hold the common-mode rejection,
which is why the topology dominates bridge and thermocouple front ends.

\paragraph{Underdamped RLC step from rest, carried to the first peak.}
A series RLC network with $L = 4.7\,\si{\milli\henry}$, $C = 47\,\si{\nano
\farad}$, and $R = 33\,\si{\ohm}$ is driven by a $12\,\si{\volt}$ step with
the capacitor and inductor initially at rest. The natural frequency is
$\omega_{n} = 1/\sqrt{LC} = 1/\sqrt{4.7\times 10^{-3}\cdot 47\times 10^{-9}}
= 6.73\times 10^{4}\,\si{\per\second}$, so $f_{n} = 10.7\,\si{\kilo\hertz}$.
The damping ratio is $\zeta = (R/2)\sqrt{C/L} = 16.5\cdot\sqrt{47\times
10^{-9}/4.7\times 10^{-3}} = 16.5\cdot 3.16\times 10^{-3} = 0.052$, deeply
underdamped, with quality factor $Q = 1/(2\zeta) = 9.6$. The damped frequency
is $\omega_{d} = \omega_{n}\sqrt{1-\zeta^{2}} = 6.72\times 10^{4}\,\si{\per
\second}$, almost indistinguishable from $\omega_{n}$ at this low damping,
and the envelope decays with time constant $1/(\zeta\omega_{n}) = 286\,\si{
\micro\second}$. The fractional overshoot is $\exp(-\pi\zeta/\sqrt{1-\zeta^{2}})
= \exp(-\pi\cdot 0.052/0.999) = 0.85$, so the capacitor voltage rings to
$12(1+0.85) = 22\,\si{\volt}$ at the first peak, almost double the rail. The
first peak arrives at $t_{p} = \pi/\omega_{d} = 46.7\,\si{\micro\second}$, and
the ringing takes about $Q$ cycles, here roughly ten, to fall inside a few
percent of the final $12\,\si{\volt}$. A capacitor rated only to the
$12\,\si{\volt}$ supply would see nearly twice its rated voltage on the first
swing; the lesson is that a lightly damped second-order circuit punishes
parts chosen for the steady state alone, and that even a modest series
resistance ($\zeta$ raised toward $0.5$) tames the overshoot at the cost of a
slower approach.

\paragraph{Transistor bias point by the load line.}
A common-emitter stage uses an npn transistor with current gain $\beta = 150$
and base-emitter drop $V_{BE} = 0.7\,\si{\volt}$, run from a $V_{CC} =
12\,\si{\volt}$ rail with a collector resistor $R_{C} = 2.2\,\si{\kilo\ohm}$
and an emitter resistor $R_{E} = 470\,\si{\ohm}$. The base is held by a
divider of $R_{1} = 47\,\si{\kilo\ohm}$ (to $V_{CC}$) and $R_{2} = 10\,\si{
\kilo\ohm}$ (to ground). The divider open-circuit voltage is $V_{B} =
V_{CC}R_{2}/(R_{1}+R_{2}) = 12\cdot 10/57 = 2.11\,\si{\volt}$ with a
Th\'evenin source resistance $R_{B} = R_{1}\parallel R_{2} = 8.25\,\si{\kilo
\ohm}$. The emitter current follows from the base loop $V_{B} = I_{B}R_{B} +
V_{BE} + I_{E}R_{E}$ with $I_{B} = I_{E}/(\beta+1)$, so
\[
I_{E} = \frac{V_{B}-V_{BE}}{R_{E} + R_{B}/(\beta+1)}
= \frac{2.11-0.7}{470 + 8250/151}
= \frac{1.41}{525}
= 2.69\,\si{\milli\ampere}.
\]
The collector current differs from the emitter current only by the small
base current, $I_{C}\approx 2.67\,\si{\milli\ampere}$, and the collector-emitter voltage is
$V_{CE} = V_{CC} - I_{C}(R_{C}+R_{E}) = 12 - 2.67\times 10^{-3}\cdot 2670 =
4.9\,\si{\volt}$, comfortably in the active region and near the centre of the
load line, where the stage has the most symmetric swing before clipping. The
emitter degeneration is what stabilises the operating point: the $R_{B}/
(\beta+1) = 55\,\si{\ohm}$ term is small against $R_{E} = 470\,\si{\ohm}$, so
$I_{E}$ depends weakly on $\beta$, and a transistor with $\beta = 300$ instead
of $150$ shifts $I_{E}$ by under a percent. Bias stability is bought by
trading voltage gain (the emitter resistor lowers it) for insensitivity to
the part-to-part scatter in $\beta$.

\paragraph{Transformer impedance matching to a loudspeaker.}
An audio output stage with a Th\'evenin source resistance of $R_{s} =
500\,\si{\ohm}$ must deliver maximum power to an $8\,\si{\ohm}$ loudspeaker.
Connected directly, the speaker would receive a fraction $R_{L}/(R_{s}+R_{L})
= 8/508 = 1.6\,\%$ of the available voltage and a small fraction of the
available power. An ideal transformer matches them: the load reflected to the
primary is $Z_{\text{ref}} = Z_{L}/n^{2}$, and setting $Z_{\text{ref}} =
R_{s}$ gives the required turns ratio $n = \sqrt{Z_{L}/R_{s}} = \sqrt{8/500} =
0.126$, that is $N_{1}:N_{2} = 1/0.126 \approx 7.9:1$ primary-to-secondary.
With the match in place the primary sees $500\,\si{\ohm}$, so half the
available power reaches the reflected load and, the transformer being
lossless, all of that crosses to the speaker. The power delivered is
$V_{s}^{2}/(4R_{s})$, the maximum-power-transfer ceiling, instead of the
$V_{s}^{2}R_{L}/(R_{s}+R_{L})^{2}$ the direct connection would manage, a
ratio of $(R_{s}+R_{L})^{2}/(4R_{s}R_{L}) = 508^{2}/(4\cdot 500\cdot 8) =
16.1$, a twelve-decibel improvement. Figure~\ref{fig:vol04ch09:transformer-match}
draws the matching transformer and the reflected impedance. The match holds at
the frequency where the transformer's magnetising inductance is large and its
leakage small; away from that band the reflected impedance drifts, which is
why a real audio transformer is specified over a stated frequency range.

\paragraph{Balanced three-phase load from nameplate data.}
A three-phase induction motor nameplate reads $400\,\si{\volt}$ line-to-line,
$50\,\si{\hertz}$, $15\,\si{\kilo\watt}$ shaft output, efficiency $0.90$,
power factor $0.86$ lagging. The electrical input power is the shaft power
divided by efficiency, $P_{\text{in}} = 15/0.90 = 16.7\,\si{\kilo\watt}$. The
line current follows from $P_{\text{in}} = \sqrt{3}V_{L}I_{L}\cos\phi$, so
\[
I_{L} = \frac{P_{\text{in}}}{\sqrt{3}V_{L}\cos\phi}
= \frac{16\,700}{\sqrt{3}\cdot 400\cdot 0.86}
= \frac{16\,700}{595.7}
= 28.0\,\si{\ampere}.
\]
The apparent power is $S = \sqrt{3}V_{L}I_{L} = 19.4\,\si{\kilo\volt\ampere}$
and the reactive power is $Q = S\sin\phi = 19.4\cdot\sin(30.7^\circ) =
9.9\,\text{kvar}$, the magnetising demand the supply must carry without
delivering shaft work. Sizing the feeder and the protective device from the
nameplate, an engineer would round the $28\,\si{\ampere}$ up for starting
inrush (a direct-on-line induction motor draws six to eight times run current
at start) and select a contactor and overload relay accordingly. The
$\sqrt{3}$ factor and the efficiency-and-power-factor chain are the daily
arithmetic of reading a motor nameplate into a feeder design.

% Figure: ideal two-winding transformer used for impedance matching, with the
% reflected-impedance relation annotated. Plain TikZ paths and nodes.
\begin{figure}[htbp]
\centering
\begin{tikzpicture}[
  line width=0.8pt,
  font=\scriptsize,
  >={Stealth},
]
% --- source side ---
\coordinate (sa) at (0,2.4);
\coordinate (sb) at (0,0);
\draw (sb) -- (0,1.0);
\draw[line width=1.4pt] (-0.25,1.0) -- (0.25,1.0);
\draw[line width=0.8pt] (-0.15,1.2) -- (0.15,1.2);
\draw (0,1.2) -- (sa);
\node[anchor=east] at (-0.3,1.1) {$V_s,\,R_s$};
% series source resistor (zigzag) on top wire
\draw (sa) -- (0.6,2.4);
\draw (0.6,2.4) -- (0.75,2.55) -- (1.05,2.25) -- (1.35,2.55) -- (1.65,2.25)
  -- (1.8,2.4) -- (2.6,2.4);
\node[anchor=south] at (1.2,2.55) {$R_s$};
\draw (sb) -- (2.6,0);

% --- primary winding (vertical bumps) ---
\draw (2.6,2.4)
  .. controls (2.95,2.2) and (2.95,1.85) .. (2.6,1.7)
  .. controls (2.95,1.5) and (2.95,1.15) .. (2.6,1.05)
  .. controls (2.95,0.85) and (2.95,0.5) .. (2.6,0.35)
  .. controls (2.6,0.25) and (2.6,0.1) .. (2.6,0);
\node[anchor=east] at (2.45,1.2) {$N_1$};

% --- core (two vertical bars) ---
\draw[line width=1.2pt] (3.15,0) -- (3.15,2.4);
\draw[line width=1.2pt] (3.35,0) -- (3.35,2.4);

% --- secondary winding ---
\draw (3.9,2.4)
  .. controls (3.55,2.2) and (3.55,1.85) .. (3.9,1.7)
  .. controls (3.55,1.5) and (3.55,1.15) .. (3.9,1.05)
  .. controls (3.55,0.85) and (3.55,0.5) .. (3.9,0.35)
  .. controls (3.9,0.25) and (3.9,0.1) .. (3.9,0);
\node[anchor=west] at (4.05,1.2) {$N_2$};
\node[anchor=south] at (3.25,2.55) {$n=N_2/N_1$};

% --- load resistor on secondary ---
\draw (3.9,2.4) -- (5.2,2.4) -- (5.2,1.7);
\draw (5.2,1.7) -- (5.05,1.55) -- (5.35,1.25) -- (5.05,0.95) -- (5.35,0.65)
  -- (5.2,0.5) -- (5.2,0);
\draw (3.9,0) -- (5.2,0);
\node[anchor=west] at (5.45,1.2) {$Z_L$};

% --- reflected impedance annotation bracket on the primary terminals ---
\draw[engblue, decorate, decoration={brace,amplitude=5pt}] (2.55,-0.15) -- (2.55,2.55);
\node[engblue, anchor=east, align=right] at (2.2,-0.55)
  {$Z_{\text{ref}}=Z_L/n^2$};
\end{tikzpicture}
\caption[Transformer impedance matching]{An ideal two-winding
transformer of turns ratio $n=N_2/N_1$ presents the secondary load $Z_L$ to
the primary as $Z_{\text{ref}}=Z_L/n^2$. Choosing $n$ so that $Z_{\text{ref}}$
equals the source resistance $R_s$ matches the load to the source for maximum
power transfer without dissipating power in a matching resistor. The same
square-of-turns-ratio scaling is what lets an audio output transformer drive a
low-impedance loudspeaker from a high-impedance amplifier stage, and what sets
the impedance step in a radio-frequency matching network.}
\label{fig:vol04ch09:transformer-match}
\end{figure}


\paragraph{Sallen-Key corner and damping from component values.}
A unity-gain Sallen-Key low-pass filter (Figure~\ref{fig:vol04ch09:sallen-key})
has a transfer function whose denominator, derived in the case study below,
is $1 + j\omega(R_1+R_2)C_2 + (j\omega)^2 R_1 R_2 C_1 C_2$. Matching this to
the standard second-order form $1 + (j\omega)/(Q\omega_0) + (j\omega/\omega_0)^2$
gives a corner $\omega_0 = 1/\sqrt{R_1 R_2 C_1 C_2}$ and a quality factor
$Q = \sqrt{R_1 R_2 C_1 C_2}/((R_1+R_2)C_2)$. Take equal resistors
$R_1 = R_2 = R = 11.3\,\si{\kilo\ohm}$ and capacitors $C_1 = 22\,\si{\nano
\farad}$, $C_2 = 11\,\si{\nano\farad}$. Then $\omega_0 = 1/(R\sqrt{C_1 C_2}) =
1/(11.3\times 10^{3}\cdot\sqrt{22\times 10^{-9}\cdot 11\times 10^{-9}}) =
5.69\times 10^{3}\,\si{\per\second}$, so $f_0 = 906\,\si{\hertz}$, and with
equal resistors the quality factor simplifies to $Q = \tfrac{1}{2}\sqrt{C_1/C_2}
= \tfrac{1}{2}\sqrt{2} = 0.707$, the maximally-flat Butterworth value. The
capacitor ratio sets the damping while the resistor and capacitor product set
the corner, which is why a Sallen-Key stage is tuned by choosing the ratio
$C_1/C_2$ first for the response shape, then scaling both for the frequency.

\paragraph{Aliasing rejection of a two-pole anti-alias filter.}
An analog-to-digital converter samples at $f_s = 8\,\si{\kilo\hertz}$, so its
Nyquist frequency is $f_N = 4\,\si{\kilo\hertz}$; content above $f_N$ folds
back into the band and corrupts the data. A second-order Butterworth low-pass
with corner $f_c = 1\,\si{\kilo\hertz}$ has magnitude $|H| = 1/\sqrt{1 +
(f/f_c)^4}$. At the Nyquist frequency $f/f_c = 4$, so $|H| = 1/\sqrt{1 + 256}
= 1/16.0 = 0.0624$, that is $20\log_{10}(0.0624) = -24.1\,\si{\decibel}$;
the first frequency that aliases onto a DC signal is $f_s = 8\,\si{\kilo
\hertz}$, where $f/f_c = 8$ gives $|H| = 1/\sqrt{1 + 4096} = 0.0156$, that is
$-36.1\,\si{\decibel}$. A single-pole filter at the same corner would reach
only $-12\,\si{\decibel}$ at $f_N$ and $-18\,\si{\decibel}$ at $f_s$,
half the slope, which is why anti-alias filters ahead of a converter are at
least second order. The numbers say that a tone at $7.5\,\si{\kilo\hertz}$,
arriving $24\,\si{\decibel}$ down from the two-pole filter and then aliasing
to $0.5\,\si{\kilo\hertz}$, lands inside the pass band attenuated by only that
$24\,\si{\decibel}$; if the converter has more than about four bits of
resolution the alias is visible, and the filter order or corner must be
raised to push it below the least significant bit.

\paragraph{Wheatstone-bridge sensitivity and the small-signal output.}
A quarter-bridge strain gauge places one variable resistor $R(1+x)$ in a
four-arm bridge of nominal resistance $R$, excited by a voltage $V_{\text{exc}}$.
The bridge output is the difference of two dividers,
\[
v_o = V_{\text{exc}}\left(\frac{R(1+x)}{R(1+x)+R} - \frac{R}{R+R}\right)
= V_{\text{exc}}\left(\frac{1+x}{2+x} - \frac{1}{2}\right)
= V_{\text{exc}}\,\frac{x}{2(2+x)}.
\]
For the small fractional change $x = \text{GF}\cdot\varepsilon$ a strain gauge
delivers (gauge factor $\text{GF}\approx 2$), the output linearises to
$v_o \approx V_{\text{exc}}\,x/4$. With $V_{\text{exc}} = 5\,\si{\volt}$ and a
strain of $\varepsilon = 10^{-3}$ (one millistrain, a heavily loaded steel
member near its working limit), $x = 2\times 10^{-3}$ and $v_o = 5\cdot
2\times 10^{-3}/4 = 2.5\,\si{\milli\volt}$. The bridge converts a part-per-
thousand resistance change into a few millivolts, a signal far too small to
digitise directly and the reason the next stage is a high-gain instrumentation
amplifier. The residual nonlinearity from the dropped $x/2$ term is
$x/4 \approx 5\times 10^{-4}$, half a part per thousand, negligible against
the gauge's own tolerance but a term a half-bridge or full-bridge arrangement
cancels when the application needs it.

\paragraph{Output-impedance reduction by feedback.}
A voltage follower built from an op-amp with open-loop output resistance
$R_{o} = 100\,\si{\ohm}$ and open-loop gain $A = 10^{5}$ presents a
closed-loop output resistance reduced by the return difference,
$R_{\text{out}} = R_{o}/(1 + A\beta)$. For a follower $\beta = 1$, so
$R_{\text{out}} = 100/(1 + 10^{5}) = 1.0\,\si{\milli\ohm}$ at DC. The
reduction tracks the loop gain, which falls with frequency at
$-20\,\si{\decibel}$ per decade above the op-amp's dominant pole, so by the
gain-bandwidth frequency the output impedance has climbed back toward $R_o$;
a follower driving a fast load looks like a near-ideal source at audio
frequencies and like a $100\,\si{\ohm}$ source in the megahertz range. This
frequency-dependent output impedance is why a buffer that measures clean on a
multimeter can still ring when it drives a capacitive load, the load
resonating with the rising output inductance the feedback no longer
suppresses.

\paragraph{Step response by Laplace transform and partial fractions.}
A series RC stage with $R = 10\,\si{\kilo\ohm}$ and $C = 47\,\si{\nano
\farad}$ is driven from rest by a $5\,\si{\volt}$ step. Working in the
$s$-domain, the source is $5/s$, the capacitor impedance is $1/(sC)$, and
the capacitor voltage is the divider
\[
V_C(s) = \frac{5}{s}\cdot\frac{1/(sC)}{R + 1/(sC)}
= \frac{5}{s}\cdot\frac{1}{1 + sRC}
= \frac{5}{s(1 + s\tau)},
\quad \tau = RC = 470\,\si{\micro\second}.
\]
Partial fractions split the product of poles into a sum,
$5/[s(1 + s\tau)] = 5/s - 5\tau/(1 + s\tau) = 5/s - 5/(s + 1/\tau)$, and
each term inverts by inspection: $5/s \to 5$ and $5/(s + 1/\tau) \to
5e^{-t/\tau}$, giving $v_C(t) = 5(1 - e^{-t/\tau})$, the familiar charging
curve, now obtained by algebra rather than by guessing the exponential.
The method earns its keep when the network is second order or higher,
where the guess-and-match approach grows unwieldy and the partial-fraction
recipe stays mechanical. The pole at $s = -1/\tau = -2130\,\si{\per
\second}$ is the single natural frequency, sitting on the negative real
axis at the reciprocal of the time constant; the pole at the origin is the
step itself, the forced DC term.

\paragraph{Harmonic content and RMS of a clipped sine.}
An amplifier overdriven into hard clipping turns a sinusoid into a flat-
topped wave, and the Fourier picture says what that costs. Take a
$1\,\si{\volt}$-amplitude sine symmetrically clipped at $\pm 0.8\,\si{
\volt}$. The clipping is an odd, half-wave-symmetric distortion, so it adds
only odd harmonics. Estimating the first few coefficients of the clipped
wave by projecting it onto the harmonics gives a fundamental of about
$0.94\,\si{\volt}$ amplitude, a third harmonic of about $0.08\,\si{\volt}$,
and a fifth of about $0.02\,\si{\volt}$, with higher terms falling fast.
The total harmonic distortion is then $\text{THD} =
\sqrt{0.08^{2} + 0.02^{2}}/0.94 = 0.083/0.94 = 8.8\,\%$, a figure an
audible distortion test would flag at once. The RMS value of the clipped
wave is the root-sum-square of the harmonic RMS values, $\sqrt{(0.94^{2} +
0.08^{2} + 0.02^{2})/2} = 0.667\,\si{\volt}$, slightly below the
$1/\sqrt{2} = 0.707\,\si{\volt}$ of the unclipped sine, the small power
loss the flat top represents. The lesson is that even gentle clipping,
invisible on a casual oscilloscope trace, plants harmonics a spectrum
analyser reads immediately, which is why distortion is specified in the
frequency domain and not from the time-domain waveform.

\paragraph{Cascade of two stages by the transmission matrix.}
Two L-sections cascade into a ladder, and the transmission matrix turns
the cascade into a product. A single L-section of a series resistor $R$
followed by a shunt capacitor $C$ has the transmission matrix that is the
product of the series and shunt elementary blocks,
\[
T = \begin{bmatrix} 1 & R \\ 0 & 1 \end{bmatrix}
\begin{bmatrix} 1 & 0 \\ j\omega C & 1 \end{bmatrix}
= \begin{bmatrix} 1 + j\omega RC & R \\ j\omega C & 1 \end{bmatrix}.
\]
Two identical such sections in cascade have transmission matrix $T^{2}$,
and the voltage transfer ratio of the loaded cascade is read from the
combined matrix. With $R = 1\,\si{\kilo\ohm}$ and $C = 10\,\si{\nano
\farad}$ each section corners near $f_c = 1/(2\pi RC) = 15.9\,\si{\kilo
\hertz}$, and the cascade rolls off at $-40\,\si{\decibel}$ per decade
well above that corner. The product also exposes the loading the second
section places on the first: the off-diagonal $B$ term of $T^{2}$ is not
simply twice the single-section value, so a designer who cascades two
identical RC sections expecting their corner frequencies to stay put finds
the combined response sags below the naive product of two isolated
first-order responses, the interstage-loading error that a buffer between
the sections removes. The transmission-matrix product carries that loading
automatically, which is its advantage over multiplying isolated transfer
functions.

\paragraph{Sampled-noise floor and the converter bit budget.}
A sample-and-hold samples a sensor voltage onto a $C = 10\,\si{\pico
\farad}$ hold capacitor at room temperature, $T = 300\,\si{\kelvin}$. The
sampled thermal noise is $\sqrt{k_B T/C} = \sqrt{1.38\times 10^{-23}\cdot
300/10\times 10^{-12}} = \sqrt{4.14\times 10^{-10}} = 20.4\,\si{\micro
\volt}$ RMS, independent of the switch resistance that charged the
capacitor. For a converter with a $2\,\si{\volt}$ full-scale range, the
noise floor in least-significant-bit terms tells how many bits are usable:
the peak-to-peak noise is roughly six times the RMS value, about
$122\,\si{\micro\volt}$, so the noise spans one least significant bit when
$2/2^{N} = 122\,\si{\micro\volt}$, that is $N = \log_2(2/122\times
10^{-6}) = 14.0$ bits. A converter of more than fourteen bits sampling
through this hold capacitor digitises noise in its lowest bits; to use a
sixteen-bit converter the designer raises the hold capacitor to
$160\,\si{\pico\farad}$, dropping the floor by a factor of four to
$5.1\,\si{\micro\volt}$, at the cost of a slower acquisition through the
same switch resistance. The kTC floor, not the converter, sets the
resolution the chain can deliver.

\paragraph{Input impedance of a two-port from open and short tests.}
An unknown passive two-port is characterised on the bench by two
impedance measurements at its input port: with the output open-circuited
the input impedance reads $Z_{\text{oc}} = z_{11} = 600\,\si{\ohm}$, and
with the output short-circuited it reads $Z_{\text{sc}} = z_{11} -
z_{12}z_{21}/z_{22} = 150\,\si{\ohm}$. The geometric mean of the two is
the image impedance, $\sqrt{Z_{\text{oc}}Z_{\text{sc}}} = \sqrt{600\cdot
150} = 300\,\si{\ohm}$, the impedance that, terminating the output, makes
the input impedance equal to itself, and the value a designer matches a
source to for maximum power transfer through the network. The ratio
$Z_{\text{oc}}/Z_{\text{sc}} = 4$ fixes the magnitude of the transmission
through the box: a reciprocal network ($z_{12} = z_{21}$) with these two
numbers has $z_{12}z_{21}/z_{22} = Z_{\text{oc}} - Z_{\text{sc}} =
450\,\si{\ohm}\cdot z_{22}$, so a third measurement, the output impedance
with the input open, closes the parameter table. The open-and-short pair
is the oldest two-port characterisation and still the first an engineer
reaches for at a bench without a network analyser.

\begin{mastery}
The phasor calculus of Section~\ref{sec:vol04ch09:ac} is the steady-state
slice of the Laplace transform. Replacing $j\omega$ by the complex
frequency $s = \sigma + j\omega$ turns each element's impedance into a
function of $s$: $Z_{R} = R$, $Z_{L} = sL$, $Z_{C} = 1/(sC)$. Any linear
circuit then has a transfer function $H(s) = V_{\text{out}}(s)/
V_{\text{in}}(s)$ that is a ratio of polynomials in $s$,
\[
H(s) = K\,\frac{(s-z_{1})(s-z_{2})\cdots}{(s-p_{1})(s-p_{2})\cdots},
\]
whose roots of the numerator are the \engterm{zeros} and roots of the
denominator the \engterm{poles}. The poles are the natural frequencies
of the circuit: a pole at $s = -1/\tau$ is a first-order decay of time
constant $\tau$; a complex-conjugate pole pair at $-\zeta\omega_{n}\pm
j\omega_{d}$ is the underdamped ringing of
Section~\ref{sec:vol04ch09:transient}. A circuit is stable when every
pole lies in the left half
of the $s$-plane ($\Re(p_{k}) < 0$); a pole on the imaginary axis is a
sustained oscillation, and a pole in the right half is exponential
growth. The steady-state frequency response is the transfer function
evaluated on the imaginary axis, $H(j\omega)$, which is why the Bode
plot and the pole-zero map carry the same information. The state-space
formulation of Volume IX generalises this from a single transfer
function to a matrix of them; the pole locations become the eigenvalues
of the system matrix.
\end{mastery}

\section{An end-to-end case study: a regulated bench supply}\pathtag{standard}
\label{sec:vol04ch09:psu-case}

The chapter's machinery earns its keep when an engineer must turn a
specification into a working power supply with every component value
chosen and every margin checked. We carry a small linear bench supply
from the wall socket to the load: a transformer secondary, a full-wave
bridge rectifier, a reservoir capacitor, and a linear regulator feeding
a fixed load. The supply is a good vehicle because every section of the
chapter appears in it once. The rectifier and reservoir are a transient
problem; the ripple is an AC steady-state problem; the regulator
headroom is a DC operating-point problem; the dissipation is a thermal
problem; and the verdict is a number an engineer would sign before
ordering parts. Figure~\ref{fig:vol04ch09:psu-chain} sketches the chain
and the waveform at each stage.

% Figure: the signal chain of a regulated DC power supply, transformer
% secondary through rectifier, reservoir capacitor, and linear regulator to
% load, with the rectified-and-smoothed waveform sketched beneath. Plain TikZ.
\begin{figure}[htbp]
\centering
\begin{tikzpicture}[
  scale=0.92, transform shape,
  line width=0.8pt,
  font=\scriptsize,
  blk/.style={draw, fill=white, minimum width=20mm, minimum height=10mm,
    align=center},
  >={Stealth},
]
% === block chain ===
\node[blk] (xfmr) at (0,0) {transformer\\secondary\\$12.6\,V_{\text{rms}}$};
\node[blk] (rect) at (3,0) {full-wave\\bridge};
\node[blk] (cap)  at (6,0) {reservoir\\$C$};
\node[blk] (reg)  at (9,0) {linear\\regulator};
\node[blk] (load) at (12,0) {load\\$R_L$};
\draw[->, engblue] (xfmr) -- (rect);
\draw[->, engblue] (rect) -- (cap);
\draw[->, engblue] (cap)  -- (reg);
\draw[->, engblue] (reg)  -- (load);
% === waveform panel beneath ===
\begin{scope}[yshift=-3.2cm]
  \draw[->, enggray] (0,-1.1) -- (12.3,-1.1) node[right]{$t$};
  \draw[->, enggray] (0,-1.1) -- (0,1.3) node[above]{$v$};
  % full-wave rectified ripple (humps) at the cap node: approximate by a
  % sawtooth-like ripple riding on a DC level. Draw as line segments.
  \def\dc{0.7}
  % rising charge spikes then linear droop, four periods
  \draw[engorange, thick]
    (0,0.55) -- (0.6,0.9) -- (1.8,0.55) -- (2.4,0.9) -- (3.6,0.55)
    -- (4.2,0.9) -- (5.4,0.55) -- (6.0,0.9) -- (7.2,0.55);
  \node[engorange, anchor=west] at (7.4,0.9) {ripple on $C$};
  % regulated flat output line
  \draw[englime, very thick] (0,0.05) -- (7.2,0.05);
  \node[englime, anchor=west] at (7.4,0.05) {regulated $V_{\text{out}}$};
  % ripple band annotation
  \draw[enggray, <->] (7.05,0.55) -- (7.05,0.9);
  \node[enggray, anchor=west, font=\tiny] at (7.4,0.55) {$\Delta V_r$};
\end{scope}
\end{tikzpicture}
\caption[Regulated DC supply signal chain]{The signal chain of a
regulated DC power supply. The transformer steps the mains down to a
low-voltage secondary; the bridge rectifies it to a unidirectional
pulsating waveform; the reservoir capacitor $C$ holds charge between
conduction peaks, leaving a sawtooth ripple $\Delta V_r$ on the
unregulated rail; the linear regulator drops the rippled rail to a clean
fixed $V_{\text{out}}$, absorbing the ripple and the line and load
variation in its headroom. Each stage is sized in the case study of
Section~\ref{sec:vol04ch09:psu-case}.}
\label{fig:vol04ch09:psu-chain}
\end{figure}


\subsection{Specification}

The supply must deliver $V_{\text{out}} = 5.0\,\si{\volt}$ at a load
current up to $I_{L} = 0.50\,\si{\ampere}$, holding the output within
$\pm 1\,\%$ over the full load range and over a mains voltage that
varies by $\pm 10\,\%$. The mains is $230\,\si{\volt}$ RMS at
$50\,\si{\hertz}$. We choose a linear three-terminal regulator that
needs a dropout headroom of $2.0\,\si{\volt}$, meaning its input must
stay at least $2.0\,\si{\volt}$ above its $5.0\,\si{\volt}$ output, that
is above $7.0\,\si{\volt}$ at the bottom of the ripple trough and at the
low end of the mains swing.

\subsection{Transformer secondary}

The unregulated rail before the regulator must stay above
$7.0\,\si{\volt}$ at its lowest, allowing a ripple amplitude and a
margin. We target a nominal unregulated rail of about $10\,\si{\volt}$
at the top of the ripple. A full-wave bridge drops two diode forward
voltages, $2V_{D}\approx 1.4\,\si{\volt}$, below the secondary peak, so
the secondary peak must reach $10 + 1.4 = 11.4\,\si{\volt}$. The
secondary RMS voltage is the peak divided by $\sqrt{2}$, so
$V_{\text{sec}} = 11.4/\sqrt{2} = 8.1\,\si{\volt}$ RMS at nominal mains.
A standard $9\,\si{\volt}$ secondary at nominal mains gives a peak of
$9\sqrt{2} = 12.7\,\si{\volt}$ and a post-bridge peak of $11.3\,\si{
\volt}$; at the $-10\,\%$ mains corner this falls to $11.3\cdot 0.90 =
10.2\,\si{\volt}$, which we carry forward as the worst-case top of the
unregulated rail.

\subsection{Rectifier and reservoir capacitor}

The bridge charges the reservoir capacitor to the post-bridge peak on
each half-cycle, then the load discharges it between peaks. A full-wave
bridge refreshes the capacitor twice per mains cycle, so the discharge
interval is half the mains period, $T/2 = 1/(2\cdot 50) = 10\,\si{\milli
\second}$. During that interval the load draws a near-constant
$0.50\,\si{\ampere}$, so the ripple voltage on the capacitor follows
from $i = C\,\dd v/\dd t$ in its averaged form $\Delta V_{r} = I_{L}\,
\Delta t/C$. Setting a ripple budget of $\Delta V_{r} = 1.0\,\si{\volt}$
peak-to-peak,
\[
C = \frac{I_{L}\,\Delta t}{\Delta V_{r}} = \frac{0.50\cdot 10\times
10^{-3}}{1.0} = 5.0\times 10^{-3}\,\si{\farad} = 5000\,\si{\micro\farad}.
\]
We select a standard $4700\,\si{\micro\farad}$ part and accept a ripple
of $\Delta V_{r} = 0.50\cdot 0.010/4.7\times 10^{-3} = 1.06\,\si{\volt}$.
The bottom of the ripple at the worst mains corner is then $10.2 - 1.06
= 9.1\,\si{\volt}$, comfortably above the $7.0\,\si{\volt}$ the regulator
demands. The reservoir capacitor carries a large pulsed charging
current: the diodes conduct only near the peak, for a conduction angle
of perhaps $3\,\si{\milli\second}$ out of the $10\,\si{\milli\second}$
interval, so the peak charging current is roughly $I_{L}\cdot(10/3)
\approx 1.7\,\si{\ampere}$, and the capacitor's ripple-current rating
and the bridge's surge rating must cover it.

\subsection{Regulator headroom and load regulation}

The regulator holds its $5.0\,\si{\volt}$ output as long as its input
stays above $7.0\,\si{\volt}$. The input ranges from the worst-case
trough $9.1\,\si{\volt}$ up to the nominal-mains peak $11.3\,\si{\volt}$,
so the headroom ranges from a worst case of $9.1 - 5.0 = 4.1\,\si{\volt}$
up to $6.3\,\si{\volt}$, with a $2.1\,\si{\volt}$ margin over the
$2.0\,\si{\volt}$ dropout even at the trough. The output stays in
regulation across the full mains and load range. The load regulation,
the output change from no load to full load, comes from the regulator's
finite output impedance; a typical linear regulator holds its output to
within a few millivolts across $0.50\,\si{\ampere}$, well inside the
$\pm 1\,\%$ ($\pm 50\,\si{\milli\volt}$) specification.

\subsection{Efficiency and thermal verdict}

The linear regulator dissipates the headroom voltage times the load
current as heat. At the nominal-mains peak input of $11.3\,\si{\volt}$
the headroom is $11.3 - 5.0 = 6.3\,\si{\volt}$, and at full load the
regulator dissipates
\[
P_{\text{reg}} = (V_{\text{in}}-V_{\text{out}})\,I_{L} = 6.3\cdot 0.50 =
3.15\,\si{\watt}.
\]
The useful output power is $V_{\text{out}}I_{L} = 5.0\cdot 0.50 =
2.5\,\si{\watt}$, so the regulator-stage efficiency is $2.5/(2.5+3.15) =
44\,\%$; the transformer and rectifier losses lower the wall-plug
efficiency further, to roughly $35\,\%$. The linear topology trades
efficiency for simplicity and low output noise. A regulator dissipating
$3.15\,\si{\watt}$ needs a heat sink: with a junction-to-ambient path
that must hold the junction below, say, $125\,\si{\celsius}$ in a
$40\,\si{\celsius}$ ambient, the allowed thermal resistance is
$(125-40)/3.15 = 27\,\si{\celsius}/\si{\watt}$, which a small clip-on
heat sink meets. A switching regulator in place of the linear part would
reach $85\,\%$ efficiency and need no heat sink, at the cost of switching
noise on the output and a more demanding layout; for a quiet bench
reference the linear supply is the right choice, and the case study ends
with a supply an engineer would build.

\section{A second case study: a switching buck converter}\pathtag{standard}
\label{sec:vol04ch09:buck-case}

The linear supply of the previous case study throws away the headroom
voltage as heat; its $44\,\%$ regulator-stage efficiency is the price of
simplicity. When the load is large or the input-output gap is wide, that
waste becomes intolerable, and the design moves to a switching converter
that stores energy in an inductor and meters it to the load. We carry a
\engterm{buck} (step-down) converter end to end, choosing the switching
frequency, the inductor, and the output capacitor, checking the
continuous-conduction boundary and the ripple, estimating the efficiency,
and noting what the control loop must do. The buck converter is the
counterpoint to the linear supply: the same specification, a different
topology, and a different set of trade-offs. Figure~\ref{fig:vol04ch09:buck-topology}
draws the power stage.

% Figure: synchronous (or diode) buck converter power stage drawn with plain
% TikZ paths and nodes (no circuitikz). Input source, high-side switch, catch
% diode, inductor, output capacitor, and load, with the switching node marked.
\begin{figure}[htbp]
\centering
\begin{tikzpicture}[
  line width=0.8pt,
  font=\scriptsize,
  >={Stealth},
]
% --- rails ---
\coordinate (vinp) at (0,3);     % input positive
\coordinate (gnd)  at (0,0);     % ground rail left
\coordinate (sw)   at (3,3);     % top of switch / switching node feed
\coordinate (lx)   at (3,1.6);   % switching node Lx
\coordinate (lout) at (6.5,1.6); % inductor right terminal
\coordinate (capt) at (6.5,3);   % unused upper; keep layout
\coordinate (outp) at (8.5,1.6); % output positive node
\coordinate (gndr) at (8.5,0);   % ground rail right

% --- input source (battery symbol) ---
\draw (vinp) -- (sw);
\draw (gnd) -- (0,1.5);
\draw[line width=1.4pt] (-0.25,1.5) -- (0.25,1.5);     % long plate
\draw[line width=0.8pt] (-0.15,1.3) -- (0.15,1.3);     % short plate
\draw (-0.25,1.65) -- (0.25,1.65);                     % long plate (upper cell)
\draw[line width=0.8pt] (-0.15,1.85) -- (0.15,1.85);
\draw (0,1.85) -- (vinp);
\node[anchor=east] at (-0.3,1.6) {$V_{\text{in}}$};

% --- high-side switch (drawn as a labelled box with an arrow contact) ---
\draw (sw) -- (3,2.4);
\draw[fill=white] (2.6,1.9) rectangle (3.4,2.4);
\node at (3,2.15) {\tiny SW};
\draw (3,1.9) -- (lx);
\node[anchor=west, engblue] at (3.5,2.55) {duty $D$};
\draw[engblue,->] (3.45,2.5) -- (3.05,2.25);

% --- catch diode from ground to switching node ---
\draw (lx) -- (3,1.6);
\draw (3,0) -- (lx);
% diode triangle (cathode up toward Lx)
\draw[fill=white] (2.75,0.7) -- (3.25,0.7) -- (3,1.1) -- cycle;
\draw[line width=1.2pt] (2.75,1.1) -- (3.25,1.1);   % cathode bar
\node[anchor=west] at (3.35,0.9) {$D$};

% --- inductor (series of bumps) ---
\draw (lx) -- (3.7,1.6);
\draw (3.7,1.6)
  .. controls (3.85,2.05) and (4.15,2.05) .. (4.3,1.6)
  .. controls (4.45,2.05) and (4.75,2.05) .. (4.9,1.6)
  .. controls (5.05,2.05) and (5.35,2.05) .. (5.5,1.6)
  .. controls (5.65,2.05) and (5.95,2.05) .. (6.1,1.6);
\draw (6.1,1.6) -- (lout);
\node[anchor=south] at (4.9,2.0) {$L$};
\node[engorange, anchor=north] at (3,1.45) {$L_x$};

% --- output capacitor ---
\draw (lout) -- (outp);
\draw (7.4,1.6) -- (7.4,1.0);
\draw[line width=1.4pt] (7.15,1.0) -- (7.65,1.0);
\draw[line width=1.4pt] (7.15,0.8) -- (7.65,0.8);
\draw (7.4,0.8) -- (7.4,0) ;
\node[anchor=west] at (7.7,0.9) {$C_{\text{out}}$};

% --- load resistor (zigzag) ---
\draw (outp) -- (8.5,1.6);
\draw (8.5,1.6) -- (8.5,1.35)
  -- (8.35,1.27) -- (8.65,1.13) -- (8.35,0.99) -- (8.65,0.85)
  -- (8.5,0.77) -- (8.5,0);
\node[anchor=west] at (8.7,1.1) {$R_L$};
\node[anchor=west, englime] at (8.7,1.7) {$V_{\text{out}}$};

% --- ground rail ---
\draw (gnd) -- (gndr);
\draw (3,1.6) -- (3,0);   % ensure diode/cap tie to ground (visual)
% ground symbol
\draw (4.2,0) -- (4.2,-0.25);
\draw (4.0,-0.25) -- (4.4,-0.25);
\draw (4.07,-0.37) -- (4.33,-0.37);
\draw (4.15,-0.49) -- (4.25,-0.49);

\end{tikzpicture}
\caption[Buck converter power stage]{The buck (step-down) converter
power stage. The high-side switch SW connects the inductor to $V_{\text{in}}$
for a fraction $D$ of each switching period; the catch diode $D$ carries the
inductor current for the remaining $1-D$. The switching node $L_x$ swings
between $V_{\text{in}}$ and roughly ground, and the inductor and output
capacitor average that square wave to a steady $V_{\text{out}} = D\,
V_{\text{in}}$. Component values are chosen in the case study of
Section~\ref{sec:vol04ch09:buck-case}.}
\label{fig:vol04ch09:buck-topology}
\end{figure}


\subsection{Specification and operating principle}

The converter steps a $V_{\text{in}} = 12\,\si{\volt}$ input down to
$V_{\text{out}} = 5.0\,\si{\volt}$ at a load current up to $I_{L} =
2.0\,\si{\ampere}$. A high-side switch connects the inductor to the input
for a fraction $D$ of each switching period $T_{s}$, then opens and lets
the catch diode (or a synchronous switch) carry the inductor current for
the remaining $1-D$. In continuous conduction the inductor current never
falls to zero, and volt-second balance across the inductor over one
period fixes the duty cycle: the inductor sees $(V_{\text{in}}-V_{\text{
out}})$ for $DT_{s}$ and $-V_{\text{out}}$ for $(1-D)T_{s}$, and the two
volt-second areas must cancel in steady state,
\[
(V_{\text{in}}-V_{\text{out}})\,DT_{s} = V_{\text{out}}\,(1-D)T_{s}
\quad\Longrightarrow\quad
D = \frac{V_{\text{out}}}{V_{\text{in}}} = \frac{5.0}{12} = 0.42.
\]
The output voltage is the input scaled by the duty cycle, set by the
fraction of time the switch is closed rather than by a dissipative drop.
A real diode-rectified buck loses one diode forward drop, raising the
ideal duty slightly, but the volt-second relation is the design anchor.

\subsection{Switching frequency and inductor sizing}

The switching frequency trades inductor size against switching loss: a
higher frequency shrinks the inductor and capacitor but costs more energy
per switching transition. We pick $f_{s} = 250\,\si{\kilo\hertz}$, a
common value for a converter of this power, giving a period $T_{s} =
1/f_{s} = 4.0\,\si{\micro\second}$. The inductor is sized from the
peak-to-peak ripple current $\Delta I_{L}$ we are willing to accept. With
the switch closed, the inductor voltage is $V_{\text{in}}-V_{\text{out}} =
7.0\,\si{\volt}$ for the on-time $DT_{s} = 0.42\cdot 4.0 = 1.67\,\si{\micro
\second}$, and $v = L\,\dd i/\dd t$ in averaged form gives
\[
\Delta I_{L} = \frac{(V_{\text{in}}-V_{\text{out}})\,DT_{s}}{L}.
\]
A standard rule sets the ripple to about $30\,\%$ of the maximum load
current, $\Delta I_{L} = 0.30\cdot 2.0 = 0.60\,\si{\ampere}$. Solving for
$L$,
\[
L = \frac{(V_{\text{in}}-V_{\text{out}})\,DT_{s}}{\Delta I_{L}}
= \frac{7.0\cdot 1.67\times 10^{-6}}{0.60}
= 19.5\,\si{\micro\henry}.
\]
We select a standard $22\,\si{\micro\henry}$ part, which lowers the ripple
to $\Delta I_{L} = 7.0\cdot 1.67\times 10^{-6}/22\times 10^{-6} = 0.53\,
\si{\ampere}$. The peak inductor current is $I_{L}+\Delta I_{L}/2 = 2.0 +
0.27 = 2.27\,\si{\ampere}$, the value the inductor's saturation rating and
the switch's current rating must clear with margin.
Figure~\ref{fig:vol04ch09:buck-ripple} shows the switching-node voltage and
the triangular inductor current.

% Figure: inductor current and output-voltage ripple waveforms of a buck
% converter in continuous conduction mode. Triangular inductor current riding
% on its average, switching-node voltage above it. Plain coordinate paths.
\begin{figure}[htbp]
\centering
\begin{tikzpicture}[
  line width=0.8pt,
  font=\scriptsize,
  >={Stealth},
]
% === switching-node voltage (square wave) ===
\begin{scope}[yshift=2.7cm]
  \draw[->, enggray] (0,-0.1) -- (8.4,-0.1) node[right]{$t$};
  \draw[->, enggray] (0,-0.1) -- (0,1.5) node[above]{$v_{L_x}$};
  % three periods, duty D about 0.42
  \def\Ton{0.85}
  \def\Tper{2.0}
  \draw[engblue, thick]
    (0,1.1) -- (\Ton,1.1) -- (\Ton,0) -- (\Tper,0)
    -- (\Tper,1.1) -- ({\Tper+\Ton},1.1) -- ({\Tper+\Ton},0) -- ({2*\Tper},0)
    -- ({2*\Tper},1.1) -- ({2*\Tper+\Ton},1.1) -- ({2*\Tper+\Ton},0) -- ({3*\Tper},0)
    -- ({3*\Tper},1.1) -- ({3*\Tper+\Ton},1.1);
  \node[engblue, anchor=west, font=\tiny] at (6.3,1.1) {$\approx V_{\text{in}}$};
  \draw[<->, enggray, font=\tiny] (0,-0.05) -- (\Ton,-0.05);
  \node[enggray, font=\tiny, anchor=north] at (0.42,-0.05) {$DT$};
  \draw[<->, enggray, font=\tiny] (\Ton,1.25) -- (\Tper,1.25);
  \node[enggray, font=\tiny, anchor=south] at (1.42,1.25) {$(1-D)T$};
\end{scope}

% === inductor current (triangular ripple) ===
\draw[->, enggray] (0,-0.1) -- (8.4,-0.1) node[right]{$t$};
\draw[->, enggray] (0,-0.1) -- (0,2.0) node[above]{$i_L$};
% average level
\draw[englime, thick, dashed] (0,1.0) -- (8.0,1.0);
\node[englime, anchor=west, font=\tiny] at (8.0,1.0) {$I_L$};
% triangular ripple: rise during DT, fall during (1-D)T
\def\Ton{0.85}
\def\Tper{2.0}
\def\ihi{1.45}
\def\ilo{0.55}
\draw[engorange, thick]
  (0,\ilo)
  -- (\Ton,\ihi) -- (\Tper,\ilo)
  -- ({\Tper+\Ton},\ihi) -- ({2*\Tper},\ilo)
  -- ({2*\Tper+\Ton},\ihi) -- ({3*\Tper},\ilo)
  -- ({3*\Tper+\Ton},\ihi);
% ripple band annotation
\draw[engred, <->] (\Ton,\ihi) -- (\Ton,\ilo);
\node[engred, anchor=west, font=\tiny] at (\Ton,1.0) {$\Delta I_L$};
\node[enggray, font=\tiny, anchor=north] at (4.0,-0.15) {time (switching periods)};
\end{tikzpicture}
\caption[Buck converter ripple waveforms]{Switching-node voltage (top)
and inductor current (bottom) of a buck converter in continuous conduction
mode. While the switch is closed for $DT$, the node sits near $V_{\text{in}}$
and the inductor current ramps up at $(V_{\text{in}}-V_{\text{out}})/L$; while
the switch is open for $(1-D)T$, the node falls toward ground and the current
ramps down at $V_{\text{out}}/L$. The current rides as a triangle of
peak-to-peak ripple $\Delta I_L$ on the average load current $I_L$; the output
capacitor integrates that triangle to a small voltage ripple. The boundary to
discontinuous conduction is reached when $\Delta I_L/2$ equals $I_L$, so the
trough touches zero.}
\label{fig:vol04ch09:buck-ripple}
\end{figure}


\subsection{Continuous-conduction boundary}

The converter stays in continuous conduction as long as the inductor
current trough stays above zero, $I_{L} > \Delta I_{L}/2$. The trough
touches zero at the boundary load current $I_{\text{crit}} = \Delta I_{L}/2
= 0.27\,\si{\ampere}$. Below that the converter enters
\engterm{discontinuous conduction}: the inductor current reaches zero
before the period ends, the diode stops conducting, and the duty-cycle-to-
voltage relation changes from the clean $D = V_{\text{out}}/V_{\text{in}}$
to a load-dependent expression. A converter that must regulate down to no
load either tolerates the mode change in its control law or adds a
synchronous switch that allows the inductor current to reverse, keeping
the conduction continuous at any load. Our specification holds the load
above $I_{\text{crit}}$ in normal operation, so the design sits in
continuous conduction across its rated range; the control note below
records what happens at light load.

\subsection{Output capacitor and voltage ripple}

The output capacitor absorbs the triangular ripple current that the
inductor delivers and the load does not take, leaving a small voltage
ripple. The charge delivered to the capacitor over the half-period when
the inductor current exceeds the average is the area of a triangle of base
$T_{s}/2$ and height $\Delta I_{L}/2$, so $\Delta Q = \tfrac{1}{2}\cdot
(T_{s}/2)\cdot(\Delta I_{L}/2) = \Delta I_{L}T_{s}/8$, and the
ideal-capacitor voltage ripple is
\[
\Delta V_{\text{out}} = \frac{\Delta Q}{C_{\text{out}}}
= \frac{\Delta I_{L}}{8 f_{s} C_{\text{out}}}.
\]
Setting a ripple budget of $\Delta V_{\text{out}} = 25\,\si{\milli\volt}$
(half a percent of the $5.0\,\si{\volt}$ output),
\[
C_{\text{out}} = \frac{\Delta I_{L}}{8 f_{s}\,\Delta V_{\text{out}}}
= \frac{0.53}{8\cdot 250\times 10^{3}\cdot 0.025}
= 10.6\,\si{\micro\farad}.
\]
A standard $22\,\si{\micro\farad}$ ceramic part gives margin against the
DC-bias capacitance loss of a class II dielectric and against the second
ripple term that the capacitor's equivalent series resistance contributes:
the ESR adds $\Delta V_{\text{ESR}} = \Delta I_{L}\cdot R_{\text{ESR}}$,
which for a $20\,\si{\milli\ohm}$ ceramic is $0.53\cdot 0.02 = 11\,\si{
\milli\volt}$, comparable to the charge-balance ripple and additive with
it. Choosing a low-ESR ceramic keeps the total inside the budget; an
electrolytic of the same capacitance would let the ESR term dominate and
blow the ripple spec.

\subsection{Efficiency and the control loop}

The buck converter's losses are the conduction loss in the switch and
inductor resistance, the switching loss in the switch, and the diode loss
in a non-synchronous design. Take a switch on-resistance $R_{\text{ds}} =
40\,\si{\milli\ohm}$, an inductor winding resistance $R_{L} = 30\,\si{
\milli\ohm}$, and a Schottky diode forward drop $V_{F} = 0.4\,\si{\volt}$.
The conduction loss in the switch over the on-time is $I_{L}^{2}R_{\text{
ds}}D = 4.0\cdot 0.04\cdot 0.42 = 67\,\si{\milli\watt}$; the diode
conducts for $1-D$ and dissipates $V_{F}I_{L}(1-D) = 0.4\cdot 2.0\cdot
0.58 = 0.46\,\si{\watt}$; the inductor copper loss is $I_{L}^{2}R_{L} =
4.0\cdot 0.03 = 0.12\,\si{\watt}$. Estimating the switching loss at
$\tfrac{1}{2}V_{\text{in}}I_{L}(t_{\text{on}}+t_{\text{off}})f_{s}$ with
transition times of $20\,\si{\nano\second}$ each gives $\tfrac{1}{2}\cdot
12\cdot 2.0\cdot 40\times 10^{-9}\cdot 250\times 10^{3} = 0.12\,\si{
\watt}$. The total loss is about $0.77\,\si{\watt}$ against an output power
of $V_{\text{out}}I_{L} = 10\,\si{\watt}$, an efficiency of $10/(10+0.77)
= 93\,\%$, with the diode the dominant loss. Replacing the diode with a
synchronous switch removes most of the $0.46\,\si{\watt}$ diode term and
lifts the efficiency toward $96\,\%$, the gain that justifies the extra
switch and its gate drive. The linear supply of the previous case study
reached $44\,\%$ on the same input-output ratio; the buck converter more
than doubles it, which is why every laptop and phone charger is a
switcher.

The output voltage is held against line and load disturbance by a
feedback loop that measures $V_{\text{out}}$, compares it to a reference,
and adjusts the duty cycle. The power stage from duty cycle to output is a
second-order system: the inductor and output capacitor form an LC filter
with a resonant frequency $f_{\text{LC}} = 1/(2\pi\sqrt{L C_{\text{out}}})
= 1/(2\pi\sqrt{22\times 10^{-6}\cdot 22\times 10^{-6}}) = 7.2\,\si{\kilo
\hertz}$, a complex-conjugate pole pair lightly damped by the load and
parasitic resistances. A controller that simply integrated the error would
ring at that pole pair; the working design adds a zero (a series
resistor-capacitor in the feedback compensator, the same pole-zero shaping
of the mastery box on transfer functions) to roll the loop gain through
the LC resonance with adequate phase margin. The crossover frequency is
set a decade below the switching frequency, around $25\,\si{\kilo\hertz}$
here, so the loop rejects load steps faster than the linear regulator can
while staying well clear of the switching ripple. The control synthesis
belongs to Volume IX; the circuit reasoning that sizes the LC filter and
locates its resonance belongs here, and it is the bridge between the two
volumes.

\begin{mastery}
The switching action that buys the buck converter its efficiency also makes
it a source of conducted and radiated interference. The switching-node
voltage of Figure~\ref{fig:vol04ch09:buck-ripple} is a trapezoidal wave with
fast edges, and its Fourier series carries harmonics of the switching
frequency $f_{s}$ whose envelope falls at $-20\,\si{\decibel}$ per decade up
to the corner $f_{1} = 1/(\pi t_{r})$ set by the edge rise time $t_{r}$,
then at $-40\,\si{\decibel}$ per decade beyond it. With $t_{r} = 20\,\si{
\nano\second}$ the first corner sits near $16\,\si{\mega\hertz}$, so
significant harmonic content reaches well into the band that conducted-
emissions standards police. The current the converter draws from its input
is pulsed at $f_{s}$, and that pulsed current flowing through the finite
impedance of the supply wiring is the conducted-emissions source; an input
filter (a small series inductor and a shunt capacitor, a low-pass of the
kind analysed in Section~\ref{sec:vol04ch09:ac}) attenuates it before it
leaves the board. Slowing the switch edges lowers the high-frequency
harmonics but raises the switching loss, the same edge-rate trade-off that
governs signal integrity in the failure section. The working designer reads
the harmonic envelope off the edge rate, places the input-filter corner
below the first troublesome harmonic, and confirms the result on a spectrum
analyser with a line-impedance stabilisation network, the bench instrument
that presents the standard source impedance the emissions limits assume.
\end{mastery}

\section{A third case study: a sensor signal-conditioning chain}\pathtag{standard}
\label{sec:vol04ch09:signal-chain-case}

The two power case studies move energy; this one moves information. A
sensor that produces a small analog voltage must be amplified, filtered,
and handed to an analog-to-digital converter without losing the
information the bit budget can carry or adding error the application
cannot tolerate. We carry a strain-gauge front end end to end: a
Wheatstone bridge source, an instrumentation amplifier for gain and
common-mode rejection, a second-order Sallen-Key anti-alias filter
designed to a stated stop-band figure, a noise and offset budget
referred to the converter's least significant bit, a Bode check of the
chain, and a verdict an engineer would sign before laying out the board.
The chain exercises every part of the chapter that the power supplies did
not: the op-amp feedback theory, the second-order filter transfer
function, and the AC frequency response.
Figure~\ref{fig:vol04ch09:signal-chain} sketches the four stages and the
signal level at each interface. The treatment follows the working
references on sensor conditioning
\cite{text:ch05-pallas-areny-webster2001, text:ch05-northrop2014}.

% Figure: block diagram of the sensor signal-conditioning chain, from a
% Wheatstone bridge strain sensor through an instrumentation amplifier and a
% second-order Sallen-Key anti-alias filter into an analog-to-digital
% converter. Plain TikZ boxes and arrows; no circuitikz.
\begin{figure}[htbp]
\centering
\begin{tikzpicture}[
  line width=0.8pt,
  blk/.style={draw, engblue, fill=white, minimum width=21mm, minimum height=12mm,
    align=center, font=\scriptsize},
  sig/.style={-{Latex[length=2mm]}, engblack},
  font=\small,
]
  \node[blk] (src) at (0,0) {Bridge\\sensor\\$2\,$mV/V};
  \node[blk] (ina) at (3.0,0) {Instr.\\amp\\$G=100$};
  \node[blk] (filt) at (6.0,0) {Sallen--Key\\$2$-pole\\$f_c=1\,$kHz};
  \node[blk] (adc) at (9.0,0) {ADC\\$16$-bit\\$f_s=8\,$kHz};
  \draw[sig] (src) -- (ina);
  \draw[sig] (ina) -- (filt);
  \draw[sig] (filt) -- (adc);
  % excitation arrow into the bridge
  \draw[sig, engorange] (0,1.4) -- (src.north);
  \node[font=\tiny, anchor=south, engorange] at (0,1.45) {$V_{\text{exc}}=5\,$V};
  % annotated signal levels under each interface
  \node[font=\tiny, anchor=north, enggray] at (1.5,-0.95) {$\le 10\,$mV};
  \node[font=\tiny, anchor=north, enggray] at (4.5,-0.95) {$\le 1\,$V};
  \node[font=\tiny, anchor=north, enggray] at (7.5,-0.95) {band-limited};
\end{tikzpicture}
\caption[Sensor signal-conditioning chain]{The four stages of the
signal-conditioning chain carried in the case study. A Wheatstone bridge
turns strain into a small differential voltage; an instrumentation amplifier
raises it by a factor of one hundred while rejecting the common-mode
excitation; a second-order Sallen--Key low-pass filter removes content above
the converter's Nyquist frequency; and the analog-to-digital converter
samples the band-limited result. Each stage is specified by a gain, a
bandwidth, and a noise contribution, and the chain is designed so that no
stage dominates the error budget set by the converter's least significant
bit.}
\label{fig:vol04ch09:signal-chain}
\end{figure}


\subsection{Specification}

The sensor is a quarter-bridge strain gauge of nominal arm resistance
$350\,\si{\ohm}$ and gauge factor $2.0$, excited at $V_{\text{exc}} =
5.0\,\si{\volt}$, measuring strain up to $\varepsilon_{\max} =
1.0\times 10^{-3}$ (one millistrain). The physical quantity of interest
varies slowly, with no real content above $1\,\si{\kilo\hertz}$. The
converter is a $16$-bit device with a $\pm 1.0\,\si{\volt}$ input range
sampling at $f_s = 8\,\si{\kilo\hertz}$. The chain must present the
full-scale strain as close to the converter's full scale as the budget
allows, hold the total referred-to-input error below one least
significant bit, and attenuate any content at or above the Nyquist
frequency $f_N = 4\,\si{\kilo\hertz}$ enough that an aliased tone falls
below that same least significant bit.

The converter's least significant bit spans $\text{LSB} = 2.0\,\si{\volt}
/2^{16} = 30.5\,\si{\micro\volt}$ at the input, the unit against which
every error in the chain is measured.

\subsection{Bridge source and required gain}

The quarter-bridge output at full strain, from the worked example above,
is $v_o \approx V_{\text{exc}}\,\text{GF}\,\varepsilon/4 = 5.0\cdot 2.0
\cdot 1.0\times 10^{-3}/4 = 2.5\,\si{\milli\volt}$ at the top of the
range. To map this $2.5\,\si{\milli\volt}$ full-scale signal onto the
converter's $1.0\,\si{\volt}$ full scale we need a gain of $G =
1.0/2.5\times 10^{-3} = 400$. The bridge's source resistance is the
$350\,\si{\ohm}$ bridge impedance, so the amplifier that follows must
present a high input impedance to avoid loading the bridge and a high
common-mode rejection to ignore the $2.5\,\si{\volt}$ common-mode
voltage that sits at the bridge midpoints. The instrumentation amplifier
of the earlier worked example is the standard answer; we set its gain to
$400$ with its single gain resistor and carry that gain forward.

\subsection{Anti-alias filter design}

The converter samples at $8\,\si{\kilo\hertz}$, so anything above
$4\,\si{\kilo\hertz}$ aliases. The signal has no real content above
$1\,\si{\kilo\hertz}$, so a low-pass filter cornered at $f_c =
1\,\si{\kilo\hertz}$ passes the signal and begins rolling off well before
Nyquist. We require that a full-scale interfering tone at $f_N =
4\,\si{\kilo\hertz}$ be attenuated below one least significant bit, a
ratio of $\text{LSB}/\text{full scale} = 1/2^{16} = -96\,\si{\decibel}$
in the worst case, relaxed in practice to the converter's effective
number of bits.

The Sallen-Key topology of Figure~\ref{fig:vol04ch09:sallen-key} builds
a second-order low-pass in one op-amp. Writing KCL at nodes A and B with
the virtual short ($v_B = v_{\text{out}}$ for the unity-gain follower)
gives the transfer function. The current balance at node A is
\[
\frac{v_{\text{in}}-v_A}{R_1} = \frac{v_A - v_{\text{out}}}{1/(j\omega C_1)}
+ \frac{v_A - v_B}{R_2},
\]
and at node B, where the op-amp input draws no current,
\[
\frac{v_A - v_B}{R_2} = \frac{v_B - 0}{1/(j\omega C_2)}.
\]
Setting $v_B = v_{\text{out}}$ and eliminating $v_A$ yields
\[
H(j\omega) = \frac{v_{\text{out}}}{v_{\text{in}}}
= \frac{1}{1 + j\omega(R_1+R_2)C_2 + (j\omega)^2 R_1 R_2 C_1 C_2}.
\]
This is the second-order low-pass form with $\omega_0 = 1/\sqrt{R_1 R_2
C_1 C_2}$ and $1/(Q\omega_0) = (R_1+R_2)C_2$, so $Q = \sqrt{R_1 R_2 C_1
C_2}/((R_1+R_2)C_2)$. Choosing equal resistors $R_1 = R_2 = R$ collapses
the damping to $Q = \tfrac{1}{2}\sqrt{C_1/C_2}$, set by the capacitor
ratio alone. The maximally-flat Butterworth response, the right choice
for an anti-alias filter because it has no pass-band ripple, needs $Q =
1/\sqrt{2}$, hence $C_1/C_2 = 2$. We take $C_1 = 22\,\si{\nano\farad}$
and $C_2 = 11\,\si{\nano\farad}$, then solve $\omega_0 = 2\pi\cdot
1000 = 1/(R\sqrt{C_1 C_2})$ for the resistor:
\[
R = \frac{1}{\omega_0\sqrt{C_1 C_2}}
= \frac{1}{2\pi\cdot 1000\cdot\sqrt{22\times 10^{-9}\cdot 11\times 10^{-9}}}
= 10.2\,\si{\kilo\ohm},
\]
which we round to a standard $10\,\si{\kilo\ohm}$ part, nudging the corner
to $f_c = 1/(2\pi\cdot 10\times 10^{3}\sqrt{22\times 10^{-9}\cdot 11
\times 10^{-9}}) = 1.03\,\si{\kilo\hertz}$, inside tolerance.

The second-order Butterworth magnitude is $|H| = 1/\sqrt{1+(f/f_c)^4}$.
At Nyquist, $f/f_c = 4000/1030 = 3.88$, so $|H| = 1/\sqrt{1 + 3.88^4} =
1/15.1 = 0.066$, that is $-23.6\,\si{\decibel}$, and at the first alias
frequency $f_s = 8\,\si{\kilo\hertz}$, $f/f_c = 7.77$ gives $-36\,\si{
\decibel}$. Figure~\ref{fig:vol04ch09:antialias-bode} plots the response
with these marks. Two poles do not reach the $-96\,\si{\decibel}$ a
full-scale worst-case interferer would demand; the design closes the gap
by noting that the physical signal has no energy near Nyquist in the
first place, so the worst credible out-of-band tone is far below full
scale, and the $-24\,\si{\decibel}$ the two-pole filter delivers brings
even a half-scale interferer to within a few least significant bits. An
application that cannot make that argument moves to a fourth-order filter
(two cascaded Sallen-Key stages) or oversamples and filters in the
digital domain.

% Figure: unity-gain Sallen-Key second-order low-pass filter, drawn with plain
% TikZ (resistor and capacitor symbols hand-drawn, op-amp as a triangle). No
% circuitikz dependency.
\begin{figure}[htbp]
\centering
\begin{tikzpicture}[
  line width=0.8pt,
  res/.style={draw, engblue, fill=white, minimum width=10mm, minimum height=4mm, font=\scriptsize},
  font=\small,
]
  % input node
  \node[font=\scriptsize, anchor=east] at (-0.3,2.0) {$v_{\text{in}}$};
  \draw[engblue] (-0.3,2.0) -- (0.2,2.0);
  % R1
  \node[res] (R1) at (0.9,2.0) {$R_1$};
  \draw[engblue] (R1) -- (1.6,2.0);
  % node A (between R1 and R2)
  \node[circle, fill=engblack, inner sep=0.7pt] (A) at (1.6,2.0) {};
  \node[font=\tiny, anchor=south] at (1.6,2.1) {A};
  % R2
  \node[res] (R2) at (2.3,2.0) {$R_2$};
  \draw[engblue] (1.6,2.0) -- (R2) -- (3.0,2.0);
  % node B (op-amp non-inverting input)
  \node[circle, fill=engblack, inner sep=0.7pt] (B) at (3.0,2.0) {};
  \node[font=\tiny, anchor=south] at (3.0,2.1) {B};
  % C2 from node B to ground
  \draw[engblue] (3.0,2.0) -- (3.0,1.0);
  \draw[engblue, line width=1.1pt] (2.75,1.0) -- (3.25,1.0);
  \draw[engblue, line width=1.1pt] (2.75,0.85) -- (3.25,0.85);
  \node[font=\scriptsize, anchor=west] at (3.35,0.93) {$C_2$};
  \draw[engblue] (3.0,0.85) -- (3.0,0.5);
  \node[font=\tiny] at (3.0,0.3) {gnd};
  % op-amp triangle
  \draw[engblue] (3.4,1.4) -- (3.4,3.2) -- (5.2,2.3) -- cycle;
  \node[font=\scriptsize, anchor=west] at (3.5,2.8) {$-$};
  \node[font=\scriptsize, anchor=west] at (3.5,1.8) {$+$};
  % connect node B to the + input
  \draw[engblue] (3.0,2.0) -- (3.4,2.0);
  % output
  \draw[engblue] (5.2,2.3) -- (6.2,2.3);
  \node[font=\scriptsize, anchor=west] at (6.25,2.3) {$v_{\text{out}}$};
  % feedback wire from output to inverting input (unity gain)
  \draw[engblue] (5.7,2.3) -- (5.7,3.6) -- (3.2,3.6) -- (3.2,2.6) -- (3.4,2.6);
  % C1 from node A to the output node (positive feedback)
  \draw[engblue] (1.6,2.0) -- (1.6,3.0);
  \draw[engblue, line width=1.1pt] (1.35,3.0) -- (1.85,3.0);
  \draw[engblue, line width=1.1pt] (1.35,3.15) -- (1.85,3.15);
  \node[font=\scriptsize, anchor=east] at (1.3,3.07) {$C_1$};
  \draw[engblue] (1.6,3.15) -- (1.6,4.0) -- (5.7,4.0) -- (5.7,3.6);
  \node[font=\small] at (3.0,-0.1) {unity-gain Sallen--Key low-pass};
\end{tikzpicture}
\caption[Unity-gain Sallen--Key low-pass filter]{The unity-gain Sallen--Key
second-order low-pass filter. The op-amp is wired as a follower, so the
output equals the voltage at node B. The series resistors $R_1$ and $R_2$
roll the signal off, $C_2$ shunts node B to ground, and $C_1$ feeds the
output back to node A, the positive-feedback path that shapes the response
near the corner. The transfer function is second order, with a corner
frequency set by $R_1 R_2 C_1 C_2$ and a damping set by the resistor ratio
and the capacitor ratio, derived in the case study.}
\label{fig:vol04ch09:sallen-key}
\end{figure}


% Figure: magnitude response of the second-order Butterworth Sallen-Key
% anti-alias filter (fc = 1 kHz, Q = 1/sqrt(2)) used in the signal-chain case
% study, with the converter Nyquist frequency (4 kHz) marked and the
% attenuation there read off. pgfplots semilog magnitude.
\begin{figure}[htbp]
\centering
\begin{tikzpicture}
\begin{semilogxaxis}[
  width=11.5cm, height=5.2cm,
  xlabel={frequency $f$ (Hz)},
  ylabel={$|H|$ (dB)},
  xmin=10, xmax=100000, ymin=-80, ymax=10,
  axis lines=left,
  tick label style={font=\tiny},
  label style={font=\scriptsize},
  ytick={0,-20,-40,-60},
  clip=false,
]
% second-order Butterworth: |H| = 1/sqrt(1 + (f/fc)^4), fc = 1000 Hz
% 20 log10 |H| = -10 log10(1 + (f/1000)^4)
\addplot[engblue, very thick, domain=10:100000, samples=400]
  {-10*log10(1 + (x/1000)^4)};
% -40 dB/dec asymptote past the corner
\addplot[engred, dashed, domain=1000:100000, samples=2] {-40*log10(x/1000)};
% corner marker
\addplot[enggray, dotted] coordinates {(1000,-80) (1000,10)};
\node[engorange, font=\tiny, anchor=south west] at (axis cs:1000,-12)
  {$f_c=1\,$kHz};
% Nyquist marker at fs/2 = 4 kHz
\addplot[enggray, dotted] coordinates {(4000,-80) (4000,10)};
\node[engamber, font=\tiny, anchor=south west, align=left] at (axis cs:4000,-46)
  {$f_N=4\,$kHz\\$-48\,$dB};
\end{semilogxaxis}
\end{tikzpicture}
\caption[Anti-alias filter response]{Magnitude of the second-order
Butterworth Sallen--Key low-pass filter of the signal-chain case study,
with corner $f_c = 1\,\si{\kilo\hertz}$ and damping $\zeta = 1/\sqrt{2}$.
The response is flat in band, breaks at the corner, and falls at
$-40\,\si{\decibel}$ per decade beyond it (dashed asymptote). At the
converter's Nyquist frequency $f_N = 4\,\si{\kilo\hertz}$ the attenuation is
about $48\,\si{\decibel}$, the figure the alias-rejection budget must clear.
A single first-order stage would reach only half that slope and would let
aliased content fold back into the band.}
\label{fig:vol04ch09:antialias-bode}
\end{figure}


\subsection{Noise and offset budget}

The chain's error budget is referred to the input, the $2.5\,\si{\milli
\volt}$ full-scale bridge signal, and compared against the input-referred
least significant bit, which is $\text{LSB}/G = 30.5\,\si{\micro\volt}
/400 = 76\,\si{\nano\volt}$. Three error sources matter.

The instrumentation amplifier's input offset voltage, a few tens of
microvolts for a precision part, sits directly in series with the bridge
signal and is amplified by the full gain. A $25\,\si{\micro\volt}$ offset
is a third of the $76\,\si{\nano\volt}$ budget multiplied up, so it
dominates unless removed; the cure is a single DC measurement at zero
strain subtracted in software, which removes the static offset and leaves
only its drift with temperature, of order $0.1\,\si{\micro\volt}$ per
kelvin, a few tenths of a least significant bit over a ten-kelvin
swing.

The amplifier's input-referred voltage noise, of order $10\,\si{\nano
\volt}/\sqrt{\si{\hertz}}$ for a low-noise part, integrates over the
chain's noise bandwidth. The anti-alias filter sets that bandwidth: for a
second-order Butterworth the noise bandwidth is $1.11\,f_c =
1.14\,\si{\kilo\hertz}$. The integrated noise is then $10\,\si{\nano
\volt}/\sqrt{\si{\hertz}}\cdot\sqrt{1140\,\si{\hertz}} = 0.34\,\si{\micro
\volt}$ RMS referred to the input, about four and a half input-referred
least significant bits RMS. To bury the noise under one bit the design
either narrows the bandwidth (the signal allows it, since real content
stops at $1\,\si{\kilo\hertz}$, but the anti-alias corner is already
there) or averages in software: averaging $N$ samples drops the noise by
$\sqrt{N}$, so averaging $25$ samples brings the $0.34\,\si{\micro\volt}$
down to $68\,\si{\nano\volt}$, just under the bit, at the cost of slowing
the effective output rate to $8000/25 = 320\,\si{\hertz}$. That output
rate still resolves the signal, whose content stops at $1\,\si{\kilo
\hertz}$ only because the anti-alias filter put it there, so the working
compromise lowers the anti-alias corner to match the slower output rate,
trading response speed for the bit of noise headroom averaging buys. The
budget exposes the trade rather than hiding it.

The bridge's own Johnson noise, $\sqrt{4k_B T R\,\Delta f}$ with $R =
350\,\si{\ohm}$ and $\Delta f = 1140\,\si{\hertz}$, is $\sqrt{4\cdot
1.38\times 10^{-23}\cdot 300\cdot 350\cdot 1140} = 0.083\,\si{\micro
\volt}$ RMS, about one input-referred least significant bit, the floor no
amplifier can beat. The amplifier noise sits above it, so the front end
is amplifier-limited, and the design lever is the amplifier choice, not
the bridge resistance.

\subsection{Verdict}

The chain delivers the full-scale strain as a full-scale converter
reading, rejects the common-mode excitation through the instrumentation
amplifier, holds the static offset to a software-removable level, and
attenuates a Nyquist-frequency interferer by $24\,\si{\decibel}$, enough
for a signal that carries no real energy near Nyquist. The residual error
is dominated by amplifier voltage noise at about four least significant
bits RMS, which a quieter amplifier or modest averaging reduces, and by
offset drift at a few tenths of a bit. An engineer would sign this design
for an application needing twelve to thirteen effective bits out of the
sixteen the converter offers, and would reach for a chopper-stabilised
amplifier and a fourth-order filter only if the application demanded the
full sixteen. The chain is the information counterpart of the two power
supplies: the same chapter machinery, turned to moving a measurement
instead of moving energy.

\begin{mastery}
The noise budget above treated the amplifier's input-referred voltage
noise as white, but a real op-amp's noise spectral density rises at low
frequency as $1/f$, the \engterm{flicker} or pink-noise region, below a
\engterm{corner frequency} $f_{\text{nc}}$ where the $1/f$ density crosses
the flat broadband floor. The total input-referred noise integrates the
density over the band, $\overline{v_n^2} = \int_{f_1}^{f_2} S_v(f)\,\dd f$
with $S_v(f) = S_0(1 + f_{\text{nc}}/f)$, giving $\overline{v_n^2} =
S_0[(f_2 - f_1) + f_{\text{nc}}\ln(f_2/f_1)]$. The flat term grows with
bandwidth; the $1/f$ term grows only logarithmically, so it matters most
for slow measurements that integrate down to fractions of a hertz, which
is the regime a strain gauge reading a static load sits in. A corner of
$10\,\si{\hertz}$ and a measurement extending from $0.1\,\si{\hertz}$ to
$1140\,\si{\hertz}$ contributes a $1/f$ term of $f_{\text{nc}}\ln(1140/0.1)
= 10\cdot 9.3 = 93\,\si{\hertz}$ of equivalent noise bandwidth, comparable
to the flat $1140\,\si{\hertz}$ over a wide band but dominant for a
measurement confined below a hertz. The architectural answer for
near-DC measurement is the chopper-stabilised or auto-zero amplifier,
which modulates the signal above $f_{\text{nc}}$, amplifies it where the
noise is white, and demodulates back to baseband, moving the
$1/f$ corner toward zero and leaving only the flat floor. The same
modulate-amplify-demodulate idea is the lock-in amplifier of
instrumentation practice, and it is why a precision bridge front end for
static strain is built around an auto-zero part rather than the lowest
broadband-noise op-amp the catalogue offers.
\end{mastery}

\section{An applied vignette: the 555-timer astable oscillator}\pathtag{standard}
\label{sec:vol04ch09:555}

The case studies above each carry a full specification end to end; this
shorter vignette takes one circuit that every engineer meets early, the
555-timer astable oscillator, and shows how the first-order RC transient
of Section~\ref{sec:vol04ch09:transient} sets its period directly. The 555
is an integrated timer whose internal comparators watch the voltage on an
external capacitor against two fixed thresholds, $\tfrac{1}{3}V_{CC}$ and
$\tfrac{2}{3}V_{CC}$, set by an internal resistor divider, and switch an
internal transistor and the output accordingly.
Figure~\ref{fig:vol04ch09:555-astable} draws the timing network and the
waveforms it produces.

% Figure: the 555-timer astable oscillator schematic (timing network R_A,
% R_B, C) on the left and its capacitor and output waveforms on the right.
% Plain TikZ for the schematic; pgfplots for the waveforms.
\begin{figure}[htbp]
\centering
% === left: timing network schematic ===
\begin{tikzpicture}[
  line width=0.8pt,
  res/.style={draw, engblue, fill=white, minimum width=4mm, minimum height=9mm, font=\scriptsize},
  font=\small,
]
  % 555 block
  \draw[engblue] (2.4,0.0) rectangle (4.4,3.0);
  \node[font=\scriptsize] at (3.4,2.6) {555};
  \node[font=\tiny, anchor=east] at (2.4,2.4) {DIS};
  \node[font=\tiny, anchor=east] at (2.4,1.5) {THR};
  \node[font=\tiny, anchor=east] at (2.4,0.6) {TRG};
  \node[font=\tiny, anchor=west] at (4.4,1.5) {OUT};
  % Vcc rail
  \draw[engred] (0.0,4.0) -- (5.0,4.0);
  \node[engred, font=\scriptsize, anchor=south] at (2.5,4.0) {$V_{CC}$};
  \draw[engred] (3.4,3.0) -- (3.4,4.0);
  % R_A from Vcc to DIS
  \node[res] (RA) at (0.6,3.2) {$R_A$};
  \draw[engblue] (0.6,4.0) -- (RA);
  \draw[engblue] (0.6,2.75) -- (0.6,2.4) -- (2.4,2.4);
  % R_B from DIS to THR/TRG node
  \node[res] (RB) at (0.6,1.6) {$R_B$};
  \draw[engblue] (0.6,2.75) -- (RB);
  \draw[engblue] (0.6,1.15) -- (0.6,0.6);
  % node joining THR and TRG
  \draw[engblue] (0.6,1.5) -- (2.4,1.5);
  \draw[engblue] (0.6,0.6) -- (2.4,0.6);
  \node[circle, fill=engblack, inner sep=0.7pt] at (0.6,1.5) {};
  % capacitor C from THR/TRG node to ground
  \draw[engblue] (1.5,1.5) -- (1.5,0.95);
  \node[circle, fill=engblack, inner sep=0.7pt] at (1.5,1.5) {};
  \draw[engblue, line width=1.1pt] (1.25,0.95) -- (1.75,0.95);
  \draw[engblue, line width=1.1pt] (1.25,0.80) -- (1.75,0.80);
  \node[font=\scriptsize, anchor=west] at (1.8,0.87) {$C$};
  \draw[engblue] (1.5,0.80) -- (1.5,0.3) -- (3.4,0.3);
  \node[font=\tiny] at (3.4,0.05) {gnd};
  % output
  \draw[engblue] (4.4,1.5) -- (5.0,1.5);
  \node[font=\scriptsize, anchor=west] at (5.0,1.5) {$v_o$};
\end{tikzpicture}
\hspace{0.3cm}
% === right: waveforms ===
\begin{tikzpicture}
\begin{axis}[
  width=6.6cm, height=5.0cm,
  xlabel={$t/T$}, ylabel={voltage},
  xmin=0, xmax=2, ymin=-0.1, ymax=1.25,
  axis lines=left,
  tick label style={font=\tiny},
  label style={font=\scriptsize},
  ytick={0.333,0.667,1.0},
  yticklabels={$\tfrac{1}{3}V_{CC}$,$\tfrac{2}{3}V_{CC}$,$V_{CC}$},
  yticklabel style={font=\tiny},
  legend style={font=\tiny, at={(0.98,0.30)}, anchor=east, draw=none, fill=none},
  clip=false,
]
% capacitor voltage: charges 1/3 -> 2/3 over fraction 0.64 of period, then
% discharges 2/3 -> 1/3 over 0.36 of period. Sketch as piecewise exponentials.
% Charge phase 0 .. 0.64 : 0.333 + (1 - 0.333) approached; use shaped curve to 0.667
\addplot[engblue, very thick, domain=0:0.64, samples=80]
  {1 - (1 - 0.333)*exp(-1.72*x/0.64)};
\addplot[engblue, very thick, domain=0.64:1.0, samples=60]
  {0.667*exp(-1.10*(x-0.64)/0.36)};
\addplot[engblue, very thick, domain=1.0:1.64, samples=80]
  {1 - (1 - 0.333)*exp(-1.72*(x-1.0)/0.64)};
\addplot[engblue, very thick, domain=1.64:2.0, samples=60]
  {0.667*exp(-1.10*(x-1.64)/0.36)};
\addlegendentry{$v_C$}
% output: high during charge (0 .. 0.64), low during discharge
\addplot[engorange, thick] coordinates {
  (0,1.1) (0.64,1.1) (0.64,-0.05) (1.0,-0.05) (1.0,1.1) (1.64,1.1)
  (1.64,-0.05) (2.0,-0.05)};
\addlegendentry{$v_o$}
\draw[enggray, dashed] (axis cs:0,0.333) -- (axis cs:2,0.333);
\draw[enggray, dashed] (axis cs:0,0.667) -- (axis cs:2,0.667);
\end{axis}
\end{tikzpicture}
\caption[555-timer astable oscillator and its waveforms]{The 555-timer
astable oscillator (left) and its waveforms (right). The capacitor
charges toward $V_{CC}$ through $R_A+R_B$ until it reaches the upper
threshold $\tfrac{2}{3}V_{CC}$, at which point the discharge transistor
turns on and the capacitor discharges through $R_B$ alone toward the
lower threshold $\tfrac{1}{3}V_{CC}$, where the cycle restarts. The
output is high during the charge phase and low during the discharge
phase, so the asymmetry between the charge path ($R_A+R_B$) and the
discharge path ($R_B$) sets a duty cycle always above one half.}
\label{fig:vol04ch09:555-astable}
\end{figure}


In the astable connection the capacitor $C$ charges from the supply
through the series pair $R_A + R_B$ until its voltage reaches the upper
threshold $\tfrac{2}{3}V_{CC}$, at which point the internal transistor
turns on and pulls the junction of $R_A$ and $R_B$ to ground, so the
capacitor now discharges through $R_B$ alone until it falls to the lower
threshold $\tfrac{1}{3}V_{CC}$, where the transistor turns off and the
charge phase begins again. Each phase is a first-order RC transient
between two fixed voltages, and the time for each is read straight from
the charging law $v_C(t) = v_\infty + (v_0 - v_\infty)e^{-t/\tau}$. The
charge phase runs from $\tfrac{1}{3}V_{CC}$ toward $V_{CC}$ with time
constant $\tau_1 = (R_A + R_B)C$, and solving for the time to reach
$\tfrac{2}{3}V_{CC}$ gives
\[
t_{\text{high}} = (R_A + R_B)C\,\ln\!\frac{V_{CC} - \tfrac{1}{3}V_{CC}}
{V_{CC} - \tfrac{2}{3}V_{CC}}
= (R_A + R_B)C\,\ln 2 = 0.693\,(R_A + R_B)C,
\]
the supply voltage cancelling out of the ratio so the timing is set by the
resistors and capacitor alone, not by $V_{CC}$. The discharge phase runs
from $\tfrac{2}{3}V_{CC}$ toward zero through $R_B$ with time constant
$\tau_2 = R_B C$, and the time to fall to $\tfrac{1}{3}V_{CC}$ is, by the
same ratio of thresholds,
\[
t_{\text{low}} = R_B C\,\ln\!\frac{\tfrac{2}{3}V_{CC}}{\tfrac{1}{3}V_{CC}}
= R_B C\,\ln 2 = 0.693\,R_B C.
\]
The total period is $T = t_{\text{high}} + t_{\text{low}} = 0.693\,(R_A +
2R_B)C$, so the oscillation frequency is
\[
f = \frac{1}{T} = \frac{1.44}{(R_A + 2R_B)\,C}.
\]
Because the charge path includes $R_A + R_B$ and the discharge path only
$R_B$, the high time always exceeds the low time and the duty cycle
$t_{\text{high}}/T = (R_A + R_B)/(R_A + 2R_B)$ is always above one half; a
duty cycle below one half needs a diode across $R_B$ to route the charge
current around it, or a different topology.

A worked instance fixes the numbers. For a $1\,\si{\kilo\hertz}$ square
wave at roughly equal duty, choose $R_A$ small against $R_B$: with $R_A =
1\,\si{\kilo\ohm}$, $R_B = 68\,\si{\kilo\ohm}$, and $C = 10\,\si{\nano
\farad}$, the period is $T = 0.693\,(1 + 136)\times 10^{3}\cdot 10\times
10^{-9} = 0.95\,\si{\milli\second}$, so $f = 1.05\,\si{\kilo\hertz}$, and
the duty cycle is $(1 + 68)/(1 + 136) = 50.4\,\%$, close to symmetric
because $R_A \ll R_B$. The design lever is the capacitor for coarse
frequency and the resistor pair for fine frequency and duty, and the
timing's independence of $V_{CC}$ is what makes the 555 a stable clock
across a supply that sags, the property that kept the part in catalogues
for half a century. The same first-order transient that charges the
reservoir capacitor of the bench supply, run between two comparator
thresholds instead of toward a single final value, is the whole of the
oscillator's timing.

\section{Power: real, reactive, apparent}\pathtag{standard}

In sinusoidal steady state the instantaneous power $p(t) = v(t)i(t)$
oscillates at $2\omega$. Its time average over a cycle is the
\engterm{real power} or \engterm{average power}:
\[
P = \frac{1}{T}\int_{0}^{T} v(t)i(t)\,\dd t = \frac{1}{2}V_{m}I_{m}
\cos\phi = V_{\text{rms}}I_{\text{rms}}\cos\phi,
\]
where $\phi$ is the phase angle by which the voltage leads the
current. The factor $\cos\phi$ is the \engterm{power factor}, the
ratio of real power to the product of RMS voltage and RMS current
(the \engterm{apparent power} $S = V_{\text{rms}}I_{\text{rms}}$ in
volt-amperes). The \engterm{reactive power} is $Q = V_{\text{rms}}I_{
\text{rms}}\sin\phi$ in volt-amperes-reactive (VAR); it represents
energy oscillating between source and load without net transfer.

% Figure: the power triangle, real P horizontal, reactive Q vertical, apparent
% S as the hypotenuse, with the power-factor angle phi. Plain TikZ.
\begin{figure}[htbp]
\centering
\begin{tikzpicture}[
  line width=1.0pt,
  font=\small,
  >={Stealth[length=2.5mm]},
]
\coordinate (O) at (0,0);
\coordinate (P) at (5,0);
\coordinate (S) at (5,3);
% real power along base
\draw[engblue, ->] (O) -- (P);
\node[engblue, anchor=north] at (2.5,-0.1) {$P$ (W), real power};
% reactive power vertical
\draw[engred, ->] (P) -- (S);
\node[engred, anchor=west] at (5.1,1.5) {$Q$ (VAR), reactive power};
% apparent power hypotenuse
\draw[engorange, ->] (O) -- (S);
\node[engorange, anchor=south east] at (2.5,1.55) {$S$ (VA), apparent power};
% angle arc
\draw[engblack] (1.3,0) arc (0:31:1.3);
\node[font=\scriptsize] at (1.7,0.35) {$\phi$};
% right-angle mark at P
\draw[engblack] (4.7,0) -- (4.7,0.3) -- (5,0.3);
% relation note
\node[font=\scriptsize, anchor=north west] at (0,-0.9)
  {$S = \sqrt{P^2 + Q^2}, \qquad \text{power factor} = \cos\phi = P/S$};
\end{tikzpicture}
\caption[The power triangle]{The power triangle of AC steady state. Real
power $P$ does the useful work; reactive power $Q$ shuttles energy back and
forth between source and load without net transfer; apparent power $S$ is the
product of RMS voltage and RMS current that the distribution network must
carry. The power factor $\cos\phi = P/S$ measures how much of the apparent
power becomes real work. Power-factor correction shrinks $Q$ toward zero,
rotating $S$ down onto $P$ and reducing the line current.}
\label{fig:vol04ch09:power-triangle}
\end{figure}


The three powers form a right triangle: $S^{2} = P^{2}+Q^{2}$, or
$S = P + jQ$ as a complex number (the \engterm{complex power}), drawn
in Figure~\ref{fig:vol04ch09:power-triangle}.
Inductive loads (motors, transformers) absorb reactive power
($Q>0$); capacitive loads supply reactive power ($Q<0$).
\engterm{Power factor correction}, the standard utility-grid practice
of installing capacitor banks at industrial loads, brings the overall
power factor near unity by cancelling the inductive reactive power of
the motors, reducing the apparent current the distribution network
must carry.

A motor rated $10\,\si{\kilo\watt}$ at power factor $0.8$ lagging
draws $S = 12.5\,\si{\kilo\volt\ampere}$, $Q = 7.5\,\text{kvar}$.
Adding a $7.5\,\text{kvar}$ capacitor bank in parallel cancels $Q$,
reducing the apparent power
to $10\,\si{\kilo\volt\ampere}$ and the line current by $20\,\%$. The
capacitor pays for itself in reduced distribution losses and demand
charges within a typical industrial-billing cycle.

\begin{mastery}
Complex power has a structural property worth carrying: Tellegen's
theorem states that for any network satisfying KCL and KVL, the sum of
the complex powers delivered to all branches is zero, regardless of the
elements' constitutive laws. Conservation of complex power follows as a
corollary, separating into conservation of real power $P$ and
conservation of reactive power $Q$. The reactive-power balance is the
formal basis of power-factor correction: a capacitor supplies the
reactive power an inductor demands, so installing the capacitor next to
the inductive load lets the reactive power circulate locally instead of
loading the whole distribution path. The same balance underlies the
reactive-power dispatch that grid operators run continuously to hold bus
voltages within tolerance.
\end{mastery}

\section{Three-phase introduction}\pathtag{standard}

A \engterm{three-phase} power system uses three sinusoidal voltages
of equal amplitude offset by $120^{\circ}$ in phase,
\[
v_{a}(t) = V_{m}\cos(\omega t), \quad v_{b}(t) = V_{m}\cos(\omega t -
120^{\circ}), \quad v_{c}(t) = V_{m}\cos(\omega t + 120^{\circ}).
\]
The three voltages sum to zero at every instant, $v_{a}+v_{b}+v_{c}=0$,
a consequence of the trigonometric identity for three equally-spaced
sinusoids. The three-phase system has three operational advantages
over single-phase: constant instantaneous power delivery to a balanced
three-phase load (rather than $2\omega$-oscillating power as in
single-phase); use of three rather than six conductors for the same
power; and the natural production of a rotating magnetic field in
three-phase motors, which simplifies motor design.

% Figure: three-phase waveforms (left) summing to zero, and the wye/delta
% phasor stars (right). pgfplots for waveforms plus a small TikZ inset.
\begin{figure}[htbp]
\centering
\begin{tikzpicture}
\begin{axis}[
  width=8.5cm, height=6cm,
  xlabel={phase angle $\omega t$ (deg)}, ylabel={voltage / $V_m$},
  xmin=0, xmax=360, ymin=-1.25, ymax=1.25,
  xtick={0,90,180,270,360},
  axis lines=left,
  tick label style={font=\scriptsize},
  label style={font=\small},
  legend pos=outer north east,
  legend style={font=\scriptsize},
  clip=false,
]
\addplot[engblue, very thick, domain=0:360, samples=120] {cos(x)};
\addlegendentry{$v_a$}
\addplot[engred, very thick, domain=0:360, samples=120] {cos(x-120)};
\addlegendentry{$v_b$}
\addplot[engorange, very thick, domain=0:360, samples=120] {cos(x+120)};
\addlegendentry{$v_c$}
\addplot[enggray, dashed] coordinates {(0,0) (360,0)};
\end{axis}
\end{tikzpicture}
\caption[Three-phase voltages]{The three line voltages of a balanced
three-phase system, equal in amplitude and offset by $120^\circ$. At every
instant the three sum to zero, $v_a + v_b + v_c = 0$, which is why a balanced
wye load needs no neutral return current. The constant total power delivered
to a balanced three-phase load (rather than the $2\omega$-pulsing power of a
single phase) is the reason rotating machinery and bulk distribution run on
three phases.}
\label{fig:vol04ch09:three-phase}
\end{figure}


Figure~\ref{fig:vol04ch09:three-phase} plots the three voltages and
shows their instantaneous sum vanishing.

Two connection topologies are standard:
\begin{itemize}[leftmargin=*]
  \item \engterm{Wye} (or star, $Y$): the three sources share a common
    neutral. The \engterm{line voltage} (between two phases) is
    $\sqrt{3}$ times the \engterm{phase voltage} (between one phase
    and neutral). North American residential service is wye:
    $120\,\si{\volt}$ phase-to-neutral, $208\,\si{\volt}$ phase-to-
    phase in a three-phase building.
  \item \engterm{Delta} ($\Delta$): the three sources form a closed
    triangle; no neutral. Line and phase voltages are equal; line and
    phase currents differ by $\sqrt{3}$. Industrial motor wiring is
    often delta.
\end{itemize}

A balanced three-phase load consumes total real power $P =
\sqrt{3}V_{L}I_{L}\cos\phi$ where $V_{L}$ and $I_{L}$ are line voltage
and line current (RMS) and $\phi$ is the power factor angle. The
formula is identical for wye and delta loads. The $\sqrt{3}$ factor is
the working multiplier the practising engineer carries when estimating
three-phase loads from nameplate data.

\section{Failure: grounding, EMC, signal integrity}\pathtag{core}

Three failure modes recur in lumped-circuit practice. They are the
counterpart in this chapter to the EMI and ground-loop failures of
Chapter 8; here we treat them at the circuit-design level rather than
the field level.

\engterm{Improper grounding} is the most common failure mode of
electronics that fail in the field after passing bench tests. The
common pattern is that the designer assumes a single ideal ground
node in the schematic, but in the physical layout the ground is a
finite-impedance return path with measurable inductance and
resistance. A high-current pulse along the ground trace drops a
voltage that appears as common-mode interference on every other
signal that referenced the same ground point. The cure is a
\engterm{ground plane}: a large uninterrupted copper area on the
printed circuit board whose low impedance keeps the ground reference
nearly uniform across the board. For mixed-signal designs (digital
and sensitive analog on the same board), separate analog and digital
grounds joined at a single point near the data-converter chip prevent
digital switching noise from contaminating the analog reference.

\engterm{Power-rail decoupling} is the second-most-common failure
mode. Every active device on a board draws current pulses from its
supply when it switches; the pulses traverse the impedance of the
supply distribution and induce voltage transients on the rail. The
cure is decoupling capacitors: small ceramic capacitors (typically
$100\,\si{\nano\farad}$) placed directly at each chip's supply pin,
providing the high-frequency charge that the slower bulk supply
cannot. The bulk supply has larger electrolytic capacitors
($10$ to $100\,\si{\micro\farad}$) for the lower-frequency
disturbances. The two-tier decoupling scheme covers from $1\,\si{
\kilo\hertz}$ to $100\,\si{\mega\hertz}$ at the chip; above $100\,\si{
\mega\hertz}$, ESL of the ceramic capacitors begins to dominate and
the working approach shifts to embedded-capacitance layer stacks.

\engterm{Signal integrity} failures arise when the lumped-circuit
approximation breaks down. A digital signal with a $1\,\si{\nano
\second}$ rise time has spectral content into $300\,\si{\mega\hertz}$
and beyond; at that frequency, a $30\,\si{\centi\meter}$ trace is a
significant fraction of a wavelength, and the trace must be treated
as a transmission line. Reflections from impedance mismatches distort
the signal; the receiver may interpret the distortion as spurious
transitions. The cure is impedance-matched transmission lines,
terminated at the receiver (or the source) with their characteristic
impedance to absorb incident power without reflection. Chapter 10
develops the transmission-line theory in detail; in this chapter the
principle is that above a few tens of megahertz, every trace is a
transmission line whether the designer planned for it or not.

A representative engineering standard is IEC 61000-4-2 (electrostatic
discharge immunity) and IEC 61000-4-4 (electrical fast-transient/
burst immunity), which specify the test waveforms a piece of
electrical equipment must withstand to be EMC-compliant on the
European market. The standards are not academic abstractions; a
device that fails the test cannot legally be sold, and the test
itself is the working specification for the input-protection circuit
on every commercial interface.

\section{Project}\pathtag{core}

\begin{project}[Bandpass filter design and characterisation]
Design, build, and characterise a passive RLC bandpass filter with
specified centre frequency and bandwidth, and compare measured
performance to theory.

\textbf{Specifications.} Centre frequency $f_{0} = 10\,\si{\kilo
\hertz}$, quality factor $Q = 5$ (bandwidth $\Delta f = 2\,\si{\kilo
\hertz}$), input and output impedance $\approx 1\,\si{\kilo\ohm}$.

\textbf{Design.} A series RLC tuned to $f_{0} = 10\,\si{\kilo\hertz}$
has $\omega_{0} = 2\pi\times 10^{4} = 6.28\times 10^{4}\,\si{\per
\second}$. With $L = 10\,\si{\milli\henry}$ (standard preferred
value), $C = 1/(\omega_{0}^{2}L) = 25\,\si{\nano\farad}$ (use $24$ or
$27\,\si{\nano\farad}$ standard value, retune iteratively); for
$Q = 5$, $R = \omega_{0}L/Q = 126\,\si{\ohm}$.

\textbf{Build.} Solder the components on a piece of strip-board or
ground-plane PCB. Mount BNC or pin-header connectors for input and
output. Photograph the build.

\textbf{Characterise.} Drive the input with a function generator at
swept frequencies from $1$ to $100\,\si{\kilo\hertz}$ and measure the
output amplitude and phase with an oscilloscope. Plot the magnitude
response on log-log axes and the phase response on semi-log axes. Fit
the data to the theoretical magnitude response and extract the
measured centre frequency, bandwidth, and $Q$. Report the relative
error of each parameter.

\textbf{Report.} One-page report with the schematic, the build
photograph, the measured frequency response plotted alongside theory,
and a discussion of the dominant error sources. Expected: $f_{0}$
within $5\,\%$ of design (limited by capacitor tolerance); $Q$ within
$15\,\%$ of design (limited by inductor ESR being larger than the
designed $R$ at audio frequencies).
\end{project}

\section{Exercises}

\subsection*{Calculation}\pathtag{core}

\begin{exercise}
A series circuit has a $12\,\si{\volt}$ source, a $100\,\si{\ohm}$
resistor, and a $50\,\si{\ohm}$ resistor. Compute the current, the
voltage across each resistor, and the power dissipated by each.
\paragraph{Solution sketch.}
The series resistance is $150\,\si{\ohm}$, so the current is
$I = 12/150 = 80\,\si{\milli\ampere}$. The voltages divide in
proportion to resistance: $V_{100} = 0.08\cdot 100 = 8\,\si{\volt}$,
$V_{50} = 0.08\cdot 50 = 4\,\si{\volt}$, summing to the source as KVL
requires. The dissipated powers are $P_{100} = I^{2}R = 0.0064\cdot
100 = 0.64\,\si{\watt}$ and $P_{50} = 0.32\,\si{\watt}$; the total
$0.96\,\si{\watt}$ equals $VI = 12\cdot 0.08$.
\end{exercise}

\begin{exercise}
Three resistors $R_{1} = 100\,\si{\ohm}$, $R_{2} = 200\,\si{\ohm}$,
$R_{3} = 300\,\si{\ohm}$ are in parallel across a $24\,\si{\volt}$
source. Compute the equivalent resistance, the total current, the
current through each resistor, and the total power dissipated.
\paragraph{Solution sketch.}
The parallel conductance is $1/100 + 1/200 + 1/300 = 0.01833\,\si{
\siemens}$, so $R_{\text{eq}} = 54.5\,\si{\ohm}$. The total current is
$24/54.5 = 0.440\,\si{\ampere}$. Each branch sees the full
$24\,\si{\volt}$: $I_{1} = 0.24\,\si{\ampere}$, $I_{2} = 0.12\,\si{
\ampere}$, $I_{3} = 0.08\,\si{\ampere}$, summing to $0.44\,\si{\ampere}$.
The total power is $24\cdot 0.44 = 10.6\,\si{\watt}$.
\end{exercise}

\begin{exercise}
A circuit consists of a $20\,\si{\volt}$ source, a $5\,\si{\kilo\ohm}$
resistor in series with the parallel combination of a $2\,\si{\kilo
\ohm}$ resistor and a $3\,\si{\kilo\ohm}$ resistor. Compute the
voltage across the parallel combination and the current through each
resistor.
\paragraph{Solution sketch.}
The parallel pair is $2\,\text{k}\parallel 3\,\text{k} = 6/5 =
1.2\,\si{\kilo\ohm}$, in series with $5\,\si{\kilo\ohm}$ for a total of
$6.2\,\si{\kilo\ohm}$. The source current is $20/6200 = 3.23\,\si{
\milli\ampere}$, which flows through the $5\,\si{\kilo\ohm}$ resistor.
The voltage across the parallel combination is $3.23\,\text{mA}\cdot
1.2\,\text{k}\Omega = 3.87\,\si{\volt}$. The branch currents are
$3.87/2000 = 1.94\,\si{\milli\ampere}$ and $3.87/3000 = 1.29\,\si{
\milli\ampere}$, summing to the source current.
\end{exercise}

\begin{exercise}
Find the Th\'evenin equivalent of the network seen looking into
terminals $a$-$b$ for a $10\,\si{\volt}$ source in series with a
$1\,\si{\kilo\ohm}$ resistor that feeds a parallel combination of a
$2\,\si{\kilo\ohm}$ resistor and a $4\,\si{\kilo\ohm}$ resistor; the
$4\,\si{\kilo\ohm}$ resistor's far end is terminal $a$ and the source
ground is terminal $b$.
\paragraph{Solution sketch.}
With $a$-$b$ open no current flows through the $4\,\si{\kilo\ohm}$
resistor, so $V_{\text{Th}}$ is the voltage across the parallel pair
seen through the divider. The $2\,\si{\kilo\ohm}$ resistor and the
$4\,\si{\kilo\ohm}$ resistor in parallel give $R_{p} = 1.33\,\si{\kilo
\ohm}$; the open-circuit voltage at that node is $10\cdot R_{p}/(1\,
\text{k}+R_{p}) = 10\cdot 1.33/2.33 = 5.71\,\si{\volt}$, and since no
current flows in the $4\,\si{\kilo\ohm}$ leg, $V_{\text{Th}} =
5.71\,\si{\volt}$. Deactivating the source (short), the resistance from
$a$ is $4\,\text{k} + (2\,\text{k}\parallel 1\,\text{k}) = 4\,\text{k}+
0.67\,\text{k} = 4.67\,\si{\kilo\ohm}$.
\end{exercise}

\begin{exercise}
An RC circuit has $R = 1\,\si{\kilo\ohm}$, $C = 1\,\si{\micro\farad}$,
driven by a $5\,\si{\volt}$ step from initial zero. Compute the
capacitor voltage at $t = 0.5, 1, 2, 5\,\si{\milli\second}$. Compute
the time at which $v_{C}$ first exceeds $4\,\si{\volt}$.
\paragraph{Solution sketch.}
The time constant is $\tau = RC = 1\,\si{\milli\second}$ and
$v_{C}(t) = 5(1-e^{-t/\tau})$. At $t = 0.5, 1, 2, 5\,\si{\milli
\second}$ the factor $1-e^{-t/\tau}$ is $0.393, 0.632, 0.865, 0.993$,
giving $v_{C} = 1.97, 3.16, 4.32, 4.97\,\si{\volt}$. Setting $v_{C}=4$
gives $e^{-t/\tau} = 0.2$, so $t = \tau\ln 5 = 1.61\,\si{\milli
\second}$.
\end{exercise}

\begin{exercise}
An RLC series circuit has $L = 1\,\si{\milli\henry}$, $C = 100\,\si{
\nano\farad}$, $R = 10\,\si{\ohm}$. Compute the natural frequency, the
damping ratio, the quality factor, and classify the response.
\paragraph{Solution sketch.}
$\omega_{n} = 1/\sqrt{LC} = 1/\sqrt{10^{-3}\cdot 10^{-7}} =
10^{5}\,\si{\per\second}$, so $f_{n} = \omega_{n}/2\pi = 15.9\,\si{\kilo
\hertz}$. The damping ratio is $\zeta = (R/2)\sqrt{C/L} = 5\cdot
\sqrt{10^{-7}/10^{-3}} = 5\cdot 0.01 = 0.05$. The quality factor is
$Q = 1/(2\zeta) = 10$. Since $\zeta < 1$ the response is underdamped
and rings near $\omega_{n}$ before decaying.
\end{exercise}

\begin{exercise}
An RLC bandpass filter has $L = 10\,\si{\milli\henry}$, $C = 100\,\si{
\nano\farad}$, $R = 100\,\si{\ohm}$. Compute the resonant frequency in
hertz, the quality factor, and the $-3\,\si{\decibel}$ bandwidth.
\paragraph{Solution sketch.}
$\omega_{0} = 1/\sqrt{LC} = 1/\sqrt{10^{-2}\cdot 10^{-7}} =
3.16\times 10^{4}\,\si{\per\second}$, so $f_{0} = 5.03\,\si{\kilo
\hertz}$. The quality factor is $Q = \omega_{0}L/R = 3.16\times
10^{4}\cdot 10^{-2}/100 = 3.16$. The $-3\,\si{\decibel}$ bandwidth is
$\Delta f = f_{0}/Q = 5030/3.16 = 1.59\,\si{\kilo\hertz}$.
\end{exercise}

\begin{exercise}
A capacitor $C = 10\,\si{\micro\farad}$ has a voltage that varies as
$v(t) = 5\sin(1000 t)\,\si{\volt}$. Compute the current waveform, the
RMS current, and the impedance magnitude.
\paragraph{Solution sketch.}
$i = C\,\dd v/\dd t = 10^{-5}\cdot 5\cdot 1000\cos(1000t) =
0.05\cos(1000t)\,\si{\ampere}$, leading the voltage by $90^\circ$. The
peak current is $50\,\si{\milli\ampere}$, so the RMS current is
$0.05/\sqrt{2} = 35.4\,\si{\milli\ampere}$. The impedance magnitude is
$|Z_{C}| = 1/(\omega C) = 1/(1000\cdot 10^{-5}) = 100\,\si{\ohm}$,
consistent with the $5\,\si{\volt}$ peak driving a $50\,\si{\milli
\ampere}$ peak.
\end{exercise}

\begin{exercise}
A motor draws $50\,\si{\ampere}$ at $230\,\si{\volt}$ RMS, $50\,\si{
\hertz}$, with a power factor of $0.75$ lagging. Compute the real,
reactive, and apparent power. Compute the capacitance needed in
parallel with the motor to correct the power factor to $0.95$
lagging.
\paragraph{Solution sketch.}
The apparent power is $S = VI = 230\cdot 50 = 11.5\,\si{\kilo\volt
\ampere}$. With $\cos\phi = 0.75$, the real power is $P = 0.75 S =
8.63\,\si{\kilo\watt}$ and the reactive power is $Q_{1} = P\tan\phi_{1}
= 8.63\cdot\tan(41.4^\circ) = 7.61\,\text{kvar}$. To reach
$\cos\phi = 0.95$ ($\phi_{2} = 18.2^\circ$) the load reactive power
becomes $Q_{2} = P\tan\phi_{2} = 2.84\,\text{kvar}$, so the capacitor
supplies $Q_{C} = 4.77\,\text{kvar}$. The capacitance is $C =
Q_{C}/(\omega V^{2}) = 4770/(2\pi\cdot 50\cdot 230^{2}) = 287\,\si{
\micro\farad}$. The companion script \texttt{code/power\_factor.py}
returns the same value.
\end{exercise}

\begin{exercise}
A balanced three-phase load on a $400\,\si{\volt}$ line-to-line system
consumes $50\,\si{\kilo\watt}$ at power factor $0.9$ lagging. Compute
the line current.
\paragraph{Solution sketch.}
From $P = \sqrt{3}V_{L}I_{L}\cos\phi$, solve for the line current:
$I_{L} = P/(\sqrt{3}V_{L}\cos\phi) = 50000/(\sqrt{3}\cdot 400\cdot 0.9)
= 50000/623.5 = 80.2\,\si{\ampere}$. The $\sqrt{3}$ factor is the
working multiplier between single-phase and balanced three-phase.
\end{exercise}

\begin{exercise}
Compute the transfer function of an RC low-pass filter with $R =
1\,\si{\kilo\ohm}$, $C = 100\,\si{\nano\farad}$. Report the $-3\,\si{
\decibel}$ frequency and the magnitude in decibels at $f = 10\,\si{
\kilo\hertz}$ and $f = 100\,\si{\kilo\hertz}$.
\paragraph{Solution sketch.}
$H(\omega) = 1/(1+j\omega RC)$ with $RC = 10^{-4}\,\si{\second}$. The
cutoff is $f_{c} = 1/(2\pi RC) = 1.59\,\si{\kilo\hertz}$. At
$10\,\si{\kilo\hertz}$, $f/f_{c} = 6.28$, so $|H| = 1/\sqrt{1+6.28^{2}}
= 0.157$, that is $-16.1\,\si{\decibel}$. At $100\,\si{\kilo\hertz}$,
$f/f_{c} = 62.8$, $|H| = 0.0159$, that is $-36.0\,\si{\decibel}$, one
decade further down at the $-20\,\si{\decibel}$-per-decade slope.
\end{exercise}

\subsection*{Derivation}\pathtag{standard}

\begin{exercise}
Derive the impedance of an ideal capacitor in sinusoidal steady state
starting from $i(t) = C\,\dd v/\dd t$ and applying the phasor
representation $v(t) = \Re(\tilde V e^{j\omega t})$.
\paragraph{Solution sketch.}
Write $v(t) = \Re(\tilde V e^{j\omega t})$ and differentiate:
$\dd v/\dd t = \Re(j\omega\tilde V e^{j\omega t})$. Then $i(t) =
C\,\dd v/\dd t = \Re(j\omega C\tilde V e^{j\omega t})$, which in phasor
form reads $\tilde I = j\omega C\tilde V$. The impedance is the ratio
$Z_{C} = \tilde V/\tilde I = 1/(j\omega C) = -j/(\omega C)$. The
$-j$ shows the current leads the voltage by $90^\circ$, and the
magnitude $1/(\omega C)$ falls with frequency.
\end{exercise}

\begin{exercise}
Derive Th\'evenin's theorem for a linear two-terminal network
containing only independent sources and resistors. State precisely
where linearity is used.
\paragraph{Solution sketch.}
Apply a test current $I$ at the terminals and let $V$ be the resulting
terminal voltage. By superposition (which holds because the network is
linear), $V$ is the sum of the response to the internal sources with
$I = 0$, call it $V_{\text{oc}}$, plus the response to $I$ with the
internal sources deactivated, which is $-IR_{\text{eq}}$ where
$R_{\text{eq}}$ is the resistance seen from the terminals. Hence
$V = V_{\text{oc}} - IR_{\text{eq}}$, the terminal relation of a
voltage source $V_{\text{Th}} = V_{\text{oc}}$ in series with
$R_{\text{Th}} = R_{\text{eq}}$. Linearity enters at the superposition
step: without it the two contributions could not be added.
\end{exercise}

\begin{exercise}
Derive the unit-step response of a second-order series RLC circuit in
its three damping regimes by solving the characteristic equation
$s^{2}+2\zeta\omega_{n}s+\omega_{n}^{2}=0$ and applying initial
conditions $v_{C}(0) = 0$, $i_{L}(0) = 0$.
\paragraph{Solution sketch.}
The roots are $s = -\zeta\omega_{n}\pm\omega_{n}\sqrt{\zeta^{2}-1}$.
For $\zeta>1$ both roots are real and negative, giving $v_{C} = V_{s} +
A e^{s_{1}t} + B e^{s_{2}t}$; the two constants come from $v_{C}(0)=0$
and $\dd v_{C}/\dd t(0) = i_{L}(0)/C = 0$. For $\zeta=1$ the repeated
root $-\omega_{n}$ gives $v_{C} = V_{s} - V_{s}(1+\omega_{n}t)
e^{-\omega_{n}t}$. For $\zeta<1$ the complex roots
$-\zeta\omega_{n}\pm j\omega_{d}$ with $\omega_{d} =
\omega_{n}\sqrt{1-\zeta^{2}}$ give the ringing form
$v_{C} = V_{s}\!\left[1 - \frac{e^{-\zeta\omega_{n}t}}{\sqrt{1-\zeta^{2}}}
\sin(\omega_{d}t + \arccos\zeta)\right]$. Each form satisfies both
initial conditions; the script \texttt{code/rlc\_transient.py} confirms
all three numerically.
\end{exercise}

\begin{exercise}
Derive the maximum power transfer theorem for a Th\'evenin source
with open-circuit voltage $V_{\text{Th}}$ and source resistance
$R_{\text{Th}}$ driving a real load $R_{L}$. State the efficiency at
the matched condition and the conditions under which matched delivery
is appropriate.
\paragraph{Solution sketch.}
The load current is $I = V_{\text{Th}}/(R_{\text{Th}}+R_{L})$ and the
load power is $P_{L} = I^{2}R_{L} = V_{\text{Th}}^{2}R_{L}/(R_{\text{Th}}
+R_{L})^{2}$. Set $\dd P_{L}/\dd R_{L} = 0$: the numerator of the
derivative is $(R_{\text{Th}}+R_{L})^{2} - R_{L}\cdot 2(R_{\text{Th}}+
R_{L}) = (R_{\text{Th}}+R_{L})(R_{\text{Th}}-R_{L})$, which vanishes at
$R_{L} = R_{\text{Th}}$. The maximum power is $V_{\text{Th}}^{2}/
(4R_{\text{Th}})$. At the match the source resistor dissipates as much
as the load, so the efficiency is $50\%$, acceptable for signal
transmission but not for power delivery, where $R_{L}\gg R_{\text{Th}}$
keeps efficiency near $100\%$.
\end{exercise}

\begin{exercise}
Derive the resonance and bandwidth of a parallel RLC tuned circuit
driven by an ideal current source. Show that $Q_{\parallel} =
R/(\omega_{0}L)$ and that the bandwidth is $\omega_{0}/Q$.
\paragraph{Solution sketch.}
The parallel admittance is $Y = 1/R + j\omega C + 1/(j\omega L) = 1/R +
j(\omega C - 1/\omega L)$. The susceptance vanishes at $\omega_{0} =
1/\sqrt{LC}$, where $|Z|$ is maximum and equal to $R$; the voltage
response driven by a fixed current source peaks there. Writing the
admittance magnitude near resonance and finding the half-power points
where the susceptance equals the conductance $1/R$ gives the two
$-3\,\si{\decibel}$ frequencies separated by $\Delta\omega = 1/(RC)$.
Then $Q = \omega_{0}/\Delta\omega = \omega_{0}RC = R/(\omega_{0}L)$,
using $\omega_{0}^{2}LC = 1$, and the bandwidth is $\omega_{0}/Q$.
\end{exercise}

\begin{exercise}
Show that an ideal lossless transformer with turns ratio $n$ reflects
a secondary load impedance $Z_{L}$ to the primary as $Z_{\text{ref}}
= Z_{L}/n^{2}$.
\paragraph{Solution sketch.}
For an ideal transformer with $n = N_{2}/N_{1}$, the secondary voltage
is $V_{2} = nV_{1}$ and the currents satisfy $I_{1} = nI_{2}$ (power
conservation, $V_{1}I_{1} = V_{2}I_{2}$). The load fixes $V_{2} =
Z_{L}I_{2}$. The impedance seen at the primary is $Z_{\text{ref}} =
V_{1}/I_{1} = (V_{2}/n)/(nI_{2}) = (1/n^{2})(V_{2}/I_{2}) = Z_{L}/n^{2}$.
Impedance scales with the square of the turns ratio, which is why
transformers serve as impedance-matching networks.
\end{exercise}

\begin{exercise}
Derive the formula $P = \sqrt{3}V_{L}I_{L}\cos\phi$ for the total
real power delivered to a balanced three-phase wye-connected load,
starting from the single-phase formula and summing over the three
phases.
\paragraph{Solution sketch.}
Each phase of a balanced wye load carries the phase voltage
$V_{\phi} = V_{L}/\sqrt{3}$ and the line current $I_{L}$ at power
factor $\cos\phi$, so the per-phase real power is $V_{\phi}I_{L}\cos\phi$.
The total over three identical phases is $P = 3V_{\phi}I_{L}\cos\phi =
3(V_{L}/\sqrt{3})I_{L}\cos\phi = \sqrt{3}V_{L}I_{L}\cos\phi$. The same
result follows for a delta load by the dual substitution
$I_{\phi} = I_{L}/\sqrt{3}$, $V_{\phi} = V_{L}$, so the line-quantity
formula is connection-independent.
\end{exercise}

\subsection*{Estimation}\pathtag{core}

\begin{exercise}
Estimate the resistance of a typical household incandescent lamp from
its nameplate power and voltage ratings. State whether the resistance
is the cold resistance, the hot resistance, or an average; explain
the difference.
\paragraph{Reference estimate.}
A $60\,\si{\watt}$ lamp on a $230\,\si{\volt}$ supply draws $P/V =
0.26\,\si{\ampere}$, so $R = V^{2}/P = 230^{2}/60 \approx 880\,\si{
\ohm}$. This is the hot (operating) resistance, computed from the
rated steady-state dissipation. The cold filament resistance is roughly
ten to fifteen times smaller (tungsten's resistance rises with
temperature), which is why an incandescent lamp draws a large inrush
current at switch-on before the filament heats and the resistance
climbs to its operating value.
\end{exercise}

\begin{exercise}
Estimate the time constant of a smartphone's battery-charging circuit
during the constant-voltage phase of charging. State assumptions about
the battery's internal resistance and the supply voltage.
\paragraph{Reference estimate.}
Model the cell as a capacitance-like reservoir behind an internal
resistance. A $4000\,\si{\milli\ampere\hour}$ cell at $3.7\,\si{\volt}$
stores about $5\times 10^{4}\,\si{\joule}$; treating it as an
equivalent capacitor $C \sim Q/V = (4\,\text{Ah}\cdot 3600)/3.7
\approx 4000\,\si{\farad}$. With an internal resistance of order
$50\,\si{\milli\ohm}$, the electrical time constant $RC$ is of order
$200\,\si{\second}$, a few minutes. The constant-voltage phase trails
off over tens of minutes because the charge controller deliberately
tapers the current well beyond one time constant to protect the cell;
the $RC$ estimate sets the lower bound on how fast the terminal voltage
could relax, not the full charge time.
\end{exercise}

\begin{exercise}
Estimate the quality factor and bandwidth of a typical AM radio
receiver's tuned circuit, given that adjacent AM stations are spaced
$10\,\si{\kilo\hertz}$ in frequency and that strong adjacent-channel
rejection is required.
\paragraph{Reference estimate.}
A station near the middle of the AM band sits at about
$1\,\si{\mega\hertz}$. To pass one $10\,\si{\kilo\hertz}$ channel while
rejecting its neighbours, the tuned circuit needs a bandwidth of order
$10\,\si{\kilo\hertz}$, so $Q = f_{0}/\Delta f \approx 10^{6}/10^{4} =
100$. A single LC tank with $Q\approx 100$ is achievable but marginal
for sharp adjacent-channel rejection, which is why real receivers move
selectivity to a fixed intermediate-frequency stage with cascaded tuned
stages rather than relying on the antenna-stage $Q$ alone.
\end{exercise}

\begin{exercise}
Estimate the reactive power that a domestic refrigerator's motor
demands during start-up, given that the motor draws roughly $7$ times
its run current at start. State the duration of the start-up
transient and its impact on the household's instantaneous demand.
\paragraph{Reference estimate.}
A refrigerator compressor runs at perhaps $150\,\si{\watt}$, drawing
about $1\,\si{\ampere}$ at $230\,\si{\volt}$ with a run power factor near
$0.6$. At start the locked-rotor current is roughly $7\,\si{\ampere}$,
mostly reactive because the stalled motor looks like an inductor. The
start-up apparent power is then $230\cdot 7 \approx 1.6\,\si{\kilo\volt
\ampere}$, of which the reactive part is of order $1.3\,\text{kvar}$.
The transient lasts a few hundred milliseconds until the rotor reaches
speed. The household sees a brief sag and a current spike that the
circuit breaker is sized to tolerate but that can dim incandescent
lamps momentarily.
\end{exercise}

\begin{exercise}
Estimate the equivalent series resistance of a $1\,\si{\milli\farad}$
aluminium electrolytic capacitor at $120\,\si{\hertz}$ and at
$100\,\si{\kilo\hertz}$. State the dominant physical mechanism in
each frequency range.
\paragraph{Reference estimate.}
A $1\,\si{\milli\farad}$ aluminium electrolytic typically lists an ESR
of order $0.1$ to $0.3\,\si{\ohm}$ at $120\,\si{\hertz}$, dominated by
the electrolyte's ionic conductivity and the dielectric loss of the
oxide film. At $100\,\si{\kilo\hertz}$ the ESR falls to roughly
$0.02$ to $0.05\,\si{\ohm}$ as the dielectric loss term shrinks, then
the impedance begins to rise again because the equivalent series
inductance of the wound foil takes over above the self-resonant
frequency. The data file \texttt{data/capacitor\_dielectrics.csv}
records the qualitative ESR-versus-frequency trend for several
dielectric classes.
\end{exercise}

\begin{exercise}
Estimate the total energy a $230\,\si{\volt}$, $50\,\si{\hertz}$
household consumes during the first cycle after a power outage's
restoration, given that the dominant inrush is the cold filaments of
incandescent lamps and the capacitor charging in switched-mode power
supplies. State the assumed installed wattage of each load type.
\paragraph{Reference estimate.}
Assume $500\,\si{\watt}$ of incandescent lamps and $300\,\si{\watt}$ of
switched-mode supplies installed. The lamps' cold filaments draw about
ten times their steady current, so the first half-cycle
($10\,\si{\milli\second}$) passes a peak of order
$5\,\si{\kilo\watt}$, delivering roughly $500\,\text{W}\cdot 0.01\,
\text{s}\cdot 10 \approx 50\,\si{\joule}$ in that first cycle before the
filaments heat. The switched-mode supplies charge their input bulk
capacitors through a near-short, drawing a narrow current spike that
deposits a comparable few tens of joules. The first-cycle energy is
therefore of order $100\,\si{\joule}$, dominated by inrush rather than
steady draw; this is the load the inrush-limiting resistors and
breakers must survive.
\end{exercise}

\subsection*{Simulation}\pathtag{mastery}

\begin{exercise}
Implement a SPICE-style nodal-analysis solver in a language of choice.
The input is a netlist (resistors, voltage sources, current sources);
the output is the steady-state node voltages. Test on a network of
$10$ nodes with $20$ resistors and verify against hand calculation on
a $3$-node sub-circuit.
\paragraph{Model-answer sketch.}
Stamp each element into a conductance matrix $G$ and a source vector
$z$: a resistor between nodes $a$ and $b$ adds $1/R$ to $G_{aa}$ and
$G_{bb}$ and subtracts it from $G_{ab}$ and $G_{ba}$; a current source
adds to $z$; a voltage source needs modified nodal analysis with one
extra unknown current per source. Solve $Gv = z$. The reference
implementation \texttt{code/nodal\_solver.py} does exactly this and
reproduces the $v_{1} = 4\,\si{\volt}$ divider result on the
three-node test circuit. The graders should check that the
$10$-node solution satisfies KCL at every node to within the solver's
numerical tolerance.
\end{exercise}

\begin{exercise}
Implement a transient simulator for first- and second-order linear
circuits using trapezoidal-rule integration with step size $\Delta
t$. Simulate the unit-step response of an RLC circuit at $\zeta = 0.1,
0.5, 1, 2$ and compare to the analytical solution. Identify the step
size at which the trapezoidal-rule truncation error becomes visible.
\paragraph{Model-answer sketch.}
Use the state vector $(v_{C}, i_{L})$ with $\dd v_{C}/\dd t = i_{L}/C$
and $\dd i_{L}/\dd t = (V_{s}-Ri_{L}-v_{C})/L$, and advance with the
trapezoidal rule (solve the implicit $2\times 2$ system each step). The
reference \texttt{code/rlc\_transient.py} matches the closed-form step
response to under a millivolt at $\omega_{n}\Delta t = 0.01$. Truncation
error grows as $(\omega_{n}\Delta t)^{2}$; it becomes visible to the eye
once $\omega_{n}\Delta t$ approaches $0.2$ to $0.5$, where the ringing
period is sampled fewer than about a dozen times. Below that the
A-stable trapezoidal scheme tracks all four damping ratios without
artificial growth.
\end{exercise}

\begin{exercise}
Implement the AC sweep of a passive bandpass filter and plot its Bode
diagram. Compare to the analytical $H(\omega)$ and identify the
component tolerances that would shift the centre frequency by
$\pm 5\,\%$ and the $Q$ by $\pm 10\,\%$.
\paragraph{Model-answer sketch.}
Compute $H(\omega) = R/(R + j\omega L + 1/(j\omega C))$ over a
logarithmic frequency grid and plot $20\log_{10}|H|$ and $\angle H$, as
\texttt{code/ac\_sweep.py} does. Since $f_{0} = 1/(2\pi\sqrt{LC})$, a
$\pm 5\%$ shift in $f_{0}$ follows from about $\pm 10\%$ combined
tolerance in $LC$ (roughly $\pm 10\%$ on either $L$ or $C$ alone). Since
$Q = \omega_{0}L/R$, a $\pm 10\%$ swing in $Q$ comes from about
$\pm 10\%$ in $R$ or, after the $f_{0}$ shift, the residual sensitivity
to $L$. The synthetic noisy sweep in
\texttt{data/bandpass\_sweep.csv} lets the student practise extracting
$f_{0}$ and $Q$ from imperfect data.
\end{exercise}

\subsection*{Design}\pathtag{standard}

\begin{exercise}
Design a first-order low-pass filter to attenuate the $50\,\si{\hertz}$
mains-frequency component of a sensor signal by $40\,\si{\decibel}$
while passing DC-to-$1\,\si{\hertz}$ within $1\,\si{\decibel}$.
Specify $R$, $C$, the cutoff frequency, and the resulting in-band
ripple. Identify the trade-off between attenuation at $50\,\si{\hertz}$
and the settling time of the filter to a step input.
\paragraph{Solution sketch.}
A first-order filter rolls off at $20\,\si{\decibel}$ per decade, so
$40\,\si{\decibel}$ of attenuation at $50\,\si{\hertz}$ needs the cutoff
two decades below, $f_{c} \approx 0.5\,\si{\hertz}$. That places the
$1\,\si{\hertz}$ pass-band edge above $f_{c}$, costing more than
$1\,\si{\decibel}$; a single pole cannot meet both edges, so either
relax the pass-band spec or move to a higher-order filter. Taking
$f_{c} = 0.5\,\si{\hertz}$ with $C = 1\,\si{\micro\farad}$ gives $R =
1/(2\pi f_{c}C) = 318\,\si{\kilo\ohm}$. The trade-off is direct: a
lower $f_{c}$ buys more $50\,\si{\hertz}$ rejection but lengthens the
step settling time $\tau = 1/(2\pi f_{c})$, here about $0.3\,\si{
\second}$, so the sensor responds sluggishly.
\end{exercise}

\begin{exercise}
Design an inverting op-amp amplifier with voltage gain $-10$, input
impedance $\ge 10\,\si{\kilo\ohm}$, and bandwidth $\ge 100\,\si{\kilo
\hertz}$. Specify the feedback and input resistors, identify the
required op-amp gain-bandwidth product, and select a candidate op-amp
from a current commercial catalogue (as of 2025).
\paragraph{Solution sketch.}
For an inverting amplifier the gain is $-R_{f}/R_{\text{in}}$ and the
input impedance is $R_{\text{in}}$. Set $R_{\text{in}} = 10\,\si{\kilo
\ohm}$ to meet the input-impedance floor, then $R_{f} = 100\,\si{\kilo
\ohm}$ for gain $-10$. The closed-loop bandwidth is the gain-bandwidth
product divided by the noise gain $(1+R_{f}/R_{\text{in}}) = 11$, so a
$\ge 100\,\si{\kilo\hertz}$ bandwidth needs $\text{GBW}\ge 1.1\,\si{
\mega\hertz}$. A common general-purpose part with a few-megahertz GBW
and rail-to-rail output (as of 2025, a part in the class of the
MCP6002 or TLV9001 family) clears the requirement with margin.
\end{exercise}

\begin{exercise}
Design a power-factor-correction capacitor bank for an industrial
load consuming $200\,\si{\kilo\watt}$ at $400\,\si{\volt}$ line-to-
line, three-phase, $50\,\si{\hertz}$, with a measured power factor
of $0.7$ lagging. Specify the capacitance per phase needed to bring
the power factor to $0.95$ lagging, the kvar rating, and the
operational considerations (over-voltage on light load, switching
transients).
\paragraph{Solution sketch.}
At $\cos\phi_{1} = 0.7$ ($\phi_{1} = 45.6^\circ$) the reactive power is
$Q_{1} = P\tan\phi_{1} = 200\cdot 1.02 = 204\,\text{kvar}$; at
$\cos\phi_{2} = 0.95$ ($\phi_{2} = 18.2^\circ$), $Q_{2} = 200\cdot
0.329 = 65.7\,\text{kvar}$. The bank supplies $Q_{C} = 138\,\text{kvar}$,
or $46\,\text{kvar}$ per phase. For a delta-connected bank each
capacitor sees the line voltage $400\,\si{\volt}$: $C = Q_{C,\text{ph}}
/(\omega V_{L}^{2}) = 46000/(2\pi\cdot 50\cdot 400^{2}) = 915\,\si{
\micro\farad}$ per phase. Operationally, a fixed bank over-corrects on
light load (raising voltage), so step-switched banks or detuned
reactors are used; inrush on switching is limited by series reactors or
pre-charge.
\end{exercise}

\begin{exercise}
Design the decoupling-capacitor network for a digital circuit board
carrying eight $3.3\,\si{\volt}$ logic chips, each drawing peak pulse
currents of $200\,\si{\milli\ampere}$ with $1\,\si{\nano\second}$
edges. Specify the bulk and bypass capacitor values, placement, and
the target supply-rail impedance from DC to $100\,\si{\mega\hertz}$.
\paragraph{Solution sketch.}
Allow a rail droop of, say, $50\,\si{\milli\volt}$ for a
$200\,\si{\milli\ampere}$ pulse with a $1\,\si{\nano\second}$ edge. The
target impedance is $Z_{\text{target}} = \Delta V/\Delta I =
0.05/0.2 = 0.25\,\si{\ohm}$ across the band. Place one
$100\,\si{\nano\farad}$ ceramic bypass capacitor at each chip's supply
pin (eight total) to source the high-frequency edge current, plus a
shared bulk reservoir of $10$ to $47\,\si{\micro\farad}$ for the
low-frequency bulk demand. The $1\,\si{\nano\second}$ edge has spectral
content to a few hundred megahertz, so keep the bypass-capacitor loop
inductance minimal (vias straight to a ground plane); above
$100\,\si{\mega\hertz}$ the ESL of the ceramics dominates and the plane
capacitance carries the rail.
\end{exercise}

\subsection*{Diagnosis and reverse engineering}\pathtag{standard}

\begin{exercise}
A bench LCR meter measures an unmarked component as $33\,\si{\nano
\farad}$ at $1\,\si{\kilo\hertz}$, $32\,\si{\nano\farad}$ at $10\,\si{
\kilo\hertz}$, and $25\,\si{\nano\farad}$ at $1\,\si{\mega\hertz}$.
Identify the dielectric class (X7R, NPO, electrolytic) most consistent
with this signature, and the secondary measurement that would
distinguish among the candidates.
\paragraph{Solution sketch.}
The measured capacitance falls by about a quarter from
$1\,\si{\kilo\hertz}$ to $1\,\si{\mega\hertz}$, a frequency dependence
too strong for NP0/C0G (which is flat to fractions of a percent) and
too mild for an aluminium electrolytic (which loses far more and is
polarised besides). The signature fits a class II ceramic such as X7R,
whose permittivity drops with frequency and bias. The distinguishing
secondary measurement is the DC-bias dependence: apply a few volts of
DC bias and re-measure; X7R loses substantial capacitance under bias,
NP0 does not, and an electrolytic would show polarity sensitivity. The
data file \texttt{data/capacitor\_dielectrics.csv} lists the
contrasting stabilities.
\end{exercise}

\begin{exercise}
A transient on an oscilloscope shows a damped sinusoid following a
$5\,\si{\volt}$ step input to a digital circuit's supply rail, with
peak-to-peak amplitude $300\,\si{\milli\volt}$, period $25\,\si{
\nano\second}$, and approximately three visible cycles before decay.
Estimate the $L$ and $C$ of the supply-rail loop, identify the source
of $L$ (PCB-trace inductance, bond-wire inductance, capacitor ESL),
and propose a measurement that distinguishes among them.
\paragraph{Solution sketch.}
The ringing period $T = 25\,\si{\nano\second}$ gives an oscillation
frequency $f = 40\,\si{\mega\hertz}$, so $\omega = 2.5\times 10^{8}\,\si{
\per\second}$ and $LC = 1/\omega^{2} = 1.6\times 10^{-17}\,\si{
\second\squared}$. Three visible cycles means $Q \approx \pi\cdot 3
\approx 10$. If the decoupling capacitor is $100\,\si{\nano\farad}$,
then $L = 1/(\omega^{2}C) = 1.6\times 10^{-17}/10^{-7} = 0.16\,\si{
\nano\henry}$, far too small; the resonance must instead involve a
smaller effective capacitance (the capacitor's ESL series-resonating
with a few nanohenries of loop inductance). Estimating $L \sim 5\,\si{
\nano\henry}$ implies $C \sim 3\,\si{\nano\farad}$, consistent with
parasitic capacitance plus a small bypass. To distinguish the
inductance source, add a known series inductance and watch the period
shift: PCB-trace inductance changes with layout, bond-wire inductance
does not, and ESL changes when the capacitor is swapped.
\end{exercise}

\begin{exercise}
A two-port network is characterised by measuring its input-impedance
magnitude with the output open ($100\,\si{\ohm}$) and shorted ($25
\,\si{\ohm}$). Identify whether the network is more consistent with
a T-network of resistors, an L-network, or a transformer-coupled
network, and propose the supplementary measurement that would close
the identification.
\paragraph{Solution sketch.}
The open- and short-circuit input impedances $Z_{\text{oc}} =
100\,\si{\ohm}$ and $Z_{\text{sc}} = 25\,\si{\ohm}$ give a geometric
mean $\sqrt{Z_{\text{oc}}Z_{\text{sc}}} = 50\,\si{\ohm}$, the image
impedance, and a ratio $Z_{\text{oc}}/Z_{\text{sc}} = 4$. All three
candidate networks can reproduce two real measurements, so two numbers
do not close the identification. Both resistive networks are
frequency-independent, while a transformer-coupled network shows
frequency dependence and a low-frequency roll-off. The supplementary
measurement is a frequency sweep of the open-circuit impedance:
flat means a resistive T or L network (distinguished by measuring the
reverse open- and short-circuit impedances for symmetry), while a
rising or falling magnitude with frequency marks the transformer.
\end{exercise}

\subsection*{Failure analysis}\pathtag{core}

\begin{exercise}
A circuit board with separate analog and digital sections shares a
common ground at the power-supply input. The analog signal contains a
$200\,\si{\milli\volt}$ peak-to-peak hash that correlates with
digital-clock activity. Identify the failure mode and the layout
change that would eliminate it.
\paragraph{Solution sketch.}
The failure mode is a shared-ground-impedance coupling (a ground loop):
digital switching currents return through the common ground trace,
whose finite inductance and resistance drop a voltage that the analog
section reads as common-mode hash, synchronised to the clock. The cure
is to give analog and digital sections their own ground returns joined
at a single star point near the data converter, so the digital return
current never flows through the analog reference path. A continuous
ground plane with a deliberate partition under the converter
accomplishes the same separation while keeping return paths short.
\end{exercise}

\begin{exercise}
An electrolytic capacitor in an industrial drive bulges and vents
after $18$ months of service. The drive's nameplate rates the
capacitor for $10\,000$ hours at $105\,\si{\celsius}$. Identify the
failure mechanism, the likely thermal stress, and the design change
that would extend the life.
\paragraph{Solution sketch.}
The mechanism is electrolyte dry-out: ripple current heats the
capacitor through its ESR, the sealed electrolyte slowly escapes, the
ESR rises, and the heating accelerates until the vent opens. Aluminium
electrolytic life roughly doubles for every $10\,\si{\celsius}$ of
reduced core temperature (an Arrhenius rule). Eighteen months is about
$13\,000$ hours, so if the part lasted only that long against a
$10\,000$-hour-at-$105\,\si{\celsius}$ rating, the core ran far above
ambient, pointing to excessive ripple heating or poor airflow. The fix
is to lower the operating temperature: a higher-ripple-rated or
lower-ESR capacitor, paralleled units to share ripple, better cooling,
or a longer-life ($105\,\si{\celsius}$, $5000$-to-$10\,000$-hour-plus)
series part.
\end{exercise}

\begin{exercise}
A relay's contacts pit and fail after $50\,000$ operations against a
nameplate rating of $10^{6}$ operations. The load is a small
inductive solenoid. Identify the failure mode, the missing protective
element, and the SPICE-level simulation that would have predicted
it.
\paragraph{Solution sketch.}
When the contacts open, the solenoid's inductance forces the current to
continue, driving the contact voltage up until an arc strikes across
the gap; the arc erodes and pits the contacts, shortening life. The
missing protective element is a freewheeling path across the inductor: a
flyback diode for a DC solenoid, or an RC snubber for AC. A transient
simulation that opens the contact (a switch with finite off
resistance) while the inductor carries current would show the voltage
spike $v = L\,\dd i/\dd t$ rising into the hundreds of volts, predicting
the arc; adding the snubber in the model clamps it.
\end{exercise}

\begin{exercise}
An audio amplifier reproduces a clean sinusoid at low amplitudes but
clips its output at peak amplitudes corresponding to less than its
rated output power. Identify three candidate causes (supply-rail
sag, output-stage thermal limiting, feedback-loop slew limiting) and
the oscilloscope measurements that would distinguish them.
\paragraph{Solution sketch.}
Supply-rail sag flattens the output at a ceiling that tracks the supply
voltage; probe the rail during the clip and watch it droop with the
signal envelope. Output-stage thermal limiting appears only after the
stage has heated, so the clip onset drifts with time and temperature;
monitor the heat-sink temperature and the clip threshold over minutes.
Slew limiting is frequency-dependent: it produces a triangular distortion
of fast edges while slow large-signal sinusoids pass cleanly, so sweep
the input frequency at fixed amplitude and find the frequency where the
peaks first distort. The three signatures (supply-tracking ceiling,
time-drifting onset, frequency-dependent edge distortion) separate the
causes.
\end{exercise}

\subsection*{Judgment}\pathtag{standard}

\begin{exercise}
A consulting client must choose between a $1\,\si{\milli\farad}$
aluminium electrolytic capacitor and a $1\,\si{\milli\farad}$
multi-layer-ceramic stack for a DC-bus decoupling application at
$48\,\si{\volt}$. State the engineering trade-off in terms of ESR,
ESL, life, cost, and failure mode. Identify the working application
in which each is the preferred choice.
\paragraph{Model-answer sketch.}
The electrolytic offers large capacitance per unit cost but high ESR,
appreciable ESL, a finite wear-out life that shortens with temperature,
and a vent-and-dry-out failure mode. The ceramic stack offers very low
ESR and ESL and long life, but costs more per microfarad, loses
capacitance under DC bias (class II dielectrics), and can fail short
through mechanical cracking. For bulk energy storage and low-frequency
ripple on the $48\,\si{\volt}$ bus, the electrolytic is the economical
choice; for high-frequency decoupling and where lifetime and ESR
dominate, the ceramic stack wins. Many designs use both: electrolytic
for bulk, ceramics for the high-frequency edge. The data file
\texttt{data/capacitor\_dielectrics.csv} summarises the contrasting
properties.
\end{exercise}

\begin{exercise}
A regulator is considering raising the power-factor minimum for
industrial consumers from $0.85$ to $0.95$, with a tariff penalty on
consumers below the new minimum. Identify who benefits and who pays
under the change, and the technical argument for and against the new
minimum.
\paragraph{Model-answer sketch.}
The utility and the wider rate base benefit: a higher power-factor floor
reduces the reactive current the distribution network carries, cutting
$I^{2}R$ losses and deferring capacity upgrades. Consumers below the new
floor pay, either through the tariff penalty or through the capital cost
of correction equipment. The technical argument for the change is that
reactive current loads the network without delivering real energy, so
charging for it aligns incentives. The argument against is that the last
increment from $0.85$ to $0.95$ costs each consumer disproportionately
(diminishing returns on correction capacitors), and over-correction
risks light-load over-voltage and resonance with the supply
inductance. A reasonable judgment weighs the network savings against the
distributed correction cost and may phase the requirement in.
\end{exercise}

\begin{exercise}
A team must choose between using passive RC filtering and an active
op-amp-based filter for a $4$-Hz-cutoff anti-aliasing filter ahead of
an analog-to-digital converter. Identify the design trade-off in
terms of component count, noise, drift, power consumption, and
output-impedance behaviour.
\paragraph{Model-answer sketch.}
A passive RC filter is cheap, needs no power, adds no active-device
noise, and cannot become unstable, but a $4\,\si{\hertz}$ cutoff with
realistic capacitor values needs large resistors that raise the output
impedance and load the following stage, and a single pole gives only a
gentle roll-off. An active op-amp filter buffers the output to a low
impedance, allows a sharp higher-order response in one stage, and
isolates the source, but it adds component count, op-amp noise and
offset drift, and power consumption, and demands attention to
stability. For an anti-aliasing filter ahead of a converter, the
active filter's low output impedance and steep roll-off usually justify
its cost; the passive filter suits low-power or low-budget designs that
can tolerate a soft cutoff.
\end{exercise}

\begin{exercise}
A design review questions whether a $230\,\si{\volt}$ AC line input
to a small piece of laboratory equipment really needs a double-pole
switch (interrupting both line and neutral) or whether a single-pole
switch (interrupting only line) is sufficient. State the safety
argument on each side and the regional regulation (IEC 60364) that
governs the choice.
\paragraph{Model-answer sketch.}
The single-pole switch interrupts only the line conductor; it is simpler
and cheaper and is sufficient where the supply polarity is reliably
known and the protective earthing system guarantees which conductor is
live. The double-pole switch interrupts both line and neutral, so the
equipment is fully isolated even if the plug is reversed or the
installation's line and neutral are swapped, which matters for portable
laboratory gear used across sockets of uncertain wiring. The safety
argument for double-pole is fail-safe isolation under wiring faults; the
argument for single-pole is reduced cost and complexity in a fixed,
known installation. IEC 60364 (low-voltage installations) governs the
earthing system and disconnection requirements that decide which is
acceptable in a given jurisdiction and application.
\end{exercise}
